\chapter{文本分析}
\label{chap:nlp-text-analysis}

读到“小王阅读通知”，读者会自然地把它切成三个词，知道“小王”是读的人、“通知”是被读的东西，也知道这句话陈述了一件事而非提出请求。要让程序给出同样的结果，就必须把这些理解写成明确的结构：词的边界与类别、词与词之间的组合或依存、谓词与论元的语义角色、可求值的逻辑式或查询、跨句的指称与篇章关系，以及话语在交流中的作用。本章回答怎样从文本中恢复这些结构。第\ref{chap:nlp-introduction}章已经区分了任务输出与判断依据，第1部分的语言模型提供了上下文表示与生成能力；这里要做的，是把它们变成可以恢复到原文位置、可以检查合法性、可以交给下游使用的分析结果。

这些任务共同面对五项挑战。第一，分析单位本身不确定：中文书写不标词界，语言学意义上的词与模型使用的子词词元也不一致。第二，输出结构多样且受约束：BIO标签不能以I开头，依存树必须单根无环，语义图允许多入边，SQL必须能在给定数据库上执行；输出空间随句长指数增长，又必须保证合法。第三，歧义往往要靠远距离上下文消解：“新学生代表”的分组、代词的先行项、“他将来北京”中“将来”的切分，都不能只看局部。第四，人工标注昂贵，各语料的规范又不一致，同一句话在不同树库中的“正确”分析并不相同。第五，篇章和语用分析依赖文本之外的背景，例如说话人的意图和交流双方的共同知识。

方法的演进可以沿一条主线来看。生成式统计模型（HMM、PCFG）先给出可以从计数估计的概率分解，但观测独立假设限制了可用的证据；判别式全局归一化模型（CRF、对数线性解析器）去掉这一假设，允许使用任意重叠的特征；神经编码器配合结构解码，用上下文表示代替人工特征，同时保留合法结构的推断；预训练编码器进一步减少了对任务标注和特征工程的依赖；线性化生成与大语言模型提供了覆盖各种结构的统一接口，代价是需要约束解码与结果核验来保证输出合法、忠实于原文。

第\ref{sec:nlp-analysis-paradigms}节先把这些方法归入五种结构预测范式：序列标注、跨度、转移系统、图与区间动态规划、约束生成，一次性给出各范式的输出表示、得分分解、推断算法、训练目标与复杂度。此后的词法、句法、语义、篇章和语用各节按“任务挑战—方法谱系—结构与目标—对比与代价”展开，注明每种方法所属的范式，通用机制只回指第\ref{sec:nlp-analysis-paradigms}节。五个分析层次是观察角度，不是必须依次执行的流水线：上文出现的实体可能改变当前词义，交流背景也可能排除一种句法解释。第\ref{chap:nlp-information-extraction}章的信息抽取与第\ref{chap:nlp-classification-matching-organization}章的分类、匹配与组织沿用同一套范式归类。

本章统一使用从零开始、左闭右开的原文字符跨度$[s,e)$，字符按Unicode计数而非UTF-8字节；数学推导中的位置可以从1开始，但每个结果都要恢复到原文偏移后再评价。语言词、模型词元与内部向量分别命名。未注明来源的短句、概率、得分和动作轨迹均为教学构造，不是模型实测。

\section{结构预测的共同范式}
\label{sec:nlp-analysis-paradigms}

本章的任务看起来各不相同：分词输出词界，句法分析输出树，语义解析输出逻辑式或SQL，共指输出提及的分组。它们却可以写成同一个问题。给定输入$x$，记它所有合法输出组成的集合为$\mathcal{Y}(x)$，模型学习一个得分函数$s_{\bm\theta}(x,y)$，预测时求
\begin{equation}
  \hat y=\argmax_{y\in\mathcal{Y}(x)}s_{\bm\theta}(x,y).
  \label{eq:nlp-analysis-structured-argmax}
\end{equation}
这类输出是带内部约束的组合对象、而不是单个类别的问题，称为\term{结构化预测}（structured prediction）。式~\eqref{eq:nlp-analysis-structured-argmax}立即带来三项困难。其一，$\mathcal{Y}(x)$的规模随输入长度指数增长，不能逐个打分；其二，合法性约束（BIO标签不能以I开头、依存树单根无环、SQL必须可解析）必须由某个机制保证，不能指望模型自动学会；其三，训练只看到参考结构$y^\star$，学习目标怎样与$\argmax$推断配合，决定了训练与预测是否一致。

不同方法对这三项困难给出了不同回答。把结构拆成逐位置标签的是\term{序列标注}（sequence labeling），对候选区间评分的是跨度范式，把构造过程写成一串动作的是转移系统，把得分分解到弧或区间上再做精确动态规划的是图方法，把结构写成词元序列再受约束地生成的是生成范式。图~\ref{fig:nlp-analysis-paradigms}用同一个短句对照这五种输出表示。本节依次写出每种范式的输出表示、得分分解、推断算法、训练目标和复杂度，第\ref{sec:nlp-analysis-paradigm-comparison}小节用一张表收拢。此后各节只写任务特有的部分，范式的通用机制回指本节；第\ref{chap:nlp-information-extraction}章与第\ref{chap:nlp-classification-matching-organization}章也按这套归类组织。

\begin{figure*}[htbp]
  \centering
  % 第九卷 9.1 节：同一输入“小王／阅读／通知”在五种结构预测范式下的输出表示（教学示意）。
% 行依次为序列标注（BMES）、跨度（嵌套区间）、转移系统（中间状态）、图（依存弧）、约束生成（括号序列前缀与允许集合）。
\begin{tikzpicture}[
  x=1mm,y=1mm,figure picture,
  rowlabel/.style={font=\figurefont\footnotesize,text=FigureInk,anchor=east,align=right},
  infer/.style={font=\figurefont\scriptsize,text=FigureMuted,anchor=west,align=left},
  ch/.style={rectangle,rounded corners=.6mm,draw=FigureInput,fill=FigureInputFill,
    line width=.6pt,minimum width=8mm,minimum height=6mm,inner sep=0pt,
    font=\figurefont\footnotesize},
  wd/.style={rectangle,rounded corners=.8mm,draw=FigureInput,fill=FigureInputFill,
    line width=.6pt,minimum width=16mm,minimum height=6mm,inner sep=0pt,
    font=\figurefont\footnotesize},
  tag/.style={font=\figurefont\footnotesize,text=BookTeal},
  span/.style={draw=FigureOutput,line width=.9pt},
  spanlab/.style={font=\figurefont\scriptsize,text=FigureInk,fill=white,inner sep=1pt},
  cfg/.style={rectangle,rounded corners=.8mm,draw=FigureModel,fill=white,
    line width=.6pt,minimum height=6.5mm,inner xsep=2pt,font=\figurefont\scriptsize},
  arc/.style={draw=BookTeal,line width=.9pt,-{Stealth[length=1.6mm,width=1.1mm]}},
  cand/.style={draw=FigureMuted,line width=.5pt,densely dashed,-{Stealth[length=1.4mm,width=1mm]}},
  sep/.style={draw=FigureGrid,line width=.4pt}
]
  % ---------- 行1：序列标注 ----------
  \node[rowlabel] at (20,0) {序列标注};
  \foreach \x/\c/\t in {30/小/B,39/王/E,55/阅/B,64/读/E,80/通/B,89/知/E}{
    \node[ch] at (\x,0) {\c};
    \node[tag] at (\x,-6.5) {\t};
  }
  \node[infer] at (122,-2) {Viterbi＋合法转移掩码\\$O(nL^2)$};
  \draw[sep] (0,-11) -- (158,-11);

  % ---------- 行2：跨度 ----------
  \node[rowlabel] at (20,-19) {跨度};
  \node[wd] at (34.5,-17) {小王};
  \node[wd] at (59.5,-17) {阅读};
  \node[wd] at (84.5,-17) {通知};
  \draw[span] (26.5,-22.5) -- ++(0,-1.2) -- (42.5,-23.7) -- ++(0,1.2);
  \node[spanlab,anchor=west] at (43.5,-23.4) {NP};
  \draw[span] (76.5,-22.5) -- ++(0,-1.2) -- (92.5,-23.7) -- ++(0,1.2);
  \node[spanlab,anchor=west] at (93.5,-23.4) {NP};
  \draw[span] (51.5,-25.5) -- ++(0,-1.4) -- (92.5,-26.9) -- ++(0,1.4);
  \node[spanlab,anchor=west] at (93.5,-26.7) {VP};
  \draw[span] (26.5,-28.6) -- ++(0,-1.6) -- (92.5,-30.2) -- ++(0,1.6);
  \node[spanlab,anchor=west] at (93.5,-30.0) {S};
  \node[infer] at (122,-21) {枚举候选$O(nW)$并剪枝\\不重叠选择或CKY};
  \draw[sep] (0,-34) -- (158,-34);

  % ---------- 行3：转移系统 ----------
  \node[rowlabel] at (20,-42) {转移系统};
  \node[cfg] (stk) at (37,-42) {栈：ROOT\enspace 阅读};
  \node[cfg] (buf) at (64,-42) {缓冲区：通知};
  \node[cfg] (arcs) at (91,-42) {已建：阅读$\to$小王};
  \node[font=\figurefont\scriptsize,text=BookTeal] at (64,-48.5) {下一动作：SHIFT};
  \node[infer] at (122,-42) {贪心或束搜索\\$2n$个动作};
  \draw[sep] (0,-53) -- (158,-53);

  % ---------- 行4：图与区间动态规划 ----------
  \node[rowlabel] at (20,-69) {图与区间DP};
  \node[wd,minimum width=12mm,draw=FigureMuted,fill=white] (r) at (30,-71) {ROOT};
  \node[wd,minimum width=13mm] (w1) at (50,-71) {小王};
  \node[wd,minimum width=13mm] (w2) at (71,-71) {阅读};
  \node[wd,minimum width=13mm] (w3) at (92,-71) {通知};
  \draw[arc] ([xshift=2mm]r.north) to[out=55,in=125] ([xshift=-1mm]w2.north);
  \draw[arc] ([xshift=-3mm]w2.north) to[out=120,in=60] ([xshift=2mm]w1.north);
  \draw[arc] ([xshift=3mm]w2.north) to[out=60,in=120] ([xshift=-2mm]w3.north);
  \draw[cand] (w1.south) to[out=-35,in=-145] (w3.south);
  \node[infer] at (122,-69) {Eisner $O(n^3)$\\Chu–Liu/Edmonds $O(n^2)$};
  \draw[sep] (0,-79) -- (158,-79);

  % ---------- 行5：约束生成 ----------
  \node[rowlabel] at (20,-87) {约束生成};
  \node[cfg,anchor=west,draw=FigureInput] (pre) at (24,-87)
    {已生成：(S (NP 小王) (VP 阅读 (NP 通知};
  \node[cfg,anchor=west,draw=BookTeal] (next) at ([xshift=4mm]pre.east)
    {下一词元只允许：)};
  \draw[figure flow] (pre.east) -- (next.west);
  \node[infer] at (122,-87) {束搜索＋文法掩码\\输出解析回结构};
\end{tikzpicture}

  \caption{同一输入“小王／阅读／通知”在五种范式下的输出表示。从上到下：序列标注为每个字符给一个BMES标签；跨度范式给出带类别的区间（NP、VP、S），区间可以嵌套；转移系统展示一个中间状态（栈、缓冲区、已建的弧）及下一动作；图方法在“根＋词”的有向图上选出一棵依存树，虚线为未入选的候选弧；约束生成把成分树线性化为括号序列，图中为生成到一半时的前缀及下一词元的允许集合。右侧注出各范式的推断方式。树与弧均为教学设定。}
  \label{fig:nlp-analysis-paradigms}
\end{figure*}

\subsection{输出结构与分解方式}
\label{sec:nlp-analysis-paradigm-structures}

先看两个最小的歧义。“研究生命”可以切成“研究／生命”，也可以切成“研究生／命”；第二个字“究”是否是词尾，取决于第三个字“生”归到哪个词。“新学生代表”可以是“[新学生]代表”，也可以是“新[学生代表]”；“学生”先和谁组合，同时决定了“新”修饰谁。两例中各位置的决定互相牵连：若为每个位置独立选择最高分的标签或父节点，拼起来的结果可能根本不是一个合法切分或一棵树。因此分析输出是一个整体，评分与推断都要在整体层面进行。

为在各范式间沿用同一组记号，设输入已经过编码器（第\ref{sec:nlp-llm-computation}节的Transformer层或第\ref{sec:nlp-classical-bert}小节的BERT类编码器），得到逐位置表示$\bm H\in\mathbb R^{n\times d}$，第$i$行$\bm h_i$是第$i$个单位的上下文向量；$n$是分析单位数，$d$是隐藏维度。若编码器在子词上运行而分析单位是字或词，需先按第\ref{sec:nlp-analysis-paradigm-tagging}小节的对齐规则取出相应位置。结构$y$被看成若干\term{部件}（part）的集合，整体得分是部件得分之和：
\begin{equation}
  s_{\bm\theta}(x,y)=\sum_{p\in y}\phi_{\bm\theta}(p;\bm H).
  \label{eq:nlp-analysis-part-score}
\end{equation}
部件随范式而变：序列标注中是“位置$i$取标签$a$”和“相邻标签$(a,b)$”；跨度范式中是“区间$[i,j]$取类别$c$”；转移系统中是“在状态$c_t$执行动作$a_t$”；图方法中是“词$h$支配词$m$的弧”或“区间$[i,j]$构成一个成分”；生成范式中是“第$t$个输出词元”。部件越局部，$\argmax$越容易用动态规划精确求解，能表达的约束和交互也越少，这是贯穿全节的权衡。

训练目标分三类，区别在于把参考结构$y^\star$与哪些竞争者比较：
\begin{equation}
  \begin{aligned}
    \mathcal L_{\mathrm{local}}&=-\sum_{k}\log p_{\bm\theta}\bigl(d_k^\star\mid x,\text{第}k\text{个决策的上下文}\bigr),\\
    \mathcal L_{\mathrm{global}}&=-s_{\bm\theta}(x,y^\star)+\log\sum_{y\in\mathcal Y(x)}\exp s_{\bm\theta}(x,y),\\
    \mathcal L_{\mathrm{margin}}&=\max_{y\in\mathcal Y(x)}\bigl[s_{\bm\theta}(x,y)+\Delta(y,y^\star)\bigr]-s_{\bm\theta}(x,y^\star).
  \end{aligned}
  \label{eq:nlp-analysis-structured-objectives}
\end{equation}
局部目标$\mathcal L_{\mathrm{local}}$把结构拆成一串决策$d_k$（逐位置标签、逐步动作、逐词元），每个决策单独做softmax交叉熵。它训练快，不需要对整个$\mathcal Y(x)$求和，但每个决策只与同一步的候选竞争，训练时看到的上下文又是参考结构的一部分，与推断时的整体$\argmax$并不一致。全局目标$\mathcal L_{\mathrm{global}}$是条件对数似然，第二项是对数\term{配分函数}（partition function），它让参考结构与全部合法结构竞争；对参数求导得到“参考结构的部件特征减去模型分布下的期望部件特征”，与第\ref{sec:nlp-classical-crf}小节式~\eqref{eq:nlp-classical-crf-gradient}的形式相同，代价是要能算出部件的边缘概率。\term{结构化间隔损失}（structured margin loss）$\mathcal L_{\mathrm{margin}}$只要求参考结构比每个错误结构高出与错误程度$\Delta$相称的间隔，例如$\Delta$取错误标签数或错误跨度数\scite{nlp-analysis-taskar2004,nlp-analysis-tsochantaridis2005}。计算它需要在得分上加入$\Delta$后再做一次$\argmax$，称为\term{代价增强解码}（cost-augmented decoding）；当$\Delta$能分解到部件上时，这只是在每个部件分数上加一个常数，原来的解码算法不变。取$\Delta\equiv0$并用次梯度逐样本更新，就退化为\term{结构化感知机}（structured perceptron）\scite{collins2002}。

用一个只有三个合法结构的教学例子比较后两种目标。设参考结构$y^\star$得分为2，错误结构$y_1$得分为2.5、与参考差一个部件（$\Delta=1$），$y_2$得分为0、差两个部件（$\Delta=2$）。全局目标为$-2+\log(e^{2}+e^{2.5}+e^{0})\approx1.02$，模型分布约为$p(y^\star)=0.36$、$p(y_1)=0.59$、$p(y_2)=0.05$；梯度按这些概率把三个结构的部件都推一推，$y_1$被压低得最多，$y_2$几乎不动。间隔目标为$\max\{2+0,\ 2.5+1,\ 0+2\}-2=1.5$，代价增强解码选出$y_1$，梯度只在$y^\star$与$y_1$不同的部件上起作用，$y_2$已经满足间隔而不受影响。结构化感知机看的是$\max\{2,2.5,0\}-2=0.5$，同样只纠正$y_1$，但不要求超出间隔。三者的差别在于参考结构要与谁竞争、赢多少才算够。

这三类目标对推断算法提出的要求不同，也正好对应动态规划的不同用法。同一张图表（格点、区间表）上把“取最大”换成“对数求和”，Viterbi变成前向算法，CKY变成inside算法；用代数语言说，两者是同一递推在不同\term{半环}（semiring）上的实例。因此凡是能精确求$\argmax$的分解，通常也能精确求配分函数，三类目标都可用；转移系统和自回归生成没有这样的图表，只能用局部目标或在束内近似全局目标。

\subsection{基于序列标注的预测}
\label{sec:nlp-analysis-paradigm-tagging}

当输出是一组互不重叠的区间时，可以把它编码为逐位置标签。分词用BMES：B为词首，M为词中，E为词尾，S为单字词；实体与组块用\term{BIO标注}（BIO tagging）：B-$X$开始一个类型为$X$的区间，I-$X$延续同类区间，O表示不在任何区间内；BIOES再用E-$X$标出区间末位、S-$X$标出单位置区间。编码是可逆的：给定合法标签序列，从左到右扫描即可恢复区间集合。可逆性只对合法序列成立，例如O之后紧跟I-PER时，就无法确定这个区间从哪里开始。

序列标注模型由三层组成。编码器输出经线性层得到\term{发射分数}（emission score）$\bm E=\bm H\bm W+\bm 1\bm b^\top\in\mathbb R^{n\times L}$，$L$为标签数，$\bm E_{i,a}$是位置$i$取标签$a$的分数；线性链CRF另有可学习的转移矩阵$\bm T\in\mathbb R^{L\times L}$和起止向量$\bm\pi,\bm\rho\in\mathbb R^{L}$；第三层是固定的\term{合法转移掩码}（legal transition mask）$\bm M\in\{0,-\infty\}^{L\times L}$及起止掩码$\bm m^{\mathrm{start}},\bm m^{\mathrm{end}}\in\{0,-\infty\}^{L}$，规范允许的转移取0，禁止的转移取$-\infty$。整句得分、训练目标和解码写成
\begin{equation}
  \begin{aligned}
    s(x,y)&=\bigl(\bm\pi+\bm m^{\mathrm{start}}\bigr)_{y_1}
      +\sum_{i=1}^{n}\bm E_{i,y_i}\\
      &\quad+\sum_{i=2}^{n}\bigl(\bm T+\bm M\bigr)_{y_{i-1},y_i}
      +\bigl(\bm\rho+\bm m^{\mathrm{end}}\bigr)_{y_n},\\
    \mathcal L_{\mathrm{CRF}}&=\log\sum_{y\in\{1,\ldots,L\}^n}\exp s(x,y)-s(x,y^\star),\\
    \hat y&=\argmax_{y\in\{1,\ldots,L\}^n}s(x,y).
  \end{aligned}
  \label{eq:nlp-analysis-legal-mask}
\end{equation}
求和与$\argmax$名义上遍历全部$L^n$条序列，但任何含非法转移的序列得分为$-\infty$，在配分函数中贡献$\exp(-\infty)=0$，也不可能成为$\argmax$；于是两者实际上都只在合法集合$\mathcal Y(x)$上进行，与式~\eqref{eq:nlp-analysis-structured-objectives}的$\mathcal L_{\mathrm{global}}$一致。线性链CRF的定义、前向—后向算法、梯度及其计算步骤见第\ref{sec:nlp-classical-crf}小节与算法~\ref{alg:nlp-classical-crf-gradient}，这里只需注意两点：掩码项是常数，对它的梯度为零；被屏蔽的转移在前向递推中自然消失，不需要改写算法。

掩码由标签规范直接构造。BIO规范下，对每个类型$X$，若$b=\mathrm{I}\text{-}X$而前驱$a\notin\{\mathrm B\text{-}X,\mathrm I\text{-}X\}$，则$\bm M_{a,b}=-\infty$，并令$\bm m^{\mathrm{start}}_{\mathrm I\text{-}X}=-\infty$；BIO没有句末限制。BIOES规范更严：B-$X$与I-$X$之后只能接I-$X$或E-$X$，E-$X$、S-$X$、O之后不能接I-$X$或E-$X$，句末只允许E-$X$、S-$X$、O。BMES是无类型的BIOES，其掩码表见第\ref{sec:nlp-analysis-bmes-crf}小节。掩码为什么必须同时用于训练和解码？若只在解码时屏蔽，训练的配分函数就包含了大量非法序列，模型要花参数把它们压低，而学到的转移分数只会是很大的有限负数，不会是$-\infty$；训练目标与解码时的输出空间也就不再相同。两处使用同一个掩码，训练、解码和标签到区间的恢复程序就共享同一个结构定义。

预训练编码器常把一个词切成多个子词，标签却定义在词或字上。常用的\term{首子词对齐}（first-subword alignment）只取每个词第一个子词的表示作为该词的$\bm h_i$，其余子词不产生标签、不计损失（第\ref{sec:nlp-classical-bert-heads}小节）。接CRF时，应先按首子词位置把$\bm E$收集成长度为词数的矩阵，再在其上做转移；若把非首子词也留在链上，相邻词之间就隔着若干无标签位置，转移矩阵描述的不再是相邻词的标签关系。解码得到词级标签后，按词到原文字符的偏移恢复字符跨度。

与逐位置softmax相比，带掩码的CRF多出两样东西：标签之间的依赖，以及合法性保证。逐位置softmax的代价是$O(nL)$，可以完全并行，但独立地取最大可能得到O后接I-$X$这样的序列，只能靠事后规则修补（例如把孤立的I改成B），而修补规则本身就是一种未经训练的决策。CRF的训练与解码代价为$O(nL^2)$，在神经编码器的代价面前通常可以忽略。序列标注的根本局限在于输出表示：一层标签只能表示互不重叠的平面区间，嵌套实体（“北京大学”是组织，其中“北京”是地点）需要多层标注，多个谓词的语义角色需要为每个谓词重复编码一次。把逐位置标签换成逐区间标签的半马尔可夫CRF可以直接给整段打分\scite{pgm-energy-models-sarawagi2004}，这已经接近下一小节的跨度范式。

\subsection{基于跨度的预测}
\label{sec:nlp-analysis-paradigm-span}

当输出由可以嵌套的区间组成时，与其把区间编码进逐位置标签，不如直接枚举候选区间并逐个评分。设候选集合为
\[
  \mathcal S=\{(i,j):1\leq i\leq j\leq n,\ j-i+1\leq W\},
\]
$W$为最大宽度，$|\mathcal S|=O(nW)$；不限宽度时为$n(n+1)/2$。每个候选需要一个固定维度的\term{跨度表示}（span representation）。Lee等在共指消解中提出的写法把三类信息拼接在一起\scite{lee2017coreference}：
\begin{equation}
  \begin{aligned}
    \alpha_t&=\bm w_\alpha^\top\operatorname{FFNN}_\alpha(\bm h_t),\qquad
    a_{i,j,t}=\frac{\exp\alpha_t}{\sum_{k=i}^{j}\exp\alpha_k},\qquad
    \hat{\bm h}_{i,j}=\sum_{t=i}^{j}a_{i,j,t}\,\bm h_t,\\
    \bm g_{i,j}&=\bigl[\bm h_i;\ \bm h_j;\ \hat{\bm h}_{i,j};\ \bm\phi(j-i)\bigr]\in\mathbb R^{3d+d_w}.
  \end{aligned}
  \label{eq:nlp-analysis-span-representation}
\end{equation}
两个端点向量$\bm h_i,\bm h_j$记录边界处的上下文，决定“从哪里开始、到哪里结束”；$\hat{\bm h}_{i,j}$是区间内部按注意力加权的汇聚，$\alpha_t$对所有跨度共享，可学成“这个位置像不像中心词”的分数，再在每个区间内部归一化；$\bm\phi(j-i)\in\mathbb R^{d_w}$是\term{宽度嵌入}（width embedding），把区间长度（通常分桶）映射成可学习向量，让模型知道人名多为两三字、长跨度较少是实体等先验。若编码器为双向循环网络，也可以用端点差分代替拼接，第\ref{sec:nlp-analysis-neural-cky}小节的神经成分解析就采用这种变体。

类别得分由前馈网络给出：$s(i,j,c)=\bm w_c^\top\operatorname{FFNN}(\bm g_{i,j})$，$c$取遍任务类别和一个表示“不是任何结构单元”的NONE类。训练目标是对每个候选做$C+1$类交叉熵，$C$为任务类别数，参考区间取其类别，其余候选取NONE。NONE类远多于正类，常见的处理是对NONE候选降采样或降低权重。另一种常用变体把起点和终点分别投影后用双仿射打分，$s(i,j,c)=\bm r_i^{\mathrm s\top}\bm U_c\bm r_j^{\mathrm e}+\bm w_c^\top[\bm r_i^{\mathrm s};\bm r_j^{\mathrm e}]+b_c$，一次矩阵运算就得到形状为$n\times n\times C$的分数张量\scite{nlp-analysis-yu2020biaffine}；它省去了内部汇聚，适合候选很多、跨度较短的场合。

候选规模是跨度范式的主要代价。式~\eqref{eq:nlp-analysis-span-representation}中每个跨度的注意力汇聚需$O(Wd)$，全部候选共$O(nW^2d)$；若后续还要给跨度对评分（共指中“提及$i$的先行项是否为$j$”），候选对数是$O(|\mathcal S|^2)$，不可直接承受。剪枝分两步：先用最大宽度$W$截断，再用一元分数$s_{\mathrm m}(i,j)=\bm w_{\mathrm m}^\top\operatorname{FFNN}_{\mathrm m}(\bm g_{i,j})$排序，只保留前$\lambda n$个跨度进入配对，配对数降为$O(\lambda^2n^2)$。剪枝分数与最终目标联合训练，其梯度只来自保留下来的跨度，被剪掉的正确跨度没有机会被纠正，因此$\lambda$要在召回与代价之间权衡。

跨度分类给出的是每个候选的独立判断，跨度之间的冲突还需要一步解码。若任务要求互不重叠（平面实体），按得分降序贪心选择不冲突的跨度最简单，但不保证总分最大；精确解是\term{加权区间调度}（weighted interval scheduling）：令$F(j)$为前$j$个位置上不重叠跨度的最大总分，
\[
  F(0)=0,\qquad F(j)=\max\Bigl\{F(j-1),\ \max_{i:\,(i,j)\in\mathcal S}\bigl[F(i-1)+\tilde s(i,j)\bigr]\Bigr\},
\]
其中$\tilde s(i,j)=\max_{c\neq\mathrm{NONE}}s(i,j,c)-s(i,j,\mathrm{NONE})$是选入该跨度相对不选的得分增量，只有正增量的跨度值得选入，递推代价为$O(nW)$。举一个教学数值：设$n=4$，候选$[1,3]$、$[1,2]$、$[3,4]$的增量分别为5、3、3.5，其余为负。贪心先选$[1,3]$，另两个都与它重叠，总分为5；区间调度比较$F(4)=\max\{F(3),F(2)+3.5\}$，其中$F(2)=3$、$F(3)=5$，得到$6.5$，选$[1,2]$与$[3,4]$。若任务允许嵌套但不允许部分交叉（成分树、层次实体），则要用第\ref{sec:nlp-analysis-paradigm-graph}小节的CKY在跨度分之上求最优树；若嵌套和交叉都允许（共指中的提及），则不需要冲突解码。

与序列标注相比，跨度范式用$O(nW)$个候选换来了嵌套表示能力，并能把整段作为一个对象建模（宽度、内部汇聚），这正是共指、成分解析和嵌套实体需要的。它的代价是候选规模与类别不平衡，以及平面任务中额外的冲突解码。

\subsection{基于转移系统的预测}
\label{sec:nlp-analysis-paradigm-transition}

前两种范式都在固定的候选集合上打分。另一种思路是模仿人逐词读句、边读边搭结构的过程：把结构的构建写成一台状态机，每一步选一个动作。形式上，一个\term{转移系统}（transition system）由四部分组成\scite{nlp-analysis-nivre2008}：状态集合$\mathcal C$，每个状态$c$通常由栈$\sigma$、缓冲区$\beta$和已建部分$A$组成；动作集合$\mathcal A$，每个动作$a$有一个前提（在哪些状态下可执行）和一个效果（把$c$变成$a(c)$）；初始函数$c_{\mathrm s}(x)$，通常是空栈、整句在缓冲区、$A=\varnothing$；终止状态集合$\mathcal C_{\mathrm t}$，从中可以读出完整结构。以arc-standard依存系统为例，SHIFT把缓冲区首词压栈，LEFT-ARC在栈顶两词间加一条由栈顶指向次顶的弧并弹出次顶，RIGHT-ARC加一条反方向的弧并弹出栈顶。“小王／阅读／通知”经SHIFT、SHIFT、LEFT-ARC、SHIFT、RIGHT-ARC，得到“阅读”支配“小王”和“通知”两条弧，再连到根即完成；完整的动作—前提—效果表与轨迹见第\ref{sec:nlp-analysis-arc-transition}小节。

模型在每个状态上给合法动作打分。状态特征$\bm f(c)$由栈顶与缓冲区若干位置的编码器表示及已建部分的信息拼成，动作分布只在合法动作集合$\mathcal A(c)$上归一化：
\[
  p_{\bm\theta}(a\mid c)=\frac{\exp\bigl(\bm w_a^\top\bm f(c)\bigr)}{\sum_{a'\in\mathcal A(c)}\exp\bigl(\bm w_{a'}^\top\bm f(c)\bigr)},\qquad a\in\mathcal A(c),
\]
不满足前提的动作直接屏蔽，与序列标注中的合法掩码作用相同。训练需要把参考结构$y^\star$变成动作监督。\term{静态oracle}（static oracle）是一个确定性规则，它从初始状态出发，每一步按固定优先级选出一个能导向$y^\star$的动作，得到参考动作序列$a_1^\star,\ldots,a_m^\star$及沿途状态$c_1^\star=c_{\mathrm s}(x)$、$c_{t+1}^\star=a_t^\star(c_t^\star)$。训练目标是这些动作的局部交叉熵：
\begin{equation}
  \mathcal L_{\mathrm{static}}(\bm\theta)=-\sum_{t=1}^{m}\log p_{\bm\theta}\bigl(a_t^\star\mid c_t^\star\bigr).
  \label{eq:nlp-analysis-transition-loss}
\end{equation}
它是式~\eqref{eq:nlp-analysis-structured-objectives}中$\mathcal L_{\mathrm{local}}$在转移系统上的具体形式。推断时从$c_{\mathrm s}(x)$出发，每步执行概率最高的合法动作直到终止；arc-standard对$n$个词恰好执行$2n$个动作，所以推断时间与句长成线性。同一结构常有多条动作序列都能得到它，这称为\term{伪歧义}（spurious ambiguity）；静态oracle只挑其中一条，其余同样正确的动作在训练中被当作错误。

式~\eqref{eq:nlp-analysis-transition-loss}的状态$c_t^\star$全部来自参考历史，模型从未见过自己犯错之后的状态。推断时一旦某一步选错，后续状态就偏离训练分布，模型在陌生状态上更容易继续出错，错误逐步累积，这就是暴露偏差（exposure bias）。\term{动态oracle}（dynamic oracle）针对这个问题\scite{nlp-analysis-goldberg2012}：对任意状态$c$，不论它是否在参考轨迹上，都给出最优动作集合
\[
  \mathcal O(c;y^\star)=\Bigl\{a\in\mathcal A(c):\ \min_{y\in\mathcal Y(a(c))}\operatorname{loss}(y,y^\star)=\min_{y\in\mathcal Y(c)}\operatorname{loss}(y,y^\star)\Bigr\},
\]
其中$\mathcal Y(c)$是从$c$出发仍可到达的结构集合，$\operatorname{loss}$为错误弧数等结构损失。也就是说，最优动作不会让最好可达结果变差。对arc-eager这类“弧可分解”的系统，动作的代价等于它使之不可达的参考弧数，可以在常数或线性时间内算出。训练时让模型按自己的预测走（可设一个探索概率，并在若干轮之后才开始探索），在每个访问到的状态上提高最优动作集合的总概率：
\begin{equation}
  \mathcal L_{\mathrm{dynamic}}(\bm\theta)=-\sum_{t}\log\sum_{a\in\mathcal O(c_t;y^\star)}p_{\bm\theta}(a\mid c_t),\qquad c_{t+1}=\hat a_t(c_t).
  \label{eq:nlp-analysis-dynamic-oracle-loss}
\end{equation}
算法~\ref{alg:nlp-analysis-dynamic-oracle}列出这一过程。对集合求和同时处理了伪歧义：多个同样好的动作都不受惩罚。从模仿学习的角度看，这与DAgger的思想一致，都是在学习者自己产生的状态分布上向专家查询正确动作\scite{ross2011}；序列生成中的计划采样（scheduled sampling）也按同一思路逐步用模型预测替换参考前缀\scite{nlp-analysis-bengio2015scheduled}。

\Needspace{14\baselineskip}
\begin{algorithm}[基于动态oracle的转移系统训练]
  \label{alg:nlp-analysis-dynamic-oracle}
  \begin{algorithmic}[1]
    \Require 训练集$\{(x,y^\star)\}$，转移系统，可计算$\mathcal O(c;y^\star)$的动态oracle，探索概率$q$，开始探索的轮次$k$
    \Ensure 参数$\bm\theta$
    \For{轮次$e=1,2,\ldots$}
      \For{每个样本$(x,y^\star)$}
        \State $c\gets c_{\mathrm s}(x)$；$\mathcal L\gets0$
        \While{$c\notin\mathcal C_{\mathrm t}$}
          \State $\mathcal O\gets\mathcal O(c;y^\star)$ \Comment{当前状态下不增加损失的动作}
          \State $\mathcal L\gets\mathcal L-\log\sum_{a\in\mathcal O}p_{\bm\theta}(a\mid c)$
          \State $\hat a\gets\argmax_{a\in\mathcal A(c)}p_{\bm\theta}(a\mid c)$
          \If{$e>k$ 且以概率$q$探索}
            \State $c\gets\hat a(c)$ \Comment{按模型预测走，可能离开参考轨迹}
          \Else
            \State $c\gets\bigl(\argmax_{a\in\mathcal O}p_{\bm\theta}(a\mid c)\bigr)(c)$
          \EndIf
        \EndWhile
        \State 用$\nabla_{\bm\theta}\mathcal L$更新$\bm\theta$
      \EndFor
    \EndFor
  \end{algorithmic}
\end{algorithm}

贪心推断的另一个弱点是早期决策无法撤销。束搜索（第\ref{sec:nlp-generation-beam}小节）在每步保留$b$个得分最高的部分动作序列，可以推迟决定。但若模型逐步局部归一化，束搜索仍受第\ref{sec:nlp-classical-memm}小节所说的标注偏置影响：每个状态出发的概率之和恒为1，后面的证据无法降低一条早期路径的总分。Andor等因此对整条动作序列做全局归一化，$p(a_{1:m}\mid x)\propto\exp\sum_t s(c_t,a_t)$，配分函数在束内的候选上近似，并在参考序列掉出束时立即停止该样本、按当时的束更新参数\scite{nlp-analysis-andor2016}；这种\term{早期更新}（early update）最早用于感知机训练的增量解析\scite{nlp-analysis-collins2004}。

转移系统的优点是推断线性、可以自由使用已建结构作为特征（例如已经找到的依存子节点），并能处理难以写成固定图表的结构，如AMR图和篇章树。代价是没有全局最优保证，结果依赖动作顺序，且每个新结构类型都要单独设计动作集合并证明其完备性（任何合法结构都可达）与可靠性（只产生合法结构）。

\subsection{基于图与区间动态规划的预测}
\label{sec:nlp-analysis-paradigm-graph}

若结构得分可以分解到弧或区间上，就能在规定的结构空间内精确求最优。依存树的\term{弧分解}（arc-factored）模型令$s(x,y)=\sum_{(h\to m)\in y}s(h,m)$，$s(h,m)$为词$h$支配词$m$的分数，位置0为人工根；成分树的跨度分解模型令$s(x,y)=\sum_{(i,j,\ell)\in y}s(i,j,\ell)$。这两种分解对应三种经典算法，它们都按“子结构的最优值只依赖子结构的边界”这一性质组织递推。

CKY在区间上递推二叉树。当分数来自跨度而非文法规则时，记$\bar s(i,j)=\max_\ell s(i,j,\ell)$，
\begin{equation}
  \pi(i,i)=\bar s(i,i),\qquad
  \pi(i,j)=\bar s(i,j)+\max_{i\leq k<j}\bigl[\pi(i,k)+\pi(k+1,j)\bigr],
  \label{eq:nlp-analysis-span-cky}
\end{equation}
$\pi(1,n)$是最优树的总分，记录每个区间取得最大的$k$和$\ell$即可回溯。共有$O(n^2)$个区间，每个区间枚举$O(n)$个分裂点，总代价$O(n^3+n^2L)$。若改用PCFG的规则概率评分，状态再带上非终结符，代价为$O(n^3|G|)$，其inside形式和期望计数见第\ref{sec:nlp-classical-cky}小节与第\ref{sec:nlp-classical-inside-outside}小节。

Eisner算法在投射依存树上做同样的事\scite{nlp-analysis-eisner1996}。直接把“以$h$为中心词的区间”作为状态需要同时记住区间两端和中心词，代价为$O(n^5)$；Eisner的做法是让中心词总在区间的一端，并区分“完整”与“不完整”区间。$C^{\to}[i,j]$表示$i$支配$(i,j]$内全部词、且$j$一侧已不再接收子节点的最优子树，$I^{\to}[i,j]$表示刚刚加入弧$i\to j$、$j$的内侧子节点已经收齐的区间，$C^{\leftarrow}$与$I^{\leftarrow}$为方向相反的对应项。初值$C^{\to}[i,i]=C^{\leftarrow}[i,i]=0$，递推为
\begin{equation}
  \begin{aligned}
    I^{\to}[i,j]&=\max_{i\leq k<j}\bigl(C^{\to}[i,k]+C^{\leftarrow}[k+1,j]\bigr)+s(i,j),\\
    I^{\leftarrow}[i,j]&=\max_{i\leq k<j}\bigl(C^{\to}[i,k]+C^{\leftarrow}[k+1,j]\bigr)+s(j,i),\\
    C^{\to}[i,j]&=\max_{i<k\leq j}\bigl(I^{\to}[i,k]+C^{\to}[k,j]\bigr),\\
    C^{\leftarrow}[i,j]&=\max_{i\leq k<j}\bigl(C^{\leftarrow}[i,k]+I^{\leftarrow}[k,j]\bigr).
  \end{aligned}
  \label{eq:nlp-analysis-eisner}
\end{equation}
每个状态只有两个端点加一个方向，共$O(n^2)$个，每个枚举$O(n)$个分裂点，总代价$O(n^3)$。投射性来自递推本身：两个子区间只在端点相接，交叉弧无法拼出来。人工根只连一个真实词的单根约束、数值算例与回溯见第\ref{sec:nlp-analysis-eisner}小节。

允许交叉弧时，依存树就是有向图上以根为起点的最大生成树（最大树形图）。Chu–Liu/Edmonds算法先为每个非根词选分数最高的入边\scite{nlp-analysis-chu1965,nlp-analysis-edmonds1967}；若这些边不成环，它们就是答案；若成环，把环$C$收缩为一个节点，进入环中节点$v$的边$u\to v$改记分数$s(u,v)-s(\pi(v),v)$，$\pi(v)$为$v$在环内的入边起点，离开环的边取环内各节点出边的最大分数，在收缩图上递归求解，再展开环、去掉被替换的那条环内边。朴素实现为$O(n^3)$，Tarjan的实现在稠密图上为$O(n^2)$\scite{nlp-analysis-tarjan1977}。McDonald等把它用于依存分析，在一阶分解下，非投射解码因此反而比投射的Eisner算法便宜\scite{v9a-mcdonald2005}。算法框、单根约束与算例见第\ref{sec:nlp-analysis-mst}小节。

三种算法都依赖一阶可加分解。若要让分数依赖兄弟弧、祖父弧等高阶部件，状态就要多记信息：投射情形下加入相邻兄弟弧仍可保持$O(n^3)$，非投射情形的二阶精确解码却已是NP难问题，需要近似\scite{nlp-analysis-mcdonald2006}；再加入祖父弧，投射解码升到$O(n^4)$\scite{nlp-analysis-koo2010}。训练上，三类目标都可用：CKY与Eisner把$\max$换成对数求和即得inside算法，可计算全局似然和边缘概率；非投射树的配分函数可由矩阵树定理用一个行列式求出\scite{nlp-analysis-koo2007}；结构化间隔损失只需代价增强解码。神经图解析器常采用更简单的局部目标：每个词在所有候选中心词上做一次softmax，按交叉熵训练，测试时再用Eisner或Chu–Liu/Edmonds恢复合法树\scite{nlp-analysis-dozat2017}。这时训练不考虑树约束，由解码修正少量成环或多根的情况；第\ref{sec:nlp-analysis-biaffine}小节的双仿射解析器即采用这种组合。

图方法与转移系统针对同一类结构，取舍正好相反：图方法在分解范围内给出精确最优，但部件只能是局部的，代价至少为$O(n^2)$；转移系统线性时间、特征可以读取已建的任意部分，但没有全局保证。

\subsection{基于约束生成的预测}
\label{sec:nlp-analysis-paradigm-generation}

前四种范式都要为每类结构设计专门的标签、候选或动作。生成范式改为把结构写成一个词元序列，用统一的自回归模型生成它。这需要一个\term{线性化}（linearization）映射$\operatorname{lin}:\mathcal Y(x)\to\mathcal V^*$，它必须可逆：给定合法序列，解析程序$\operatorname{lin}^{-1}$能唯一恢复结构。成分树可以写成括号序列“(S (NP 小王) (VP 阅读 (NP 通知)))”，语义图可以写成带变量和重入标记的括号式，实体与关系可以写成“[小王｜人名]”这类结构化抽取语言，查询直接写成SQL。模型按
\[
  p_{\bm\theta}(z\mid x)=\prod_{t=1}^{m}p_{\bm\theta}(z_t\mid z_{<t},x),\qquad z=\operatorname{lin}(y)
\]
给序列$z$打分，用教师强制的逐词元交叉熵训练（第\ref{sec:nlp-sft-masked-loss}小节），即在参考前缀$z^\star_{<t}$上预测$z^\star_t$；这是$\mathcal L_{\mathrm{local}}$的又一种形式，因而也有与转移系统相同的暴露偏差。可以把转移系统看成线性化的一种特殊情形：动作序列就是一种线性化，只是动作的合法性由状态机直接给出。

生成的输出必须能落回原文与结构定义，这要靠两类机制。第一类是复制：结构中的词语、实体名大多来自输入，与其从整个词表中生成，不如让解码器直接指向输入位置。复制分布与词表分布的混合形式在第\ref{sec:nlp-analysis-seq2seq-lemma}小节式~\eqref{eq:nlp-analysis-copy}首次写出；更彻底的做法是让目标序列只含位置指针和类型标记，解析时由指针直接恢复字符偏移。第二类是约束解码：在每一步只允许能延续成合法输出的词元，约束后的分布与掩码计算见第\ref{sec:nlp-inference-fsm}小节与第\ref{sec:nlp-inference-grammar}小节，前者用有限状态机处理正则格式，后者用下推自动机处理括号嵌套。在结构分析中，掩码有三个来源：线性化文法本身（括号配对、每个开括号后必须是类别符号）；输入（复制的片段必须是原文子串，可用原文的前缀树检查）；外部模式（SQL中的列名必须属于给定数据库，实体名必须在知识库名称表中）。以图~\ref{fig:nlp-analysis-paradigms}最后一行为例，前缀已生成“(S (NP 小王) (VP 阅读 (NP 通知”，此时原文已读完，括号文法只允许生成闭括号；若不施加约束，模型可能在这里再生成一个原文中不存在的词。

线性化还引入了前四种范式没有的偏差。集合型结构（一句中的多个实体、一个事件的多个论元）本无顺序，线性化必须规定一个顺序，参考序列的顺序就成了额外的监督，模型生成了正确集合的另一种排列也会被交叉熵惩罚；常见做法是固定一种规范顺序（如按原文位置）并在评价时按集合比较。自回归解码也没有全局最优保证，束搜索只是近似；输出还须经$\operatorname{lin}^{-1}$解析并与原文对齐核验，解析失败、指向不存在的位置或复制了不在原文中的片段，都应计为错误，而不是被静默修补。

生成范式的优势在于接口统一：同一个预训练编码器—解码器或大语言模型可以覆盖树、图、逻辑式和抽取记录，新增结构类型只需新增线性化与文法，并能直接利用预训练中学到的生成能力；大语言模型的提示式分析也属于这一范式，只是不更新参数。代价是解码串行、长结构的序列较长，以及上面所说的顺序偏差和核验负担。

\subsection{范式的比较与选择}
\label{sec:nlp-analysis-paradigm-comparison}

表~\ref{tab:nlp-analysis-paradigms}把五种范式放在同一组维度上。读表时可以沿两条线比较：从上到下，输出表示越来越灵活，从平面区间到任意可线性化的结构；推断的精确性却没有单调关系，序列标注和图方法在各自的分解内精确，跨度范式的冲突解码视任务而定，转移系统和生成只有近似保证。表中最后一列列出本章及后两章的主要任务所采用的范式；同一任务出现在多行，说明范式是建模选择，而非由任务唯一决定。

\begin{table*}[htbp]
  \centering
  \small
  \begin{tabularx}{\linewidth}{@{}>{\raggedright\arraybackslash}p{13mm}>{\raggedright\arraybackslash}p{26mm}>{\raggedright\arraybackslash}p{22mm}>{\raggedright\arraybackslash}p{24mm}>{\raggedright\arraybackslash}p{18mm}>{\raggedright\arraybackslash}X@{}}
    \toprule
    范式 & 输出表示与得分分解 & 推断 & 训练目标 & 复杂度 & 典型任务\\
    \midrule
    序列标注 & 逐位置标签；发射分＋相邻转移分 & 带合法掩码的Viterbi & 逐位置交叉熵，或CRF条件似然 & $O(nL^2)$ & 本章：分词、词性与组块、语义角色BIO、话语行为；第\ref{chap:nlp-information-extraction}章：BIO实体、触发词、统一标注的关系抽取；第\ref{chap:nlp-classification-matching-organization}章：标注式关键短语抽取\\
    跨度 & 候选区间及类别；跨度分（可加跨度对分） & 枚举与剪枝；不重叠选择、CKY或直接阈值 & 含NONE类的跨度交叉熵 & 候选$O(nW)$；配对$O(\lambda^2n^2)$ & 本章：神经成分解析、跨度语义角色、跨度共指；第\ref{chap:nlp-information-extraction}章：嵌套实体、实体对关系分类；第\ref{chap:nlp-classification-matching-organization}章：跨度配对的方面情感三元组\\
    转移系统 & 动作序列；状态上的动作分 & 贪心或束搜索 & oracle动作交叉熵；动态oracle；束内全局归一化 & $O(n)$个动作 & 本章：移进–归约成分解析、arc-standard与arc-eager、转移式AMR、RST篇章解析\\
    图与区间动态规划 & 弧或区间集合；一阶可加分解 & CKY、Eisner、Chu–Liu/Edmonds & 全局似然、结构化间隔，或逐词选中心词 & $O(n^3)$；CLE可达$O(n^2)$ & 本章：CKY成分解析、投射与非投射依存、双仿射、增强依存与语义图；第\ref{chap:nlp-information-extraction}章：超图与词对网格实体、词元对关系标记；第\ref{chap:nlp-classification-matching-organization}章：网格标注的方面情感三元组\\
    约束生成 & 线性化词元序列；逐词元条件概率 & 束搜索加文法、前缀树或模式掩码 & 教师强制交叉熵 & $m$步串行解码，掩码代价另计 & 本章：序列化解析、AMR生成、逻辑式与SQL、话语改写；第\ref{chap:nlp-information-extraction}章：指针式与结构化生成抽取、生成式实体链接、复制式三元组；第\ref{chap:nlp-classification-matching-organization}章：指针生成的方面情感三元组、关键短语生成\\
    \bottomrule
  \end{tabularx}
  \caption{结构预测的五种范式。$n$为分析单位数，$L$为标签数，$W$为最大跨度宽度，$\lambda$为剪枝比例，$m$为输出序列长度；复杂度只计结构层，编码器代价另计。第\ref{chap:nlp-classification-matching-organization}章的文本分类与匹配输出单个类别或分数，相当于$\mathcal Y(x)$为有限标签集的退化情形，不需要结构解码；该章的方面情感三元组与关键短语抽取仍是结构输出，列在相应行中。}
  \label{tab:nlp-analysis-paradigms}
\end{table*}

选择范式时，可以依次检查四个问题。第一，结构是否允许嵌套或交叉：只有平面区间时序列标注最经济；有嵌套而无交叉时用跨度加CKY或分层解码；有交叉弧或多入边时用图方法或生成。第二，是否需要精确推断与可校准的概率：需要边缘概率（例如把分析结果传给下游做不确定性处理）时，优先选能算配分函数的标注或图方法。第三，监督是结构本身还是只有最终答案：只有执行结果或答案时（语义解析中的答案监督），结构是隐变量，生成范式配合采样或边缘似然更容易训练。第四，长度与时延：长文档上$O(n^3)$的图方法和逐词元生成都可能过慢，转移系统或带剪枝的跨度模型更合适。

同一任务通常有多种实现。依存分析有转移、图和生成三种做法；实体识别有序列标注、跨度、超图、词对网格和生成五种做法。后文在各小节开头注明所采用的范式；比较不同范式时，应固定编码器、训练数据与评价口径，只改变输出表示与解码，否则观察到的差异无法归因于范式本身。

\section{词法分析与局部标注}
\label{sec:nlp-analysis-lexical}

“机器学习读书会”可以作为一个活动名称，词性标注却可能需要“机器学习／读书会”两个词，模型输入又可能把每个汉字作为一个词元。\term{词法分析}（lexical analysis）处理词的边界、基本形式、形态属性、词性和浅层短语，它必须先确定采用什么语言单位，再给这些单位赋属性。本节的共同挑战有三项：分析单位本身需要预测，且语言词、模型词元与原文字符三者的边界不一致；同一表面形式在不同语境中有不同切分、基本形式或类别；每个局部决定都受相邻决定约束，例如词内的BMES标签必须连贯、同一词内的字必须属于同一词性。

三个小节分别处理边界、形态和局部类别。分词与词性、组块主要使用第\ref{sec:nlp-analysis-paradigm-tagging}小节的序列标注范式，合法转移掩码在这里有最直接的用途；词形还原把变化规则当作类别做分类，或用字符级生成直接写出基本形式，属于第\ref{sec:nlp-analysis-paradigm-generation}小节的生成范式。所有输出最终都恢复为原文字符跨度，再按所选规范评价。

图~\ref{fig:nlp-analysis-alignment}把位置作为各层共同的参照。图中的模型切分只是教学设定，不对应某个实际词表。分词和组块可以有不同边界，模型词元也不必服从其中任何一种边界。

\begin{figure}[htbp]
  \centering
  \begin{tikzpicture}[
  x=1mm,y=1mm,
  every node/.append style={font=\figurefont\footnotesize,align=center},
  unit/.style={draw=FigureOutput,fill=FigureOutput!10,line width=.85pt,
    minimum height=9mm,inner sep=1mm},
  mapping/.style={->,draw=NNLine,line width=.85pt}
]
  \node[anchor=east] at (21,17) {模型词元};
  \node[anchor=east] at (21,0) {原文字符};
  \node[anchor=east] at (21,-17) {语言词};
  \node[anchor=east] at (21,-31) {BMES};
  \node[anchor=east] at (21,-42) {词级组块};
  \node[unit,draw=FigureModel,fill=FigureModel!10,minimum width=51mm]
    at (50,17) {研究生 \quad $[0,3)$};
  \node[unit,draw=FigureModel,fill=FigureModel!10,minimum width=17mm]
    at (84,17) {命};
  \foreach \x/\char/\offset in {33/研/0,50/究/1,67/生/2,84/命/3} {
    \node[unit,draw=FigureInput,fill=FigureInput!10,minimum width=17mm]
      at (\x,0) {\char};
    \node[font=\figurefont\scriptsize,anchor=south] at (\x,4.8) {\offset};
    \draw[mapping] (\x,12.5) -- (\x,9);
    \draw[mapping] (\x,-12.5) -- (\x,-4.5);
  }
  \node[unit,minimum width=34mm] at (41.5,-17) {研究 \quad $[0,2)$};
  \node[unit,minimum width=34mm] at (75.5,-17) {生命 \quad $[2,4)$};
  \foreach \x/\tag in {33/B,50/E,67/B,84/E}
    \node at (\x,-31) {\tag};
  \node[unit,minimum width=34mm] at (41.5,-42) {B-VP};
  \node[unit,minimum width=34mm] at (75.5,-42) {B-NP};
  \node[anchor=north,align=left] at (58.5,-50)
    {组块：研究/VP，生命/NP\\箭头：单位覆盖的原文字符};
\end{tikzpicture}

  \caption{原文“研究生命”的字符、语言词、模型词元与标签对应。横向位置均指原文字符；箭头只表示覆盖位置，不表示预测依赖。语言词采用“研究／生命”，模型词元故意采用“研究生／命”。本例把“研究”视为动词、“生命”视为名词，分别形成VP与NP组块；词元标签必须经过位置映射才能恢复这些语言单位。}
  \label{fig:nlp-analysis-alignment}
\end{figure}

\subsection{分词与边界确定}
\label{sec:nlp-analysis-segmentation}

\term{分词}（word segmentation）按照选定规范把字符序列切成词语。中文分词的困难可以归为四项。第一，书写不用空格标出词界，边界本身要预测。第二，歧义有两种形态：字串“研究生命”中“研究生”与“生命”都是词且共享“生”，这类相邻候选词互相重叠的情形称为\term{交集歧义}（overlapping ambiguity）；“他将来北京”中的“将来”在别处是一个词，这里却应切成“将／来”，这类同一字串既可整体成词、又可拆开的情形称为\term{组合歧义}（combinational ambiguity）。第三，人名、机构名、领域术语和新词不断出现，训练语料或词典中没有的\term{未登录词}（out-of-vocabulary word, OOV）是错误的主要来源。第四，各语料的分词规范不一致，同一句话在不同语料中的“正确”切分并不相同。

语言学分词与模型的子词切分目标不同。BPE、WordPiece和Unigram分词模型按压缩或似然组织模型输入（第\ref{sec:nlp-llm-tokenization}小节），不会自动生成所需的语言词界。固定“研究生命”时，语言词可为“研究／生命”，模型词元可为“研究生／命”；后者跨过了语言词界，一个词元的类别不能直接复制给一个完整词。

方法沿以下谱系推进，每一步回应前一方法的一个缺陷。词典最大匹配直接利用词表，规则透明，但无法用语境消歧，也不能发现新词；HMM把分词改写为字符边界标注，用统计代替规则，新词可以由已见的边界模式组合出来，但发射只依赖当前字；CRF允许使用任意重叠的上下文特征，并对整句全局归一化；预训练编码器以上下文表示替代人工特征，提高了未登录词召回；词格模型把被字符模型丢掉的词典信息重新接入；多标准分词让规范不同的语料共享一个模型。

评价把预测词和参考词都恢复为原文跨度集合$P$与$G$，精确率、召回率分别为$|P\cap G|/|P|$与$|P\cap G|/|G|$，F1为二者的调和平均。未登录词召回只在参考词中不属于训练词表$V$的那部分上计算，即$|\{w\in P\cap G:w\notin V\}|/|\{w\in G:w\notin V\}|$，它单独反映新词处理能力。预测“研究生／命”得到$\{[0,3),[3,4)\}$，参考“研究／生命”为$\{[0,2),[2,4)\}$，严格匹配数为0，两个指标都是0。

预训练模型普遍以字或子词为输入后，还需要显式分词吗？对许多端到端任务，例如以字为单位的实体识别或文本分类，模型可以不经分词直接学习，分词错误也就不会传给下游。显式的词界仍有三类用途：需要词作为计数或检索单位的系统（倒排索引、词频统计、术语抽取）；需要语言学单位的分析（词性、依存树都定义在词上）；以及需要人能检查的中间结果。分词的价值因此取决于下游是否消费词这一单位，而不取决于模型能否在内部绕开它。

\subsubsection{词典最大匹配与候选切分图}
\label{sec:nlp-analysis-dictionary}

设词典包含“研究”“研究生”“生命”“命”，并允许单字兜底。\term{正向最大匹配}（forward maximum matching）从左边当前位置出发，尝试不超过词典最大词长的片段，取最长的已收录词，再移动到其右端，“研究生命”因此得到“研究生／命”。逆向最大匹配从右端寻找最长词，先选“生命”，再选“研究”，得到另一种切分。两种方法都只做局部选择；双向比较还需另行规定冲突时按词数、单字数或其他证据决定。

要同时比较整句候选，可把字符间的边界$0,\ldots,n$作为节点。若$x[s:e]$在词典中，就加入边$s\to e$，边上记该词及其代价$c(s,e)$。所有边都向右，构成有向无环图。令$d(e)$为从0到$e$的最小总代价，递推为
\begin{equation}
  d(0)=0,\qquad
  d(e)=\min_{s:\,s\to e}\{d(s)+c(s,e)\}.
  \label{eq:nlp-analysis-word-path}
\end{equation}
其余初值为$+\infty$，每次保存最优前驱$s$。到达$n$后逆向读取前驱，把各条边的$[s,e)$反转为正序，就得到词跨度；若$d(n)=+\infty$，应报告无合法切分，或明确启用单字兜底。代价取负对数词概率时，最短路径就是词概率乘积最大的切分，即一元语言模型下的最优切分。

设“研究”“生命”“研究生”“命”的代价依次为$1,1,0.7,2$，单字兜底代价均为3（教学设定）。于是$d(2)=1$，$d(3)=0.7$；到达4时，经“生命”的代价是$1+1=2$，经“命”的代价是$0.7+2=2.7$，其余路径更大。回溯$4\leftarrow2\leftarrow0$得到“研究／生命”。正向最大匹配先取了最长的“研究生”，全局路径则因后续代价而放弃它，这正是交集歧义需要整句比较的原因。

这类方法的知识全在词表和代价里：词表缺少“生命”会直接删掉正确候选，错误的长词会使贪心提前走错，两种故障应分别靠补候选和调代价修复。它适合词表稳定、要求能解释词项来源的受限文本，也常作为统计方法的候选来源或特征来源；面对组合歧义和未登录词，它没有可用的语境信息。

\subsubsection{HMM边界标注}
\label{sec:nlp-analysis-hmm-boundary}

词典不能枚举所有新词，字符在词内的位置却可以在不同词之间共享。\term{BMES标注}（BMES tagging）用B表示词首、M表示词中、E表示词尾、S表示单字词；“研究／生命”转成B、E、B、E，“研究生／命”转成B、M、E、S。分词由此成为字符序列标注，一个未见过的词仍可由已见过的边界模式组成\scite{nlp-analysis-xue2003}。

HMM把边界标签$y_i$当作隐藏状态、字符$x_i$当作观测，模型定义、前向算法和Viterbi见第\ref{sec:nlp-classical-hmm}小节与第\ref{sec:nlp-classical-viterbi}小节；这里只说明参数从何而来。由已分词文本转换出的BMES序列，统计句首状态次数$N_{\mathrm{start}}(a)$、相邻状态次数$N(a,b)$和状态发射字符次数$N(a,c)$，用加$\alpha$平滑在允许后继集合$A(a)$上估计
\[
  p(b\mid a)=
  \frac{N(a,b)+\alpha}
       {\sum_{b'\in A(a)}N(a,b')+\alpha|A(a)|}.
\]
发射概率按状态分别归一化，并在字符表中保留一个未知字符符号。非法转移（如B后接S）概率恒为0，不参与平滑分母，这相当于HMM版本的合法掩码。

两条训练记录“研究／生命”和“研究生／命”贡献B到E两次、B到M一次，因此不加平滑时$p(\mathrm E\mid\mathrm B)=2/3$、$p(\mathrm M\mid\mathrm B)=1/3$。第一条记录中“生”贡献B发射，第二条中“生”贡献E发射：同一字符可以处在不同的词内位置，这一不确定性要靠转移和相邻字符共同消解。设一个已平滑的模型给两条候选路径的连乘因子为（教学参数，依次为起始、发射、转移与发射交替及句末转移）
\[
\begin{aligned}
 p(x,\mathrm{BEBE})&=
 .8\cdot.5\cdot.6\cdot.3\cdot.6\cdot.3
 \cdot.6\cdot.4\cdot.2=0.00062208,\\
 p(x,\mathrm{BMES})&=
 .8\cdot.5\cdot.4\cdot.5\cdot.9\cdot.2
 \cdot.2\cdot.2\cdot.8=0.0004608,
\end{aligned}
\]
BEBE胜出，恢复为$[0,2)$、$[2,4)$。要确认整句最优，仍须让Viterbi覆盖所有合法路径，而不是只比较两条手选候选。

HMM的局限来自它的生成式假设：给定状态，字符的出现与其他字符无关。“生”是否是词尾，强烈依赖右邻字是“命”还是“活”，但HMM只能通过状态转移间接利用这一信息，无法直接使用“右邻字是什么”这类特征。它的优点同样来自这一简单结构：参数只有$4\times4$的转移表和每个状态上的字符分布，训练只需一遍计数，解码为$O(4^2n)$，平滑后的参数可以逐项检查。在标注很少或需要快速得到基线时，这些性质仍然有用。

\subsubsection{BMES--CRF边界标注与Viterbi解码}
\label{sec:nlp-analysis-bmes-crf}

CRF改为直接为“位置$i$取标签$a$”评分，评分可以读取整句。发射分数$\bm E_{i,a}$在传统做法中是特征的线性组合，特征模板包括窗口内的单字与相邻字对（$x_{i-2},\ldots,x_{i+2}$，$x_{i-1}x_i$，$x_ix_{i+1}$）、字符类型（汉字、数字、字母、标点），以及从训练语料统计的构词与未登录词特征\scite{v9a-tseng2005}；在神经做法中是编码器输出经线性层的结果。标签之间的依赖由转移矩阵$\bm T$学习，合法性由固定掩码规定。BMES的合法转移与起止条件为
\[
\begin{array}{c|c}
 \text{位置或前驱}&\text{允许标签或后继}\\ \hline
 \text{句首}&\mathrm{B,S}\\
 \mathrm{B}&\mathrm{M,E}\\
 \mathrm{M}&\mathrm{M,E}\\
 \mathrm{E}&\mathrm{B,S}\\
 \mathrm{S}&\mathrm{B,S}\\
 \text{句末}&\mathrm{E,S}
\end{array}
\]
特征模板怎样变成可学习的参数，可以看“研究生命”的第3个字“生”。窗口特征取值为$x_{i-2}=$研、$x_{i-1}=$究、$x_i=$生、$x_{i+1}=$命、$x_{i+2}=$句末符，字对特征为“究生”“生命”以及跳过当前字的“究命”，字符类型为汉字。每个特征值都与当前标签组合成一个指示特征，例如$\mathbb I[x_ix_{i+1}=\text{生命}\wedge y_i=\mathrm B]$；发射分数是这些指示特征的权重之和，$\bm E_{i,a}=\sum_k w_k f_k(x,i,a)$。在大语料上，仅字对模板与4个标签组合就可产生数以百万计的特征，绝大多数只出现几次，因此要用L2或L1正则约束权重。“生命”这个字对在训练语料中多以“生”为B出现，于是$\mathbb I[x_ix_{i+1}=\text{生命}\wedge y_i=\mathrm B]$获得正权重，这正是HMM无法直接使用的右邻证据。

表外的项在$\bm M$、$\bm m^{\mathrm{start}}$、$\bm m^{\mathrm{end}}$中取$-\infty$，整句得分、训练目标与解码按式~\eqref{eq:nlp-analysis-legal-mask}进行。与HMM相比，CRF在整句上归一化，不受第\ref{sec:nlp-classical-memm}小节所述MEMM标注偏置的影响；与逐位置分类相比，它在训练和解码中都只考虑合法序列。

解码用Viterbi。令$\delta_i(b)$为结束于位置$i$、标签为$b$的最高前缀分数，则$\delta_1(b)=\tilde\pi_b+\bm E_{1,b}$，$\delta_i(b)=\bm E_{i,b}+\max_a[\delta_{i-1}(a)+\tilde{\bm T}_{a,b}]$，其中$\tilde\pi=\bm\pi+\bm m^{\mathrm{start}}$、$\tilde{\bm T}=\bm T+\bm M$；同时记录取得最大值的前驱。末位选$\argmax_b[\delta_n(b)+\tilde\rho_b]$，再沿前驱回溯。两个末标签相同的前缀，未来新增的得分完全相同，因此只保留较大者不会丢掉全局最优；时间为$O(nL^2)$。

对“研究生命”，令所有合法转移和起止得分为0，使用表~\ref{tab:nlp-analysis-viterbi}的位置得分。逐位置选择最大值会得到S、E、B、E，分数为15，但S后不能接E，这一序列不在输出空间中。

\begin{table}[htbp]
  \centering
  \begin{tabular}{@{}c rrrr rrrr@{}}
    \toprule
    &\multicolumn{4}{c}{位置得分$\bm E_{i,\cdot}$}
    &\multicolumn{4}{c}{最优前缀$\delta_i$}\\
    字&B&M&E&S&B&M&E&S\\
    \midrule
    研&3&0&0&4&3&$-\infty$&$-\infty$&4\\
    究&0&1&4&0&4&4&7&4\\
    生&3&1&2&0&10&5&6&7\\
    命&0&0&4&1&7&10&14&8\\
    \bottomrule
  \end{tabular}
  \caption{BMES的教学得分与Viterbi图表。起止和转移采用正文掩码，合法项为0。右半表按末标签分别保留最优前缀；句末的M虽有10分，但句末掩码不允许以M结束。}
  \label{tab:nlp-analysis-viterbi}
\end{table}

例如$\delta_2(\mathrm E)=4+\delta_1(\mathrm B)=7$，而$\delta_3(\mathrm B)=3+\max\{\delta_2(\mathrm E),\delta_2(\mathrm S)\}=10$。句末只比较E的14和S的8，选E；前驱依次为位置3的B、位置2的E、位置1的B，得到BEBE，分数14。恢复程序见算法~\ref{alg:nlp-analysis-bmes-recover}：B开启跨度，M延长跨度，E闭合跨度，S独自成词，结果为“研究”$[0,2)$和“生命”$[2,4)$。

\Needspace{12\baselineskip}
\begin{algorithm}[BMES标签序列到原文字符跨度的恢复]
  \label{alg:nlp-analysis-bmes-recover}
  \begin{algorithmic}[1]
    \Require 合法BMES序列$y_{1:n}$，每个分析位置在原文中的起止偏移$[o_i,o_i')$
    \Ensure 词跨度列表$\mathcal W$
    \State $\mathcal W\gets[\,]$；$s\gets\text{未定义}$
    \For{$i=1,\ldots,n$}
      \If{$y_i=\mathrm S$}
        \State 把$[o_i,o_i')$加入$\mathcal W$
      \ElsIf{$y_i=\mathrm B$}
        \State $s\gets o_i$
      \ElsIf{$y_i=\mathrm E$}
        \State 把$[s,o_i')$加入$\mathcal W$；$s\gets\text{未定义}$
      \EndIf \Comment{M只延长当前跨度}
    \EndFor
    \State \Return $\mathcal W$ \Comment{序列合法时，每个B都有对应的E}
  \end{algorithmic}
\end{algorithm}

恢复步骤要求输入合法，这正是掩码在训练与解码中共同起作用的意义：解码器只产生合法序列，恢复程序无须猜测孤立的M或E属于哪个词。若原文在送入模型前做过全角半角转换、空白删除等规范化，偏移$o_i$必须指向规范化之前的原文，否则即使标签正确也会返回错误位置。用有限的大负数近似$-\infty$时，还要确认数值范围足以压过合法路径的分差。

特征CRF的效果取决于特征模板与语料规范的匹配。Tseng等在SIGHAN 2005评测中加入的未登录词特征，只在词较长的AS与MSR语料上带来提高；在HK语料上反而降低了F1，作者把原因追到测试集含有训练集中未见的标点\scite{v9a-tseng2005}。

\subsubsection{预训练编码器分词与词格模型}
\label{sec:nlp-analysis-pretrained-lattice}

特征CRF的模板要人工设计，窗口之外的信息也用不上。用预训练编码器替换特征层，是第\ref{sec:nlp-analysis-paradigm-tagging}小节三层结构的直接实例。中文BERT类模型的词表对常用汉字基本按单字切分，一句$n$个字的输入加上首尾的[CLS]与[SEP]，编码后得到$(n+2)\times d$的表示（基础规模$d=768$）；去掉两个特殊位置即为$\bm H\in\mathbb R^{n\times d}$。线性层$\bm W\in\mathbb R^{d\times4}$给出$\bm E\in\mathbb R^{n\times4}$，上面接逐位置softmax或带BMES掩码的CRF。CRF层只新增$4\times4$的转移矩阵和两个长度为4的起止向量，与编码器的上亿参数相比可以忽略，代价主要在编码器。整个网络端到端微调，流程见第\ref{sec:nlp-classical-bert-finetuning}小节。

它比特征CRF强在两处。其一，表示由全句自注意力得到，“生”的向量同时受“研究”和“命”影响，不受窗口宽度限制；其二，预训练让模型在大量无标注文本中见过低频字和新词的组合方式，标注语料中从未出现的人名、术语也能得到有用的表示，这直接针对未登录词这一主要错误来源。逐位置softmax与CRF的差距在强编码器上会缩小，因为上下文表示已经隐含了相邻标签的大部分信息；但softmax仍可能输出非法序列，保留带掩码的CRF或至少在解码时施加掩码，才能保证恢复程序收到合法输入。

字与模型词元并非总是一一对应。数字、拉丁字母串和部分符号会被合成一个词元，例如“14:00”可能被切成“14”“:”“00”三个词元。分词标签定义在字上，处理办法是在词元级别预测并规定合成规则：若一个词元覆盖多个字，这几个字被视为必然在同一词内，词元的标签由首字与末字的标签合成：首字为B或S表示词元开启一个词，末字为E或S表示词元闭合一个词，两者都成立记S，只开启记B，只闭合记E，都不成立记M。当参考词界落在一个词元内部时，这种表示无法输出正确答案，应在数据转换时统计并报告这类样本，而不是静默丢弃。以“会议14:00开始”为例，设词元为“会”“议”“14”“:”“00”“开”“始”，参考切分为“会议／14:00／开始”。字级标签中“1”为B、“4”为M，于是词元“14”合成为B；“:”为M；“0”为M、“0”为E，词元“00”合成为E。解码得到的词元级标签序列B、E、B、M、E、B、E仍满足BMES掩码，再按每个词元覆盖的原文偏移恢复出$[0,2)$、$[2,7)$、$[7,9)$。若某规范要求把“14”与“:00”切开，词元“00”内部没有冲突，这种切分仍可表示；若要求把“14”拆成“1／4”，就落在词元内部，无法表示。

预训练编码器仍以字为单位，词典中积累的词语知识只能通过微调数据间接学到。最简单的补救是\term{词典特征}（lexicon feature）：对每个字$x_i$，构造若干二值指示“以$x_i$开头（或结尾）的长度为$k$的片段是否在词典中”，$k=2,\ldots,5$，把这些指示拼到字向量上。它只告诉模型“这里有个词”，却不提供这个词本身的表示，也不能区分同一位置的多个候选词。

\term{词格}（lattice）模型更进一步，把词典匹配到的所有词作为额外单元接入字符序列：字符链上每个匹配词$w_{b,e}$（从第$b$个字到第$e$个字）构成一条从$b$跳到$e$的边，整体形成一个格。Lattice-LSTM是这一思想的代表，最初为中文实体识别提出\scite{nlp-analysis-zhang2018lattice}：它在普通字符LSTM之外，为每个匹配词计算一个词单元$\bm c^{\mathrm w}_{b,e}$，输入是词嵌入和第$b$个字的隐藏状态；在词尾字$e$处，字符单元把自己的候选更新$\tilde{\bm c}_e$与所有以$e$结尾的词单元按门控加权：
\[
  \bm c_e=\sum_{b}\bm\alpha_{b,e}\odot\bm c^{\mathrm w}_{b,e}+\bm\alpha_e\odot\tilde{\bm c}_e,
\]
其中各门向量$\bm\alpha_{b,e}$与$\bm\alpha_e$由各自的输入门经逐维指数归一化得到，使它们在每一维上相加为1。

图~\ref{fig:nlp-analysis-lattice}以“研究生命”和词典$\{$研究，研究生，生命$\}$为例画出这个格。三个匹配词分别是$w_{1,2}$、$w_{1,3}$与$w_{3,4}$：前两个词单元从第1个字的状态出发，第三个从第3个字的状态出发；它们分别在第2、3、4个字处并入字符链。于是“究”的单元中混入了“研究”的信息，“生”混入了“研究生”，“命”混入了“生命”。交集歧义在格中表现为两条重叠的词边：“研究生”与“生命”都覆盖“生”。门控根据语境决定各个字更依赖哪个候选词，例如在“研究生命的起源”中，模型可以学会在“生”处压低“研究生”的贡献，让后面的“生命”单元主导“命”处的表示，再由CRF输出B、E、B、E。

\begin{figure}[htbp]
  \centering
  % 第九卷 9.2.1.4 节：“研究生命”上的词格（教学词典：研究、研究生、生命）。
% 下方为字符LSTM链，上方弧为词典匹配词；词单元从起始字取状态，在词尾字处经门控并入字符单元。
\begin{tikzpicture}[
  x=1mm,y=1mm,figure picture,
  char/.style={rectangle,rounded corners=.8mm,draw=FigureInput,fill=FigureInputFill,
    line width=.6pt,minimum width=10mm,minimum height=7mm,inner sep=0pt,
    font=\figurefont\footnotesize},
  cell/.style={rectangle,rounded corners=.8mm,draw=FigureStroke,fill=white,
    line width=.6pt,minimum width=10mm,minimum height=6mm,inner sep=0pt,
    font=\figurefont\scriptsize},
  word/.style={rectangle,rounded corners=.8mm,draw=FigureModel,fill=white,
    line width=.6pt,minimum height=5.5mm,inner xsep=2pt,font=\figurefont\scriptsize},
  chain/.style={figure flow},
  wordarc/.style={draw=FigureModel,line width=.8pt},
  merge/.style={draw=BookTeal,line width=.9pt,-{Stealth[length=1.6mm,width=1.1mm]}}
]
  \foreach \i/\c/\x in {1/研/0,2/究/22,3/生/44,4/命/66}{
    \node[char] (x\i) at (\x,0) {\c};
    \node[cell] (h\i) at (\x,11) {$\bm c_{\i}$};
    \draw[figure link] (x\i) -- (h\i);
  }
  \draw[chain] (h1) -- (h2);
  \draw[chain] (h2) -- (h3);
  \draw[chain] (h3) -- (h4);
  % 研究：1→2
  \draw[wordarc] (h1.north) to[out=70,in=180] (8,22) -- (14,22);
  \node[word] (wa) at (11,22) {研究};
  \draw[merge] (wa.east) to[out=0,in=110] (h2.north);
  % 研究生：1→3
  \draw[wordarc] ([xshift=-2mm]h1.north) to[out=90,in=180] (14,33);
  \node[word] (wb) at (22,33) {研究生};
  \draw[merge] (wb.east) to[out=0,in=105] ([xshift=-1.5mm]h3.north);
  % 生命：3→4
  \draw[wordarc] ([xshift=1.5mm]h3.north) to[out=75,in=180] (52,24);
  \node[word] (wc) at (55,24) {生命};
  \draw[merge] (wc.east) to[out=0,in=110] (h4.north);
  \node[font=\figurefont\scriptsize,text=FigureMuted,anchor=west] at (74,11) {字符单元};
  \node[font=\figurefont\scriptsize,text=FigureMuted,anchor=west] at (74,0) {输入字符};
  \node[font=\figurefont\scriptsize,text=FigureMuted,anchor=west] at (74,28) {词典词单元};
\end{tikzpicture}

  \caption{“研究生命”上的词格。下方是字符链，每个字一个LSTM单元；上方的弧是词典匹配到的词，每个词从起始字取状态、计算一个词单元，在词尾字处经门控并入字符单元（实心箭头）。“研究生”与“生命”两条弧在“生”处重叠，对应交集歧义；采用哪一个由门控与上层CRF按语境决定。词典为教学设定。}
  \label{fig:nlp-analysis-lattice}
\end{figure}

Yang等把同一结构用于中文分词，并把格中的词典词换成从语料中自动抽取的子词单元，从而不再依赖人工词典\scite{nlp-analysis-yang2019lattice}。词单元的完整门控运算，以及把字与词展平为带头尾位置的序列、从而可以并行计算的FLAT，见第\ref{chap:nlp-information-extraction}章的中文实体识别。

这一谱系的进步在于把词典从硬约束变成软证据：最大匹配必须服从词典，词格模型只把词典词当作可被语境否决的候选。代价有两点。其一，效果依赖词典的覆盖与质量，领域换了往往要换词典；其二，Lattice-LSTM每个位置的输入数随匹配词数变化，难以批量并行，训练与推断明显慢于普通字符模型。

\subsubsection{多标准分词}
\label{sec:nlp-analysis-multi-criteria}

公开的中文分词语料遵循不同规范。以“姚明进入总决赛”为例，按中文树库的规范应切为“姚明／进入／总决赛”，按北大语料的规范则切为“姚／明／进入／总／决赛”，人名的姓与名、复合词的切分粒度都不同\scite{nlp-analysis-chen2017mccws}。把这些语料直接混合训练，同一字串会收到互相矛盾的标签；每个规范单独训练一个模型，又浪费了各语料共享的大部分边界知识（绝大多数词在各规范下切法一致）。\term{多标准分词}（multi-criteria word segmentation）用一个模型同时学习$K$种规范，并在预测时指定输出哪一种。

把这个例子转成BMES，冲突的位置一目了然：按前一种规范，七个字的标签为B、E、B、E、B、M、E；按后一种为S、S、B、E、S、B、E。第3、4、7个字的标签两边一致，第1、2、5、6个字不一致。若混合训练，同一输入在这些位置上有两种标签，模型最多只能在两者之间分配概率，学到的是两种规范的折中，在任何一个测试集上都会出错；若分开训练，“进入”这类两边一致的边界要在两个模型中各学一遍。理想的模型应当让一致的部分共享，只把不一致的部分交给规范信息决定。

早期做法采用共享—私有结构：一个共享编码器学习各规范共同的边界知识，每种规范另有私有编码器与CRF输出层；再加一个判别共享特征来自哪个语料的分类器，用对抗训练使共享特征不含规范信息\scite{nlp-analysis-chen2017mccws}。参数随规范数线性增长。后来的做法更简洁：所有语料共用一个Transformer编码器和一个输出层，只在输入首位加入一个表示规范的特殊词元$[\mathrm{c}_k]$，规范信息通过自注意力传到每个位置\scite{nlp-analysis-qiu2020mccws}。记第$k$种规范的训练集为$\mathcal D_k$，训练目标是各规范的CRF损失之和：
\begin{equation}
  \mathcal L=\sum_{k=1}^{K}\sum_{(x,y)\in\mathcal D_k}
  \mathcal L_{\mathrm{CRF}}\bigl(y;\ \operatorname{Enc}([\mathrm c_k];x)\bigr),
  \label{eq:nlp-analysis-multi-criteria}
\end{equation}
$\mathcal L_{\mathrm{CRF}}$为式~\eqref{eq:nlp-analysis-legal-mask}中的带掩码条件负对数似然。规范词元在不同语料中取不同值，模型因此能把“同一字串在规范$k$下怎么切”与“字串本身的边界证据”分开学习：在“进入”这类各规范一致的位置，规范词元的影响很小；在人名、复合词等规范分歧处，各位置通过注意力读取规范词元，输出相应的切法。新增一种规范时，只需新增一个规范嵌入，再用该规范的语料继续训练。与共享—私有结构相比，规范词元方法的参数几乎不随$K$增长，并且规范信息可以影响编码器的每一层，而不只是输出层；代价是规范差异完全靠一个嵌入向量传递，规范之间差别很大时，表达能力可能不如独立的输出层。

评价按规范分别进行：在规范$k$的测试集上，输入时也用$[\mathrm c_k]$。多标准模型只能在训练过的规范之间切换；各规范测试集上的F1即使平均成一个数，也只是汇总，不对应某种统一规范。

表~\ref{tab:nlp-analysis-segmentation-methods}按“解决了前一方法的什么问题”归纳本小节的方法。

\begin{table*}[htbp]
  \centering
  \begin{tabularx}{\linewidth}{@{}>{\raggedright\arraybackslash}p{24mm}>{\raggedright\arraybackslash}p{36mm}XX@{}}
    \toprule
    方法&解决的问题&核心机制&新代价或局限\\
    \midrule
    最大匹配／候选图&把词典知识直接用于切分&最长匹配，或在候选词图上求最短路径&不能用语境消歧，词典缺项即漏召回\\
    BMES--HMM&词典无法枚举新词&字符边界标注，计数估计转移与发射&发射只依赖当前字\\
    BMES--CRF&需要使用重叠的上下文特征&特征发射分＋转移分，带合法掩码全局归一化&特征模板需人工设计\\
    预训练编码器＋CRF&人工特征覆盖不足，未登录词召回低&上下文表示替代特征，端到端微调&编码代价高；字与词元对齐需处理\\
    词格模型&字符模型丢失词典信息&词典词作为候选单元，按语境门控融合&依赖词典质量；Lattice-LSTM难并行\\
    多标准分词&不同规范的语料无法合并&共享编码器＋规范标记或私有输出层&只能输出训练过的规范；须分规范评价\\
    \bottomrule
  \end{tabularx}
  \caption{中文分词方法的谱系。每一行针对上一行的主要局限；比较时须固定分词规范与评价语料，表中不跨语料排列性能。}
  \label{tab:nlp-analysis-segmentation-methods}
\end{table*}

\subsection{词形还原与形态属性}
\label{sec:nlp-analysis-morphology}

得到词界后，同一个词还可能以多种表面形式出现。英语studies与studied通常还原为study，而studies在句中可能是复数名词，也可能是动词第三人称单数。\term{基本形式}（lemma）是词典或标注规范选定的代表形式；\term{形态属性}（morphological features）记录数、时态、人称等语法变化。基本形式不同于词干提取的结果：把studies截成studi是词干提取可能产生的字符串，却不是一个词。中文的“昨天讨论”和“明天讨论”不靠动词屈折区分时间，这类任务主要针对屈折变化丰富的语言。三种输出的区别可以用同一个词看清：studies的词干提取结果可能是studi，基本形式是study，形态属性在名词读法下是“复数”，在动词读法下是“现在时、第三人称、单数”。词干提取只用于把同根词归并（例如检索时让study与studies命中同一索引项），不要求结果是合法词；基本形式要求是词典中的词；形态属性则要求与句法语境一致。

词形还原的挑战包括：屈折往往伴随拼写变化（y变ie、辅音双写）；不规则形式（went到go）无法由规则推出；同形词需要语境才能决定基本形式与属性；低资源语言缺少词典与规则；基本形式与属性彼此相关，需要联合输出。方法谱系依次为：词典加\term{有限状态转导器}（finite-state transducer, FST）精确、可解释，但覆盖受词典与规则限制；编辑树分类从数据中学习“怎样改词形”，能推广到同一变化模式的新词；字符级序列到序列模型直接生成基本形式，可处理较长或不规则的变换。评价分别报告基本形式精确匹配率、逐项属性准确率和整个属性组合的准确率。

\subsubsection{词典与有限状态形态分析}
\label{sec:nlp-analysis-fst}

只查表面词形会错过大量规则变化，把每种变化形式都写进词典又过于冗余。FST提供了一种紧凑写法：在一条状态路径上同时读入一个符号串、写出另一个符号串，从而把“基本形式＋属性”与表面形式联系起来。形式上，FST是一个五元组$T=(Q,\Sigma,q_0,F,\delta)$\scite{nlp-analysis-mohri1997}：$Q$是有限状态集；$\Sigma\subseteq(\Sigma_{\mathrm{in}}\cup\{\epsilon\})\times(\Sigma_{\mathrm{out}}\cup\{\epsilon\})$是由“输入符号:输出符号”对组成的字母表，$\epsilon$表示空串，允许一侧不读或不写；$q_0\in Q$是初始状态；$F\subseteq Q$是终止状态集；$\delta\subseteq Q\times\Sigma\times Q$是转移关系。从$q_0$到某个终止状态的一条路径，把路径上各符号对的输入侧连起来得到串$u$，输出侧连起来得到串$v$，就说$T$接受$(u,v)$。$T$接受的全部串对构成一个二元关系，交换输入与输出侧即得到逆关系，因此同一个转导器既能从分析生成表面形式，也能从表面形式反查分析。

形态分析器通常由两个转导器\term{复合}（composition）而成。若$T_1$把$u$映射到$v$，$T_2$把$v$映射到$w$，复合$T_1\circ T_2$就把$u$映射到$w$；构造时以状态对$(q_1,q_2)$为新状态，当$T_1$某转移的输出符号与$T_2$某转移的输入符号相同时，合成一条从$(q_1,q_2)$出发的转移，中间符号在结果中消失。词典转导器$T_{\mathrm{lex}}$把\code{study+N+Pl}映射为带语素边界的中间形式\code{study\^{}s}；拼写规则转导器$T_{\mathrm{rule}}$实现“辅音后的y在边界前写作ie，并删去边界符”，把\code{study\^{}s}写成\code{studies}。两者复合后得到从分析到表面形式的单一转导器，拼写规则可以先写成重写规则再编译成FST\scite{nlp-analysis-kaplan1994}。

用一个只处理本例的微型系统把复合过程走一遍。$T_{\mathrm{lex}}$中有一条路径依次读写\code{s:s}、\code{t:t}、\code{u:u}、\code{d:d}、\code{y:y}、\code{+N:}$\epsilon$、\code{+Pl:\^{}s}（最后一个符号对一次写出两个字符，可以拆成两步），它把\code{study+N+Pl}映射为\code{study\^{}s}。$T_{\mathrm{rule}}$有四个状态：在$q_0$上，除y和边界符外的字母原样复制并留在$q_0$；读到y时有两种选择，\code{y:y}留在$q_0$，或\code{y:ie}转到$q_1$；$q_1$只接受\code{\^{}:}$\epsilon$转到$q_2$；$q_2$只接受\code{s:s}转到终止状态$q_3$；另外$q_0$上的\code{\^{}:}$\epsilon$转到$q_2$，处理不以y结尾的词干。复合时，$T_{\mathrm{lex}}$写出的y与$T_{\mathrm{rule}}$读入的y相接，状态对从$(\cdot,q_0)$走到$(\cdot,q_1)$并写出ie；接着$T_{\mathrm{lex}}$写出的边界符被$T_{\mathrm{rule}}$删去，s原样保留。复合结果中这条路径的符号对为\code{s:s}、\code{t:t}、\code{u:u}、\code{d:d}、\code{y:ie}、\code{+N:}$\epsilon$、\code{+Pl:s}，中间形式\code{study\^{}s}不再出现。但若在y处走\code{y:y}分支，$q_0$读边界符后进入$q_2$，这个微型系统同样会接受错误的\code{studys}。可见它缺少语境条件；实际的拼写规则要写明“y前为辅音、后接边界与s时必须改写”，例如boys中元音后的y保持不变，这类必选条件正是重写规则编译成FST时要处理的部分。

沿逆方向查询\code{studies}，得到两条接受路径：
\[
\begin{aligned}
 \code{studies}&\leftrightarrow
   (\code{study},\mathrm{NOUN},\mathrm{Number=Plur}),\\
 \code{studies}&\leftrightarrow
   (\code{study},\mathrm{VERB},\mathrm{Tense=Pres},
                       \mathrm{Person=3},\mathrm{Number=Sing}).
\end{aligned}
\]
FST只说明两个分析都合乎词典和规则，选哪一个要看语境。在“She studies”中，为两个候选提取“左邻词为代词”“左邻为数词”等特征，设分类器给名词候选$-1$、动词候选2，softmax后动词概率约为0.953；“two studies”中得分若反为2与$-1$，则选复数名词。规则由语言资源提供，消歧分类器由语境与正确分析的配对训练，两种知识来源各自维护。另一种做法是把排序也放进转导器：给每条转移附一个权重（如负对数概率），路径权重为各转移权重之和，查询时按权重给分析排序，这就是加权有限状态转导器；复合与求最短路径在加权情形下仍有高效算法\scite{nlp-analysis-mohri1997}。但转移权重只能反映规则与词干本身的先验频率，读不到句子语境，“She studies”与“two studies”会得到相同的排序，所以同形歧义仍要交给语境分类器。

FST方法的优点是精确与可解释：错误可以追到缺失的词干（词典漏候选）或缺失的规则（拼写变化漏候选），也能保证输出的属性组合合法。缺点是覆盖完全取决于人工资源，新词、外来词和低资源语言都难以处理，这正是下面两种数据驱动方法要弥补的。

\subsubsection{编辑脚本与编辑树预测}
\label{sec:nlp-analysis-edit-lemma}

大量词语共享相同的变化，可以预测“怎样修改词形”，而不必给每个完整词形单独建类别。\term{编辑脚本}（edit script）是一段可执行的字符修改操作。直接按绝对位置记录的脚本难以推广，因为长度不同的词需要在不同位置修改；\term{编辑树}（edit tree）只记录要保留的公共部分在词首与词尾各留出多少字符，以及两端未对齐部分怎样替换，因此能用于不同长度的词\scite{nlp-analysis-chrupala2008}。

编辑树由词形$w$与基本形式$\ell$递归构造。若两者没有公共子串，就生成一个替换叶节点“把$w$替换为$\ell$”；否则找出最长公共子串，设它在$w$中前面有$a$个字符、后面有$b$个字符，在$\ell$中前面和后面分别是$\ell$的前缀和后缀，生成节点$\operatorname{Split}(a,b)$，再对两侧的前缀对和后缀对递归构造左右子树。应用时，节点$\operatorname{Split}(a,b)$要求输入长度不小于$a+b$，保留中间部分并把两端分别交给子树处理；替换叶节点要求输入恰好等于它记录的串。对studies与study，最长公共子串为stud，前侧为空，后侧为ies对y，得到$\operatorname{Split}(0,3)$，左子树为空串替换，右子树为“ies替换为y”。同理，ties与tie的最长公共子串为tie，得到$\operatorname{Split}(0,1)$，右子树为“s替换为空串”。算法~\ref{alg:nlp-analysis-edit-tree}给出构造与应用两个过程。

\Needspace{16\baselineskip}
\begin{algorithm}[编辑树的构造与应用]
  \label{alg:nlp-analysis-edit-tree}
  \begin{algorithmic}[1]
    \Require 训练时为词形$w$与基本形式$\ell$；应用时为编辑树$r$与新词形$u$
    \Ensure 编辑树$r$；或$u$的基本形式，前提不成立时返回失败
    \Function{Build}{$w,\ell$}
      \State 求$w$与$\ell$的最长公共子串，记它在$w$中为$w[a:|w|-b]$，在$\ell$中为$\ell[c:|\ell|-d]$
      \If{公共子串为空}
        \State \Return $\operatorname{Replace}(w\to\ell)$
      \EndIf
      \State 左子树$\gets\Call{Build}{w[0:a],\ell[0:c]}$
      \State 右子树$\gets\Call{Build}{w[|w|-b:],\ell[|\ell|-d:]}$
      \State \Return 以左、右子树为孩子的节点$\operatorname{Split}(a,b)$
    \EndFunction
    \Function{Apply}{$r,u$}
      \If{$r=\operatorname{Replace}(s\to t)$}
        \State \Return $t$，若$u=s$；否则返回失败
      \EndIf
      \State 设$r=\operatorname{Split}(a,b)$；若$|u|<a+b$，返回失败
      \State $p\gets\Call{Apply}{\text{左子树},u[0:a]}$；$q\gets\Call{Apply}{\text{右子树},u[|u|-b:]}$
      \State \Return $p+u[a:|u|-b]+q$；任一子调用失败则返回失败
    \EndFunction
  \end{algorithmic}
\end{algorithm}

编辑树中只存长度$a,b$和两端的替换串，不存中间的公共部分，所以Apply过程可以把studies的编辑树用到任何以ies结尾、长度不小于3的词上。前提检查就是Apply过程中的两处失败条件。

把编辑树作为类别，词形还原就成了上下文分类：用语境表示$\bm h$对编辑树$r$打分$g_r=\bm w_r^\top\bm h+b_r$，softmax只在前提成立的编辑树上归一化，以正确编辑树的交叉熵训练。对未见过的retries，同一编辑树保留retr并补y，恢复retry；对ties，前提同样成立，却得到错误的ty，正确的tie对应$\operatorname{Split}(0,1)$那棵“删去末尾s”的编辑树；反过来，后者用于studies又会得到studie。两棵树对两个词的前提都成立，只能靠特征区分，例如“候选结果是否在词表中”：study与tie在词表中，ty与studie不在。只在前提成立的编辑树上归一化，与第\ref{sec:nlp-analysis-paradigm-tagging}小节的合法掩码作用相同：它保证分类器选出的类别一定能执行。

Lemming把编辑树候选与形态标注放进一个联合对数线性模型\scite{nlp-analysis-muller2015lemming}：给定词形与语境，模型对“基本形式候选$\ell$（由所有可执行的编辑树生成）、形态标签$m$”联合打分，特征包括编辑树类别、候选是否出现在词表中、候选与形态标签的共现等，形态标签之间仍由序列CRF连接。写成打分形式，对一句中第$i$个词的候选基本形式$\ell_i$与形态标签$m_i$，模型的对数势为
\[
  \bm\theta_{\mathrm L}^\top\bm f_{\mathrm L}(\ell_i,w_i)+\bm\theta_{\mathrm J}^\top\bm f_{\mathrm J}(\ell_i,m_i)+\bm\theta_{\mathrm M}^\top\bm f_{\mathrm M}(m_i,w,i)+\bm\theta_{\mathrm T}^\top\bm f_{\mathrm T}(m_{i-1},m_i),
\]
四项分别评价“基本形式与词形是否相配”“基本形式与形态标签是否相配”“形态标签与语境是否相配”和“相邻形态标签是否相配”。整句的条件概率在所有候选组合上全局归一化，训练与解码仍是第\ref{sec:nlp-analysis-paradigm-tagging}小节的格点动态规划，只是每个位置的状态从单个标签扩展为“基本形式候选×形态标签”。联合建模让“这是动词第三人称单数”的证据直接支持“去掉s”一类编辑，反之亦然。UDPipe 2在SIGMORPHON 2019共享任务中也采用这一类可执行规则，并以词形字符表示和句子上下文共同预测规则类别，107个语料的平均基本形式准确率为95.78，任务基线为94.17；该系统同时加入了BERT表示与模型集成，差值不能全部归于编辑规则\scite{v9a-straka2019}。

编辑树方法的优点是类别数有限、预测快，并能把常见变化推广到新词；局限是只能选择训练中出现过的编辑树，遇到全新的变化模式或不规则形式（went到go）只能依赖替换叶节点，而这种叶节点几乎不能推广。

\subsubsection{字符级序列到序列还原}
\label{sec:nlp-analysis-seq2seq-lemma}

当所需变换不在编辑树集合中时，可以直接逐字符生成基本形式。编码器读取词形字符与左右若干字符的语境窗口，得到状态$\bm e_1,\ldots,\bm e_k$；解码器在第$t$步有状态$\bm s_t$，用注意力权重$\alpha_{t,i}=\operatorname{softmax}_i\bigl(\bm v^\top\tanh(\bm W_e\bm e_i+\bm W_s\bm s_t)\bigr)$汇聚出上下文向量$\bm c_t=\sum_i\alpha_{t,i}\bm e_i$，再经线性层与softmax给出字符表上的分布$p_{\mathrm{vocab}}(\cdot)$。Lematus就是这样一个带注意力的编码器—解码器，它不预先输入词性或编辑树，而让模型从语境窗口中学会选择基本形式\scite{v9a-bergmanis2018}。

具体地，输入是一串字符：左语境的若干字符、一个左边界符、词形的字符、一个右边界符、右语境的若干字符，词与词之间的空格也作为字符保留。设字符嵌入维度为$d_{\mathrm{ch}}$，双向循环编码器每个方向的隐藏维度为$d$，则输入长度为$k$时，编码器输出$\bm e_i\in\mathbb R^{2d}$，整体为$k\times2d$；解码器是单向循环网络，状态$\bm s_t\in\mathbb R^{d_s}$，每步输出字符表加结束符上的分布。语境窗口让编码器看到“She”或“two”这类决定词性的线索，边界符告诉注意力哪些字符属于待还原的词。

基本形式的大部分字符原样来自输入，只有词尾少数位置发生变化。让解码器从整个字符表中重新生成每个字符，既浪费模型容量，也容易在长词中途拼错。\term{复制机制}（copy mechanism）把“生成一个字符”与“复制输入中某个位置的字符”混合起来\scite{nlp-analysis-gu2016copynet,nlp-intro-see2017}：
\begin{equation}
  \begin{aligned}
    p_{\mathrm{gen}}&=\sigma\bigl(\bm w_c^\top\bm c_t+\bm w_s^\top\bm s_t+\bm w_y^\top\bm y_{t-1}+b\bigr),\\
    p(v)&=p_{\mathrm{gen}}\,p_{\mathrm{vocab}}(v)+\bigl(1-p_{\mathrm{gen}}\bigr)\sum_{i:\,x_i=v}\alpha_{t,i}.
  \end{aligned}
  \label{eq:nlp-analysis-copy}
\end{equation}
$p_{\mathrm{gen}}\in(0,1)$是第$t$步选择生成的概率，由上下文向量、解码器状态和上一步输出的嵌入$\bm y_{t-1}$决定；第二项把注意力权重当作复制分布，对输入中所有等于$v$的位置求和，因为同一字符可以出现在多个位置，它们都支持输出$v$。$v$不在输入中时第二项为0，只能靠生成；$v$不在字符表中时第一项为0，只能靠复制，所以输出空间是字符表与输入字符的并集。后文的复制式三元组抽取、话语改写和关键短语生成都沿用这一形式。

以retries为例，输入字符依次为r、e、t、r、i、e、s。第一步要输出r，设注意力在第1个位置为0.6、第4个位置为0.1、其余位置合计0.3，$p_{\mathrm{vocab}}(\mathrm r)=0.3$，$p_{\mathrm{gen}}=0.4$，则
\[
  p(\mathrm r)=0.4\times0.3+0.6\times(0.6+0.1)=0.54.
\]
两个r位置的注意力都计入了复制概率。第五步要输出输入中没有的y，复制项为0，模型必须提高$p_{\mathrm{gen}}$并把生成分布集中在y上：若$p_{\mathrm{gen}}=0.9$、$p_{\mathrm{vocab}}(\mathrm y)=0.8$，则$p(\mathrm y)=0.72$。这组教学数值说明了机制的分工：复制负责原样保留的部分，生成负责变化点，$p_{\mathrm{gen}}$在变化点附近升高。

训练对$(w,\ell^\star)$的目标是包括结束符在内的条件负对数似然
\[
  -\sum_{j=1}^{|\ell^\star|+1}\log p\bigl(\ell_j^\star\mid w,\text{语境},\ell^\star_{<j}\bigr),
\]
按教师强制计算；推断用束搜索，生成到长度上限仍无结束符时应标记截断。束搜索比较的是整条路径的条件概率：设正确路径依次输出r、e、t、r、y、结束符的概率为$.9,.9,.8,.8,.7,.95$，路径概率为$0.344736$；另一条把retries原样复制下来的路径概率为$0.12$（教学数值）。束宽为2时两者都保留到结束，最终选择前者。贪心解码在某一步若被复制分布带偏（例如在第五步复制了i），就无法回头。形态属性另设预测头，从编码器的汇聚表示出发，为时态、人称、数、词性各预测一个含“不适用”的类别，再检查属性组合是否合法；“She retries”应恢复$\mathrm{VERB,Pres,3,Sing}$，“two retries”则是$\mathrm{NOUN,Plur}$，相同的基本形式并不决定属性。

与编辑树相比，字符生成能组合训练中未见过的变化，也能处理不规则形式，并能直接利用语境：Lematus的实验显示，语境窗口对同形歧义词的帮助大于对未见词的帮助\scite{v9a-bergmanis2018}。代价是逐字符串行解码，并可能拼出不存在的词形；编辑树的输出至少总是某个已知变化的结果。

\subsection{词性标注与短语组块}
\label{sec:nlp-analysis-pos-chunk}

\term{词性标注}（part-of-speech tagging, POS tagging）判定一个词在当前语境中的语法类别：“研究问题”中的“研究”通常是动词，“这项研究”中则是名词。\term{短语组块}（phrase chunking）只标出不重叠的浅层短语范围，例如一个连续的名词短语，不递归恢复完整的嵌套结构。两者分别输出类别序列和带类别的区间，评价分别用词性准确率和组块F1，后者按区间与类别都正确才算命中。

兼类现象在中文里尤其普遍，“研究”“发展”“建设”这类词常同时可作动词和名词，形式上没有屈折标记，只能靠语境：前面有“这项”“他的”时多为名词，后面直接接宾语时多为动词。组块则要在词性之上决定短语范围，例如“校外来宾”是一个名词组块，“须报名”中的“须”不属于它。

这一小节的挑战是：兼类词需要语境决定类别；未知词没有发射统计；相邻标签之间有依赖（限定词后多为名词）；组块标签有BIO合法性；中文若先分词再标注，分词错误会传给词性，而词性证据本可以帮助分词。谱系依次为：HMM用转移建模标签依赖；CRF加入任意上下文特征；BiLSTM-CRF以双向上下文表示代替人工特征；预训练编码器进一步减少特征设计；字符级联合标注把边界与词性放进同一个标签，消除流水线的误差传播。若起点是原始文本而非给定词界，词性与组块都应在原文跨度上比较，不能假定错误切分后的第$i$个词仍对应参考的第$i$个词。

\subsubsection{HMM词性标注}
\label{sec:nlp-analysis-hmm-pos}

把第\ref{sec:nlp-analysis-hmm-boundary}小节的字符观测换成词、BMES状态换成词性，就得到HMM词性标注。由已标注语料统计$p(t_i\mid t_{i-1})$、$p(w_i\mid t_i)$和起止概率，Viterbi对完整词性路径求最大值。与BMES不同，词性转移没有普适的禁止表，低频组合不应自动视为不可能。未见过的词没有发射统计，最简单的处理是在训练时把低频词替换为若干未知词类别（如“含数字”“以-ed结尾”“首字母大写”），估计这些类别在各词性下的发射概率，测试时把未见词映射到相应类别。替换必须在训练时就做，若只在测试时替换，就会用到一个从未估计过的观测。

考虑给定分词的“研究／问题”，只比较动词—名词（V、N）和名词—名词（N、N）。设起始概率$p(\mathrm V)=.6$、$p(\mathrm N)=.4$，观测概率$p(\text{研究}\mid\mathrm V)=.3$、$p(\text{研究}\mid\mathrm N)=.4$，转移$p(\mathrm N\mid\mathrm V)=.8$、$p(\mathrm N\mid\mathrm N)=.2$，以及$p(\text{问题}\mid\mathrm N)=.5$（教学参数，两条路径的句末因子相同而省略）。两条路径的得分分别为$.6\cdot.3\cdot.8\cdot.5=.072$与$.4\cdot.4\cdot.2\cdot.5=.016$。虽然“研究”的名词发射概率更高，整句仍选V、N，恢复为“研究”$[0,2)$：V，“问题”$[2,4)$：N。

一阶转移只看前一个标签，而许多搭配需要前两个标签才能判断。TnT把转移扩展到二阶，并用插值平滑解决二阶计数稀疏的问题\scite{nlp-analysis-brants2000}：
\begin{equation}
  \begin{aligned}
    p(t_i\mid t_{i-2},t_{i-1})={}&\lambda_1\hat p(t_i)+\lambda_2\hat p(t_i\mid t_{i-1})\\
    &+\lambda_3\hat p(t_i\mid t_{i-2},t_{i-1}),\qquad\lambda_1+\lambda_2+\lambda_3=1,
  \end{aligned}
  \label{eq:nlp-analysis-tnt}
\end{equation}
$\hat p$为最大似然估计。三个权重对所有标签三元组共用，用\term{删除插值}（deleted interpolation）估计：遍历每个出现过的三元组$(t_1,t_2,t_3)$，假设把它的这一次出现从语料中删去，分别计算剩余计数下的三元、二元、一元估计$\frac{f(t_1t_2t_3)-1}{f(t_1t_2)-1}$、$\frac{f(t_2t_3)-1}{f(t_2)-1}$、$\frac{f(t_3)-1}{N-1}$（分母为0时记该项为0），哪一项最大，就把$f(t_1t_2t_3)$加到对应的$\lambda$上，最后归一化。“删去一次”模拟了这个三元组在新文本中出现时的情形，避免高阶估计因计入自身而被高估。举一个教学计数：某三元组出现5次，其前两个标签共现10次，后两个标签共现20次而中间标签出现100次，末标签出现500次，语料共$N=10000$个标签，则三项依次为$4/9\approx0.444$、$19/99\approx0.192$、$499/9999\approx0.050$，三元项最大，$\lambda_3$加5。若某三元组只出现1次，删去后三元估计为0，它的计数就会记到二元或一元项上；稀疏的三元组越多，$\lambda_3$越小，模型越依赖低阶估计。未知词按词尾字母估计词性分布。设未知词的字母为$l_1,\ldots,l_n$，从空后缀开始逐步加长后缀，每加一个字母就把新的计数估计与上一级的平滑估计混合：
\[
  p(t\mid l_{n-i+1:n})=\frac{\hat p(t\mid l_{n-i+1:n})+\theta\,p(t\mid l_{n-i+2:n})}{1+\theta},
\]
其中$l_{a:n}$表示从第$a$个字母到词尾的后缀，$i$从1增加到$m$，递推从$p(t)$开始，$\hat p$只在训练集的低频词上统计（低频词与未知词的分布更接近），$m$取不超过10的最长已见后缀，$\theta$取各词性无条件概率的标准差。直观上，-ing、-tion这样的长后缀一旦在训练中出现，就几乎决定词性，短后缀只提供弱证据。得到的$p(t\mid\text{后缀})$再用贝叶斯公式除以$p(t)$转为与发射概率同尺度的量；首字母大写另作条件，分开统计。二阶模型的Viterbi状态是相邻两个标签，代价为$O(nK^3)$，$K$为词性数。TnT在宾州树库上的标注准确率约为96.7\%，与当时的最大熵标注器相当，而训练只需计数\scite{nlp-analysis-brants2000}。

HMM标注的优点是训练只需计数、速度快、参数可直接检查；局限仍是生成式假设，发射只看当前词，“后缀是-ing且左邻是is”这类重叠特征无法直接使用。

\subsubsection{CRF与BiLSTM-CRF局部标注}
\label{sec:nlp-analysis-crf-chunk}

特征CRF以左右邻词、词形、后缀和大小写等特征产生发射分数，配合转移矩阵与掩码完成标注，其目标与解码与第\ref{sec:nlp-analysis-bmes-crf}小节相同，只是标签换成词性或BIO组块标签。浅层名词组块采用BIO标签时，句首禁止I-NP，I-NP只能接在B-NP或I-NP之后；多类组块还要求延续标签的类型与开头相同。以“校外／来宾／须／报名”为例，设三种标签的发射分数为
\[
\begin{array}{c|rrr}
 &\mathrm{B\!-\!NP}&\mathrm{I\!-\!NP}&\mathrm O\\ \hline
 \text{校外}&3&4&0\\
 \text{来宾}&0&4&1\\
 \text{须}&0&0&3\\
 \text{报名}&0&0&3
\end{array}
\]
合法转移为0时，逐位置取最大得到I-NP、I-NP、O、O，首标签非法；带掩码的Viterbi恢复B-NP、I-NP、O、O，总分13。读到第一个O时闭合名词组块，按词的起止偏移得到“校外来宾”$[0,4)$。

BiLSTM-CRF把手工特征换成神经编码\scite{nlp-analysis-huang2015,nlp-analysis-lample2016,v9a-ma2016}。设句长$n$、批大小$B$，逐层运算为：
\begin{enumerate}
  \item 词嵌入：每个词查表得到$\bm e^{\mathrm w}_i\in\mathbb R^{d_{\mathrm w}}$，常用预训练词向量初始化，张量形状$B\times n\times d_{\mathrm w}$。
  \item 字符卷积：把词内字符嵌入排成$B\times n\times\ell_{\max}\times d_{\mathrm{ch}}$（$\ell_{\max}$为最大词长，不足补齐），沿字符维做宽度为3的一维卷积，再在字符维取最大值，得到$\bm e^{\mathrm c}_i\in\mathbb R^{d_{\mathrm c}}$；它捕捉后缀、前缀和大小写等形态信息，对未知词尤其重要。
  \item 拼接：$\bm x_i=[\bm e^{\mathrm w}_i;\bm e^{\mathrm c}_i]$，形状$B\times n\times(d_{\mathrm w}+d_{\mathrm c})$，训练时对其施加dropout。
  \item 双向LSTM：前向与后向各有$d$维隐藏状态，拼接得$\bm h_i=[\overrightarrow{\bm h}_i;\overleftarrow{\bm h}_i]\in\mathbb R^{2d}$，即$\bm H$的形状为$B\times n\times2d$。
  \item 发射层：$\bm E=\bm H\bm W+\bm b$，$\bm W\in\mathbb R^{2d\times L}$，得到$B\times n\times L$的发射分数。
  \item CRF层：转移矩阵$\bm T\in\mathbb R^{L\times L}$加合法掩码，按式~\eqref{eq:nlp-analysis-legal-mask}训练与解码；批内不同句长用长度掩码处理，填充位置不参与求和。
\end{enumerate}
字符卷积学习词形，BiLSTM学习句内用法，CRF学习相邻类别的约束，三者由同一个参考标签端到端训练。Ma与Hovy在英语宾州树库上控制其他设置，词性准确率由仅用BiLSTM的96.93依次升到加字符卷积后的97.33和再加CRF后的97.55，说明字符表示对词性的帮助大于CRF层\scite{v9a-ma2016}。字符表示也可以用字符级双向LSTM代替卷积，取前后两个方向的末状态拼接作为词的字符向量\scite{nlp-analysis-lample2016}；两者都能从后缀与大小写推断未知词的类别，卷积计算更快，双向LSTM对整个词形建模。CRF层的作用随任务而变：词性标签之间的约束较弱，相邻依赖大多已被双向编码捕捉，CRF的增益较小；实体与组块的BIO标签有硬约束（I-$X$必须接在同类之后）和较长的区间，CRF与合法掩码的作用更明显。

改用预训练编码器时，上述第1至4步由编码器替代，输入为子词序列。按第\ref{sec:nlp-analysis-paradigm-tagging}小节的首子词对齐，先从$B\times n_{\mathrm{sub}}\times d$的子词表示中按每个词首子词的下标收集出$B\times n\times d$的词表示，再接发射层与CRF。英语studies若被切成stud和\#\#ies，只有stud的位置产生标签；中文字符模型则通常不需要这一步。

\subsubsection{字符级分词--词性联合标注}
\label{sec:nlp-analysis-joint-pos}

流水线一旦把“研究生”固定成一个词，后面的词性标注器就无法恢复“研究／生命”。联合标注把边界与词性组合成一个标签，例如B-V、E-V、B-N、E-N，让词性证据参与词界选择\scite{v9a-ng2004}。训练时把每个已标注词展开为字符标签：“研究/V，生命/N”产生B-V、E-V、B-N、E-N，“研究生/N，命/N”产生B-N、M-N、E-N、S-N。

合法转移掩码在这里要同时表达两种约束：B-$t$和M-$t$只能转到M-$t$或E-$t$，保证同一词内词性一致；E-$t$、S-$t$之后才能转到任意词性的B或S。设两条候选的完整得分为8和6（教学设定），其余路径更低，约束解码选择前者；恢复程序在E-V处闭合“研究”$[0,2)$并附V，在E-N处闭合“生命”$[2,4)$并附N。B-V后接E-N这类词内词性不一致的组合，即使局部分数高也被掩码排除。

若有$k$种词性，联合标签数为$L=4k$，转移矩阵按$L^2=16k^2$增长，CRF的$O(nL^2)$代价随之放大；不过合法转移是稀疏的，每个B-$t$只有两个合法后继，按稀疏结构实现递推可以省去大部分无用比较。以中文树库的33个词性为例，$L=132$，稠密转移矩阵有$132^2=17424$项；合法项中B-$t$与M-$t$各有2个后继，E-$t$与S-$t$各有$2k=66$个后继，合计$33\times(2+2+66+66)=4488$项，约占四分之一。联合训练需要同时有词界和词性标注。它消除了分词到词性的误差传播，但组块仍须按自己的边界规范另行预测。

字符级标签并不是唯一的联合方式。另一条路线以词为单位、用转移系统逐步构造“词＋词性”序列：每一步决定当前字是并入上一个词，还是开启一个新词并给出它的词性，词级特征（整个候选词、词与词性的搭配）因此可以直接使用；由于没有可做动态规划的图表，推断用束搜索，训练用结构化感知机\scite{nlp-analysis-zhang2008joint}。两条路线分属第\ref{sec:nlp-analysis-paradigm-tagging}小节与第\ref{sec:nlp-analysis-paradigm-transition}小节的范式：字符标注可以精确解码，但特征只能落在字上；词级转移能用整词特征，代价是只有近似搜索。

表~\ref{tab:nlp-analysis-morph-pos-methods}归纳形态还原与词性、组块标注的方法。

\begin{table*}[htbp]
  \centering
  \begin{tabularx}{\linewidth}{@{}>{\raggedright\arraybackslash}p{27mm}>{\raggedright\arraybackslash}p{34mm}XX@{}}
    \toprule
    方法&解决的问题&核心机制&新代价或局限\\
    \midrule
    词典＋FST&从表面形式恢复基本形式与属性&词典转导器与拼写规则转导器复合，逆向查询；语境分类消歧&覆盖依赖人工资源，新词难处理\\
    编辑树分类&资源不足时从数据学变化&词形—基本形式的编辑树作为类别，只在前提成立的类别上归一化&只能选训练中出现过的变化\\
    字符级生成＋复制&不规则与未见变化&注意力编码器—解码器，复制与生成混合&串行解码，可能拼出不存在的形式\\
    HMM／TnT&标签依赖&一阶或二阶转移，删除插值与后缀模型&发射只看当前词\\
    特征CRF&需要重叠的上下文特征&特征发射＋转移＋合法掩码&特征需人工设计\\
    BiLSTM-CRF／预训练编码器&人工特征覆盖不足&字符卷积与双向编码（或预训练编码，首子词对齐）给发射分&编码代价高\\
    字符级联合标注&分词错误传给词性&边界×词性的组合标签，掩码保证词内一致&标签数与转移规模增大\\
    \bottomrule
  \end{tabularx}
  \caption{形态还原与词性、组块标注的方法谱系。前三行处理基本形式与属性，后四行处理逐词类别；比较时须固定标注规范与词界条件。}
  \label{tab:nlp-analysis-morph-pos-methods}
\end{table*}

\section{句法分析}
\label{sec:nlp-analysis-syntax}

词法分析给出了词和局部标签，却没有说明短语怎样嵌套，也没有说明哪个词支配哪个词。句法分析补上这两类信息。\term{成分结构}（constituency structure）用树表示相邻词怎样组成更大的短语，树的内部节点是短语类别；\term{依存结构}（dependency structure）用有向边表示中心词与依存词之间的关系，节点就是词本身。

图~\ref{fig:nlp-analysis-syntax}固定“小王／阅读／通知”的分词。成分树把“阅读通知”组成动词短语，再与“小王”组成句子；依存树以“阅读”为中心，分别连向主语“小王”和宾语“通知”。两种表示回答的问题不同：成分树说明哪些词连成一个单位，依存树说明词与词之间的支配关系。借助中心词规则可以把成分树转换成依存树，但反方向缺少短语类别和层次信息，两者一般不能无损互换。

\begin{figure*}[htbp]
  \centering
  \begin{tikzpicture}[
  x=1mm,y=1mm,
  every node/.append style={font=\figurefont\footnotesize,align=center},
  category/.style={draw=FigureModel,fill=FigureModel!10,
    rounded corners=1mm,inner sep=1.5mm},
  word/.style={draw=FigureInput,fill=FigureInput!10,
    minimum width=16mm,minimum height=8mm,inner sep=1mm},
  branch/.style={draw=NNLine,line width=.85pt},
  dep/.style={->,draw=NNWire,line width=1.1pt}
]
  \node[font=\figurefont\footnotesize\bfseries] at (40,59) {(a) 成分：短语包含关系};
  \node[category] (s) at (39,46) {S};
  \node[category] (np1) at (13,30) {NP};
  \node[category] (vp) at (53,30) {VP};
  \node[category] (v) at (41,15) {V};
  \node[category] (np2) at (68,15) {NP};
  \node[word] (a) at (13,0) {小王};
  \node[word] (b) at (41,0) {阅读};
  \node[word] (c) at (68,0) {通知};
  \foreach \parent/\child in {s/np1,s/vp,np1/a,vp/v,vp/np2,v/b,np2/c}
    \draw[branch] (\parent.south) -- (\child.north);
  \node at (13,-10) {$[0,2)$};
  \node at (41,-10) {$[2,4)$};
  \node at (68,-10) {$[4,6)$};

  \node[font=\figurefont\footnotesize\bfseries] at (127,59) {(b) 依存：词间有向关系};
  \node[category] (root) at (127,43) {0：人工根};
  \node[word] (d) at (102,0) {小王};
  \node[word] (e) at (127,0) {阅读};
  \node[word] (f) at (152,0) {通知};
  \draw[dep] (root.south) -- node[right,fill=white,inner sep=1mm] {root} (e.north);
  \draw[dep] (e.north west) to[out=145,in=45]
    node[above left,inner sep=1mm] {nsubj} (d.north);
  \draw[dep] (e.north east) to[out=35,in=135]
    node[above right,inner sep=1mm] {obj} (f.north);
  \node at (102,-10) {$1:[0,2)$};
  \node at (127,-10) {$2:[2,4)$};
  \node at (152,-10) {$3:[4,6)$};
  \draw[NNLine,dashed,line width=.5pt] (84,-14) -- (84,52);
\end{tikzpicture}

  \caption{同一句话的成分树与基本依存树。左图的NP、VP、S分别表示名词短语、动词短语和句子，方括号是字符跨度；右图的边由中心词指向依存词，0是人工根，nsubj和obj分别标记主语与宾语关系。两图固定同一分词与一种教学分析，不表示处理的先后。}
  \label{fig:nlp-analysis-syntax}
\end{figure*}

本节按输出结构分为两部分，每部分内部按第\ref{sec:nlp-analysis-paradigms}节的预测范式组织方法。成分解析依次覆盖区间动态规划（PCFG与CKY、Earley、神经跨度解析）、转移系统（移进--归约）和约束生成（序列化解析）；依存解析覆盖转移系统（arc-standard与arc-eager）和图上的精确解码（Eisner、Chu--Liu/Edmonds），再由双仿射网络提供弧分数，最后延伸到允许多入边的增强依存图。

\subsection{成分结构与解析}
\label{sec:nlp-analysis-constituency}

成分记为$(i,j,\ell)$，表示覆盖第$i$个词边界到第$j$个词边界之间的短语，类别为$\ell$。区间按词边界编号，不是字符偏移；输出时再经词跨度表还原原文。一棵树中的成分只能嵌套或互不相交，不能同时含有$[0,2)$和$[1,3)$这样交叉的区间。

“新／学生／代表”有两种分组：“新学生”的代表，即$[[\text{新 学生}]\ \text{代表}]$；或新任的“学生代表”，即$[\text{新}\ [\text{学生 代表}]]$。两种分组的词序完全相同，差别只在内部跨度。一个评分器可以给$[0,2)$与$[1,3)$都打高分，但合法树最多保留其中一个。这个例子集中体现了成分解析的四项挑战：一是\term{结构歧义}（structural ambiguity），同一词串对应多棵树，介词短语附着是英语中最常见的情形；二是上下文无关假设过强，规则概率只看当前类别，无法利用父节点和词汇信息；三是解析代价，在全部区间上搜索需要$O(n^3)$时间；四是输出必须是一棵合法树，局部最优的跨度拼不出合法结构。

方法谱系沿这几项挑战展开。PCFG与CKY提供可估计的概率树和精确搜索，父节点标注与词汇化缓解独立假设；Earley去掉乔姆斯基范式的要求，直接处理一般文法；神经跨度解析去掉显式文法，用整句上下文给跨度评分，仍以CKY保证合法；移进--归约把解析降到线性时间，代价是失去全局最优；序列化生成用统一的序列到序列接口输出括号树，合法性改由后处理或约束解码负责。

各方法用同一套指标比较。PARSEVAL评价把树转成带标签跨度的集合，计算预测集合与参考集合的精确率、召回率和F1，并事先固定是否计入根节点、词性节点和标点。以“新学生代表”为例，若参考为右分组、预测为左分组，计入根时参考集合为$\{\mathrm{NP}[0,3),\mathrm{NP}[1,3)\}$，预测集合为$\{\mathrm{NP}[0,3),\mathrm{NP}[0,2)\}$，只有根匹配，精确率、召回率和F1都是$1/2$；不计根时两者没有共同成分，F1为0。同一棵树在两种口径下的分数差别很大，比较不同系统时必须固定口径。

\subsubsection{PCFG与CKY图表解析}
\label{sec:nlp-analysis-pcfg-cky}

第\ref{sec:nlp-classical-pcfg-definition}小节已经定义了\term{概率上下文无关文法}（probabilistic context-free grammar, PCFG）：同一左部非终结符的展开概率之和为1，一棵树的概率是所用规则概率之积，见式~\eqref{eq:nlp-classical-pcfg}；第\ref{sec:nlp-classical-cky}小节给出了对所有树求和的inside算法，第\ref{sec:nlp-classical-inside-outside}小节给出了无树库时的inside--outside估计。本小节只处理解析：给定文法，求概率最大的树，并说明怎样从图表恢复原文中的成分。

有树库时，规则概率用相对频率估计：
\[
 p(A\to\beta)=\frac{N(A\to\beta)}{\sum_{\gamma}N(A\to\gamma)},
\]
$N(A\to\beta)$是规则在参考树中出现的次数；需要平滑时保持同一左部归一化。采用只有$A\to BC$与$A\to w$两类规则的乔姆斯基范式，本例固定根类别为NP，教学文法为
\[
\begin{array}{llll}
 \mathrm{NP}\to\mathrm A\ \mathrm{NP}&0.4&
 \mathrm{NP}\to\mathrm{NP}\ \mathrm N&0.3\\
 \mathrm{NP}\to\mathrm N\ \mathrm N&0.2&
 \mathrm{NP}\to\text{学生}&0.1\\
 \mathrm A\to\text{新}&1&
 \mathrm N\to\text{学生}&0.6\\
 &&\mathrm N\to\text{代表}&0.4
\end{array}
\]
NP规则概率之和为1，N规则概率之和也为1。这是只用于手算的微型文法；没有对应规则的组合不是低概率树，而是根本不在候选集合中。

\term{CKY图表解析}（CKY chart parsing）把每个“区间加根类别”作为可复用的状态。令$V[A,i,j]$为类别$A$覆盖$[i,j)$时最优子树的概率。长度为1的区间直接查词汇规则，$V[A,i,i+1]=p(A\to w_{i+1})$，没有规则时为0。长度至少为2时，根若使用$A\to BC$，左右子树必在某个边界$k$处分开，于是
\begin{equation}
 V[A,i,j]=
 \max_{\substack{A\to BC\\i<k<j}}
 p(A\to BC)\,V[B,i,k]\,V[C,k,j].
 \label{eq:nlp-analysis-pcfg-cky}
\end{equation}
这就是inside递推式~\eqref{eq:nlp-classical-inside}把求和换成取最大的形式。按区间长度从小到大填表，每个组合所需的两个子区间都已算好；同时保存取得最大值的规则与分割点$k$，供回溯使用。最优子结构成立的理由是：若最优树的某棵子树不是同一类别、同一区间下的最优子树，把它换成最优子树会提高整棵树的概率，与最优性矛盾。

表~\ref{tab:nlp-analysis-cky}列出所有非零条目。长度为2的“学生代表”有两条推导，必须先比较它们，而不是保留第一次找到的规则。

\begin{table}[htbp]
  \centering
  \begin{tabularx}{\linewidth}{@{}l X l@{}}
    \toprule
    区间&非零类别及计算&最优回溯\\
    \midrule
    $[0,1)$&$\mathrm A:1$&新\\
    $[1,2)$&$\mathrm N:0.6,\ \mathrm{NP}:0.1$&学生\\
    $[2,3)$&$\mathrm N:0.4$&代表\\
    $[0,2)$&$\mathrm{NP}:0.4\cdot1\cdot0.1=0.04$&$\mathrm A,\mathrm{NP},k=1$\\
    $[1,3)$&$\mathrm{NP}:\max(0.2\cdot0.6\cdot0.4,\,
      0.3\cdot0.1\cdot0.4)=0.048$&$\mathrm N,\mathrm N,k=2$\\
    $[0,3)$&$\mathrm{NP}:\max(0.4\cdot1\cdot0.048,\,
      0.3\cdot0.04\cdot0.4)=0.0192$&$\mathrm A,\mathrm{NP},k=1$\\
    \bottomrule
  \end{tabularx}
  \caption{“新／学生／代表”的完整非零CKY图表。未列出的条目为0。根的另一种顶层分组概率为0.0048；右侧区间$[1,3)$内部另有概率0.012的推导$\mathrm{NP}\to\mathrm{NP}\ \mathrm N$，但最优值为0.048。}
  \label{tab:nlp-analysis-cky}
\end{table}

从根$V[\mathrm{NP},0,3]$读出$(\mathrm A,\mathrm{NP},k=1)$，再读出右子树的$(\mathrm N,\mathrm N,k=2)$，得到
\[
 [\mathrm{NP}\ [\mathrm A\ \text{新}]
       [\mathrm{NP}\ [\mathrm N\ \text{学生}]
                    [\mathrm N\ \text{代表}]]],
\]
概率为0.0192，左分组为$0.3\cdot0.04\cdot0.4=0.0048$。原文“新学生代表”中三个词的字符跨度依次为$[0,1)$、$[1,3)$、$[3,5)$，所以右侧NP还原为“学生代表”$[1,5)$。这个结论只反映本例文法的偏好；若前文是“刚入学”，读者可能偏向左分组，而PCFG读不到这条背景。

实际文法并不都是二叉的。规则$A\to BCD$可以引入专用辅助类别$A_\star$，改写成$A\to B\,A_\star$与$A_\star\to CD$，原概率放在第一条，第二条设为1，回溯后删去辅助节点；辅助类别必须记住原规则身份，否则会引入原文法不存在的组合。单分支链$A\to B\to C$可以折叠成复合类别，恢复时再展开；空产生式需要事先消除。若根值为0，返回解析失败。长句用对数概率把乘积变成和，避免下溢。设二叉规则集合为$R_2$、非终结符集合为$N$，直接实现的时间为$O(n^3|R_2|)$，图表与回溯指针的空间为$O(n^2|N|)$。

微型文法的另一个问题是NP在不同位置共用一个分布。第\ref{sec:nlp-classical-annotation}小节介绍的\term{父节点标注}（parent annotation）把类别改写成“类别加父类别”，让展开概率依赖结构位置。以下计数是教学设定：设树库中100个NP节点的展开与上面的概率一致，其中60个是句根NP，40个位于另一个NP之下。按父类别拆开后，
\[
\begin{array}{lcccc}
 &\mathrm A\ \mathrm{NP}&\mathrm{NP}\ \mathrm N&\mathrm N\ \mathrm N&\text{学生}\\
 \mathrm{NP}^{\wedge}\mathrm{ROOT}\ (60)&18&30&12&0\\
 \mathrm{NP}^{\wedge}\mathrm{NP}\ (40)&22&0&8&10
\end{array}
\]
各列相加仍是$40,30,20,10$。重新估计后，右分组的概率为$p(\mathrm{NP}^{\wedge}\mathrm{ROOT}\to\mathrm A\ \mathrm{NP}^{\wedge}\mathrm{NP})\cdot p(\mathrm{NP}^{\wedge}\mathrm{NP}\to\mathrm N\ \mathrm N)\cdot0.6\cdot0.4=0.3\cdot0.2\cdot0.24=0.0144$，左分组为$0.5\cdot(0.55\cdot0.25)\cdot0.4=0.0275$，偏好翻转。原文法把两种位置的统计平均在一起，掩盖了“句根NP常以中心名词结尾、内嵌NP常是单个名词”这类差异；标注后类别数增加，每条规则的训练样本变少，需要平滑。Klein与Manning沿这一思路系统地拆分类别，不用词汇化，在英语宾州树库第23节40词以内的句子上得到86.36的带标签F1，接近同期的词汇化解析器，说明“按什么状态共享展开概率”本身就决定了解析能力\scite{v9a-klein2003}。

\subsubsection{Earley文法解析}
\label{sec:nlp-analysis-earley}

CKY要求文法是乔姆斯基范式。文法作者未必愿意把每条规则改成二叉形式，二叉化还会改变规则的含义和概率的放置位置。\term{Earley解析}（Earley parsing）不改写文法，用带点的规则记录“一条规则已经读到了哪里”\scite{v9a-earley1970}。

\term{带点项目}（dotted item）$[A\to\alpha\cdot\beta,i]$存放在图表位置$j$上，表示这条规则从边界$i$开始，已经匹配了$w_{i+1:j}$对应的$\alpha$，还在等待$\beta$。点是进度标记，不是输入词元。图表位置$j$的项目集合记作$C_j$。为了让句根也用同一机制处理，增加一条起始规则$S'\to\mathrm{NP}$。

\begin{definition}[Earley的三条推理规则]
  \label{def:nlp-analysis-earley-rules}
  初始时$[S'\to\cdot\,\mathrm{NP},0]\in C_0$。设输入为$w_1\cdots w_n$，按以下规则向图表加入项目，同一位置上相同的项目只保留一份：
  \[
  \begin{array}{lll}
   \text{预测}&
   \dfrac{[A\to\alpha\cdot B\beta,\,i]\in C_j\qquad B\to\gamma\in G}
         {[B\to\cdot\,\gamma,\,j]\in C_j}\\[3.2ex]
   \text{扫描}&
   \dfrac{[A\to\alpha\cdot a\beta,\,i]\in C_j\qquad w_{j+1}=a}
         {[A\to\alpha a\cdot\beta,\,i]\in C_{j+1}}\\[3.2ex]
   \text{完成}&
   \dfrac{[B\to\gamma\,\cdot,\,k]\in C_j\qquad[A\to\alpha\cdot B\beta,\,i]\in C_k}
         {[A\to\alpha B\cdot\beta,\,i]\in C_j}
  \end{array}
  \]
  横线上方是前提，下方是结论。若$[S'\to\mathrm{NP}\,\cdot,0]\in C_n$，则输入被文法接受。
\end{definition}

预测说明“等待$B$的项目需要从当前位置开始识别一个$B$”；扫描让点越过与下一个输入词相同的终结符，并把项目移到下一个位置；完成把一个已经识别完的$B$交还给所有在其起点等待$B$的项目。每个位置反复应用预测与完成直到不再产生新项目，再对下一个词做扫描。

沿用上一小节的文法，位置0预测出NP的四条规则$[\mathrm{NP}\to\cdot\,\mathrm A\ \mathrm{NP},0]$、$[\mathrm{NP}\to\cdot\,\mathrm{NP}\ \mathrm N,0]$、$[\mathrm{NP}\to\cdot\,\mathrm N\ \mathrm N,0]$、$[\mathrm{NP}\to\cdot\,\text{学生},0]$，以及A、N的词汇规则。规则$\mathrm{NP}\to\mathrm{NP}\ \mathrm N$是左递归的：预测它又会预测NP，但这些项目在$C_0$中已经存在，去重使预测停止。这正是Earley能直接处理左递归、而朴素的自顶向下递归下降会无限循环的原因。

扫描“新”后，$C_1$得到$[\mathrm A\to\text{新}\,\cdot,0]$；完成操作回到位置0，推进出$[\mathrm{NP}\to\mathrm A\cdot\mathrm{NP},0]$，后者又预测从1开始的NP。扫描“学生”后，$C_2$同时有$[\mathrm N\to\text{学生}\,\cdot,1]$与$[\mathrm{NP}\to\text{学生}\,\cdot,1]$：前者推进$[\mathrm{NP}\to\mathrm N\cdot\mathrm N,1]$；后者完成$[\mathrm{NP}\to\mathrm A\ \mathrm{NP}\,\cdot,0]$，即“新学生”，并推进$[\mathrm{NP}\to\mathrm{NP}\cdot\mathrm N,1]$。由于“新学生”是一个完整的NP，它又推进$[\mathrm{NP}\to\mathrm{NP}\cdot\mathrm N,0]$。扫描“代表”后，$C_3$完成三条NP：$[\mathrm{NP}\to\mathrm N\ \mathrm N\,\cdot,1]$与$[\mathrm{NP}\to\mathrm{NP}\ \mathrm N\,\cdot,1]$都覆盖“学生代表”，$[\mathrm{NP}\to\mathrm{NP}\ \mathrm N\,\cdot,0]$覆盖整句的左分组；前两者再完成$[\mathrm{NP}\to\mathrm A\ \mathrm{NP}\,\cdot,0]$，即整句的右分组。最终$[S'\to\mathrm{NP}\,\cdot,0]\in C_3$，输入被接受。

接受项目只说明至少存在一棵树。要输出树，每次扫描和完成都要记录结论由哪些前提项目得到；同一项目可以有多组来源，这些来源指针构成\term{共享解析森林}（shared packed forest）。本例中$[\mathrm{NP},0,3]$有两组来源，$[\mathrm{NP},1,3]$也有两组。给项目附上Viterbi内部概率：预测时设为规则概率，完成时乘上被交还子项目的概率，同一项目的多组来源取最大，就得到与CKY相同的结果：右分组0.0192胜过左分组0.0048。Stolcke把这种做法推广为同时计算前缀概率的概率Earley解析\scite{nlp-analysis-stolcke1995}。

可空规则$B\to\varepsilon$需要单独处理。$B$在位置$j$预测后立即完成，若此时等待$B$的项目尚未全部加入$C_j$，完成就会漏掉它们。Aycock与Horspool的修正是：预测一个可空类别时，同时把等待它的项目的点直接越过它\scite{nlp-analysis-aycock2002}。单分支循环（如$A\to B$、$B\to A$）会使森林含环，Viterbi计算需要在环上迭代或预先消除。Earley在一般文法下的最坏时间为$O(n^3)$，在无歧义文法上为$O(n^2)$，对许多确定性文法接近线性。它与CKY的主要区别在于不必改写文法，以及只在可由左侧推出的位置上预测；共同代价是歧义越多，森林越大。

\subsubsection{神经跨度成分解析}
\label{sec:nlp-analysis-neural-cky}

父节点标注和词汇化都在手工设计“规则概率应当依赖什么”。神经跨度解析换一条路：去掉显式文法，用编码整句得到的上下文表示直接给每个跨度和类别评分，再用CKY保证输出是合法树。它属于第\ref{sec:nlp-analysis-paradigm-span}小节的跨度范式，评分对象是带标签的跨度$(i,j,\ell)$。

跨度表示的一般形式是式~\eqref{eq:nlp-analysis-span-representation}的端点拼接、内部注意力汇聚与宽度嵌入。成分解析中常用一种更紧凑的形式。用双向LSTM在句首、句末加哨兵符号后编码整句，在边界位置$t=0,\ldots,n$上取前向状态$\bm f_t\in\mathbb R^d$与后向状态$\bm b_t\in\mathbb R^d$：$\bm f_t$读完了边界$t$左侧的词，$\bm b_t$读完了右侧的词。跨度$(i,j)$的表示定义为
\begin{equation}
 \bm r_{i,j}=[\bm f_j-\bm f_i;\ \bm b_i-\bm b_j]\in\mathbb R^{2d}.
 \label{eq:nlp-analysis-span-difference}
\end{equation}
$\bm f_j$与$\bm f_i$的差近似刻画前向网络读入$w_{i+1:j}$这几个词带来的变化，$\bm b_i-\bm b_j$是后向网络读同一段时的变化，因此不需要额外的宽度嵌入或内部汇聚。图~\ref{fig:nlp-analysis-span-difference}画出了“学生代表”$(1,3)$的四个状态。这种差分是Wang与Chang在依存解析中提出的LSTM-Minus特征的双向形式；Cross与Huang把它用于转移式成分解析，Stern等用于图表解析\cite{nlp-analysis-cross2016,nlp-analysis-stern2017}。Kitaev与Klein把编码器换成自注意力，把每个位置的输出向量拆成两半分别充当$\bm f_t$与$\bm b_t$，差分形式不变\scite{nlp-intro-kitaev2018}。

\begin{figure*}[htbp]
  \centering
  \begin{tikzpicture}[
  x=1mm,y=1mm,figure picture,
  word/.style={figure node,minimum height=7mm,minimum width=16mm,
    draw=FigureInput,fill=FigureInputFill},
  post/.style={circle,draw=FigureStroke,fill=white,line width=.75pt,
    minimum size=5.5mm,inner sep=0pt,font=\figurefont\scriptsize},
  state/.style={figure node,minimum height=6mm,minimum width=11mm,
    inner xsep=1pt,font=\figurefont\scriptsize},
  hot/.style={state,draw=FigureOutput,fill=FigureOutputFill,line width=1pt},
  fwd/.style={figure flow},
  bwd/.style={figure flow,dash pattern=on 2pt off 1.5pt},
  note/.style={font=\figurefont\scriptsize,text=FigureMuted}
]
  % 词与边界位置
  \foreach \k/\w in {1/新,2/学生,3/代表}{
    \node[word] (w\k) at (\k*30-15,0) {\w};
  }
  \foreach \k in {0,1,2,3}{
    \node[post] (p\k) at (\k*30,-11) {\k};
  }
  \node[note,anchor=east] at (-6,-11) {边界位置};

  % 前向与后向状态
  \foreach \k in {0,1,2,3}{
    \node[state] (f\k) at (\k*30,13) {$\bm f_{\k}$};
    \node[state] (b\k) at (\k*30,25) {$\bm b_{\k}$};
  }
  \node[hot] at (f1) {$\bm f_{1}$};
  \node[hot] at (f3) {$\bm f_{3}$};
  \node[hot] at (b1) {$\bm b_{1}$};
  \node[hot] at (b3) {$\bm b_{3}$};
  \foreach \a/\b in {0/1,1/2,2/3}{
    \draw[fwd] (f\a) -- (f\b);
    \draw[bwd] (b\b) -- (b\a);
  }
  \node[note,anchor=east] at (-6,13) {前向（读完左侧）};
  \node[note,anchor=east] at (-6,25) {后向（读完右侧）};

  % 待评分跨度
  \draw[figure subgroup,draw=FigureOutput] (24.5,-17) rectangle (95.5,6);
  \node[note,anchor=north] at (60,-17.5) {跨度 $(1,3)$：学生代表};

  % 差分表示：四个高亮状态参与差分
  \node[figure node,draw=FigureOutput,fill=white,minimum width=96mm,
    minimum height=8mm] (r) at (45,40)
    {$\bm r_{1,3}=[\bm f_{3}-\bm f_{1};\ \bm b_{1}-\bm b_{3}]\in\mathbb R^{2d}$};
  \foreach \k in {1,3}{\draw[figure link] (b\k.north) -- (b\k.north |- r.south);}
  \node[note,anchor=south] at (45,44.5) {再经 MLP 得到各标签分数 $s(1,3,\ell)$};
\end{tikzpicture}

  \caption{边界差分的跨度表示。下排圆圈是词之间的边界位置0至3；中排前向状态$\bm f_t$已读完边界$t$左侧的词，上排后向状态$\bm b_t$已读完右侧的词。跨度$(1,3)$只用高亮的四个状态：前向差$\bm f_3-\bm f_1$与后向差$\bm b_1-\bm b_3$都对应读入“学生代表”带来的变化。}
  \label{fig:nlp-analysis-span-difference}
\end{figure*}

标签分数由一个单隐层前馈网络给出，这里采用Kitaev与Klein加入层归一化的写法：
\begin{equation}
 \bm s_{i,j}=\bm W_2
  \operatorname{ReLU}\!\bigl(
   \operatorname{LayerNorm}(\bm W_1\bm r_{i,j}+\bm b_1)
  \bigr)+\bm b_2\in\mathbb R^{L},
 \label{eq:nlp-analysis-span-label}
\end{equation}
其中$\bm W_1\in\mathbb R^{d_h\times2d}$，$\bm W_2\in\mathbb R^{L\times d_h}$，$L$是标签数，$\bm s_{i,j}$的第$\ell$个分量记为$s(i,j,\ell)$。一句话共有$n(n+1)/2$个跨度，全部分数可组织成形状为$(n+1)\times(n+1)\times L$、只用上三角的张量，一次批量计算。

树的分数是所含跨度分数之和，$S(T)=\sum_{(i,j,\ell)\in T}s(i,j,\ell)$。单分支链折叠成一个复合标签，例如S与VP的链记作S+VP；二叉化引入的辅助节点使用空标签$\varnothing$，固定$s(i,j,\varnothing)=0$。这样同一棵原始树的不同二叉化不会仅因辅助节点而得到不同总分；空标签的跨度在恢复时删去。若标签选择只依赖跨度本身，令$q(i,j)=\max_\ell s(i,j,\ell)$，最优树由式~\eqref{eq:nlp-analysis-span-cky}的跨度CKY求出。在本节的边界编号下，$C(i,j)$表示区间$[i,j)$上最优子树的总分，递推为
\begin{align}
 C(i,i+1)&=q(i,i+1),\nonumber\\
 C(i,j)&=q(i,j)+\max_{i<k<j}[C(i,k)+C(k,j)].
 \label{eq:nlp-analysis-neural-cky}
\end{align}
状态里不再有类别，复杂度从$O(n^3|R_2|)$降到$O(n^3)$加上$O(n^2L)$的标签取最大。若评分含父子类别耦合，就须像式~\eqref{eq:nlp-analysis-pcfg-cky}那样把类别留在状态中。

对同一“新／学生／代表”，设单词叶分数为0，两个长度为2的跨度的最佳NP分数分别为5和6，根NP为1，其他非空类别都低于这些值（教学数值）。独立选择高分跨度会同时取$[0,2)$与$[1,3)$，总分$5+6+1=12$，但二者交叉。图表计算$C(0,2)=5$、$C(1,3)=6$，根比较$1+0+6=7$与$1+5+0=6$，选择右分组。

训练使用式~\eqref{eq:nlp-analysis-structured-objectives}中的结构化间隔损失，具体为
\begin{equation}
 \mathcal L=\max_T[S(T)+\Delta(T,T^\star)]-S(T^\star),
 \label{eq:nlp-analysis-parse-hinge}
\end{equation}
$T^\star$是参考树，$\Delta(T,T^\star)$是$T$中不属于$T^\star$的带标签跨度个数。候选包含$T^\star$本身，所以损失非负。$\Delta$按跨度分解，代价增强解码只需把每个跨度分数改为$s(i,j,\ell)+\mathbb 1[(i,j,\ell)\notin T^\star]$再跑一遍CKY，就能找到“最违反间隔”的树；梯度提高参考跨度的分数、降低这棵竞争树中错误跨度的分数。损失为0意味着参考树比每棵其他树至少高出它们的错误跨度数。

与PCFG相比，神经跨度解析的收益来自整句上下文：跨度分数可以利用远处的词和字符形态，不再受规则局部独立的限制。Kitaev与Klein沿用同一图表解码、把编码器换成自注意力，不用外部预训练表示时单模型在宾州树库WSJ测试集上的F1为93.55\scite{nlp-intro-kitaev2018}。代价是要有树库监督，且精确解码依赖分数按跨度分解；不同跨度之间的相互作用无法进入这个CKY。

\subsubsection{移进--归约成分解析}
\label{sec:nlp-analysis-shift-reduce}

图表方法同时保存$O(n^2)$个区间的分数，解码需要$O(n^3)$时间。长句或流式输入时，可以改用第\ref{sec:nlp-analysis-paradigm-transition}小节的转移系统，按动作逐步构树。移进--归约成分解析和下一小节的arc-standard依存解析属于同一范式：状态由栈、缓冲区和已建结构组成，每个动作有前提与效果，参考结构经oracle转成动作序列，模型在合法动作上做softmax，按式~\eqref{eq:nlp-analysis-transition-loss}以交叉熵训练。两者的区别只在动作集合与栈上存放的对象：这里栈上是带类别的子树，依存解析的栈上是词。

Sagae与Lavie、Zhu等的成分系统使用表~\ref{tab:nlp-analysis-shift-reduce-actions}中的动作\scite{nlp-analysis-sagae2005,nlp-analysis-zhu2013}。归约动作中的L或R指明新节点的中心词取自左孩子还是右孩子，这是为了让特征能读取子树的中心词；X是新节点的类别。参考树需要先二叉化，并记录单分支的位置。

\begin{table*}[htbp]
  \centering
  \begin{tabularx}{\linewidth}{@{}l X X@{}}
    \toprule
    动作&前提&效果\\
    \midrule
    SHIFT&缓冲区非空&把缓冲区首词连同其词性作为叶节点压栈\\
    REDUCE-L-X / REDUCE-R-X&栈中至少两个子树，且合成后满足二叉化约束&弹出栈顶两个子树，合成类别为X的节点后压栈；中心词取自左（L）或右（R）孩子\\
    UNARY-X&栈非空，且栈顶刚才没有执行过单分支&弹出栈顶子树，套一层类别为X的节点后压栈\\
    FINISH&缓冲区为空，栈中只剩一个可作根的子树&终止，输出栈中的树\\
    \bottomrule
  \end{tabularx}
  \caption{移进--归约成分解析的动作。前提保证每一步都可执行，且不会出现无限单分支；它们只保证结构合法，不保证选中的分组正确。}
  \label{tab:nlp-analysis-shift-reduce-actions}
\end{table*}

对右分组，“新／学生／代表”的参考轨迹为五个构树动作加终止：
\[
\begin{array}{llll}
 \text{步}&\text{动作}&\text{栈}&\text{缓冲区}\\
 1&\mathrm{SHIFT}&[\mathrm A(\text{新})]&[\text{学生},\text{代表}]\\
 2&\mathrm{SHIFT}&[\mathrm A,\ \mathrm N(\text{学生})]&[\text{代表}]\\
 3&\mathrm{SHIFT}&[\mathrm A,\ \mathrm N,\ \mathrm N(\text{代表})]&[\,]\\
 4&\mathrm{REDUCE\mbox{-}R\mbox{-}NP}&[\mathrm A,\ \mathrm{NP}(\text{学生代表})]&[\,]\\
 5&\mathrm{REDUCE\mbox{-}R\mbox{-}NP}&[\mathrm{NP}(\text{新学生代表})]&[\,]
\end{array}
\]
随后FINISH。左分组的轨迹则在第2步之后立即归约：SHIFT、SHIFT、REDUCE-R-NP（新学生）、SHIFT、REDUCE-R-NP。两条轨迹在第3步分叉：读完“新／学生”之后，选择继续移进还是立即归约，决定了最终分组。这个选择是不可撤销的；一旦把“新学生”归约成NP，后续任何动作都不能再把它拆开。

动作分类器读取栈顶若干子树的类别与中心词、缓冲区前几个词等信息。设分叉处“继续移进”与“立即归约”的局部概率为$.45$与$.55$，两条后续路径的条件概率分别为$.9$与$.4$，完整轨迹概率为$.405$与$.22$（教学数值）。贪心解码在分叉处选择归约，得到左分组；束宽为2的束搜索保留两条，最终按完整概率选回右分组。束搜索仍是近似，有限束宽不能保证找到全局最优。单分支使不同树的动作序列长度不同，比较束中候选时会偏向较短的轨迹；Zhu等为此加入空动作IDLE，让所有候选在相同步数后比较。动态oracle与全局归一化等通用的训练改进见第\ref{sec:nlp-analysis-paradigm-transition}小节。

移进--归约的时间与句长成线性，每步只读取有限的状态特征，适合需要快速解析的场景。代价是早期错误无法撤销，且没有CKY那样的全局最优保证。Cross与Huang的转移式成分解析把动作分类与上一小节的跨度差分结合：动作分类器读取栈顶跨度的差分表示，并给出可证明最优的动态oracle\scite{nlp-analysis-cross2016}。

\subsubsection{序列化生成式解析}
\label{sec:nlp-analysis-seq2seq-parsing}

前三种方法都为树设计了专门的搜索。另一条路线完全不设计搜索：把树写成括号序列，交给通用的序列到序列模型生成，属于第\ref{sec:nlp-analysis-paradigm-generation}小节的约束生成范式。Vinyals等按深度优先顺序线性化树，词本身不出现在输出中，叶节点上的词性标签统一替换为占位符XX，例如右分组写成
\[
 \texttt{(NP XX (NP XX XX )NP )NP},
\]
右括号带上类别，便于模型记住当前要关闭哪一层\scite{nlp-analysis-vinyals2015grammar}。编码器读入原句，带注意力的LSTM解码器逐个生成括号与XX，训练目标是参考序列每一步的交叉熵。推断用束搜索；解码完成后检查括号是否配对、XX的个数是否等于词数，不合法时补全或删除括号，再把第$k$个XX对应回第$k$个词。原文中，只用WSJ训练时无注意力的模型F1不到70，加入注意力后才可用；再用大量自动解析的句子扩充训练集，达到了与专门解析器相当的水平。

这条路线的收益是接口统一：同一个序列到序列模型可以解析、翻译或生成摘要，不需要为树设计特征和动态规划。代价是合法性不再由构造保证。括号配对是上下文无关语言的性质，可以用第\ref{sec:nlp-inference-grammar}小节的文法约束解码在每步屏蔽不能完成合法括号串的词元，从源头排除不配对的输出；叶节点个数与原句一致则需要在约束中计数。

大语言模型可以在提示中给出括号格式和少量示例，直接输出成分树。这时还会出现一种序列化解析特有的错误：模型在括号中改写、漏掉或增加原句的词。交付前应逐一核对叶序列与原句的词序列，并用PARSEVAL在固定口径上评价；没有树库微调时，模型输出的类别体系可能与目标树库不一致，需要映射。

表~\ref{tab:nlp-analysis-constituency-methods}在连续成分树这一共同输出下比较五条路线。

\begin{table*}[htbp]
  \centering
  \begin{tabularx}{\linewidth}{@{}l XXXX@{}}
    \toprule
    方法&解决的问题&核心机制&推断与复杂度&新代价\\
    \midrule
    PCFG--CKY&概率化歧义选择&规则概率之积；区间动态规划&精确，$O(n^3|R_2|)$&独立假设强；需二叉化\\
    Earley&一般文法不必改写&带点项目；预测、扫描、完成&精确，最坏$O(n^3)$&可空规则与森林处理复杂\\
    神经跨度--CKY&利用整句上下文&边界差分表示；跨度标签分数&精确，$O(n^3)$&需树库；分数须按跨度分解\\
    移进--归约&降低解析时间&栈上子树与归约动作&近似（贪心或束），$O(n)$&早期错误不可撤销\\
    序列化生成&统一序列接口&括号线性化；seq2seq生成&近似（束），$O(n)$步&合法性靠约束或后处理\\
    \bottomrule
  \end{tabularx}
  \caption{成分解析方法对照。复杂度中的$n$为句长；输出均须去掉辅助节点、展开复合类别后，再按同一PARSEVAL口径评价。}
  \label{tab:nlp-analysis-constituency-methods}
\end{table*}

\subsection{依存结构与解析}
\label{sec:nlp-analysis-dependency}

基本依存分析以词$1,\ldots,n$和人工根0为节点，边$h\to d$表示中心词$h$支配依存词$d$，边上另有关系标签。本节要求：每个真实词恰有一个中心词，人工根没有入边，所有真实词可从根到达，并且人工根恰好连接一个真实词。满足这些条件的边集合是以0为根的有向生成树。\term{投射性}（projectivity）进一步要求每个词及其全部后代在句中占据连续区间；把人工根放在最左侧时，等价于画在句子上方的依存边互不交叉。依存标注资源以通用依存（Universal Dependencies, UD）为代表，它为多种语言提供共享的关系标签体系\scite{nivre2020universal}。

评价按词计分。\term{无标签依存准确率}（unlabeled attachment score, UAS）是中心词预测正确的词所占比例，\term{带标签依存准确率}（labeled attachment score, LAS）还要求关系标签正确。本节手算计入指向根词的根弧、示例不含标点；实际评测需声明是否排除标点，并在不同系统间保持分母一致。

依存解析的挑战有四项：输出必须满足单头、单根、无环的树约束；自由语序和长距离依赖会产生非投射结构；关系标签需要与弧一起预测；长距离弧的错误率明显高于短弧。方法谱系为：转移式解析线性时间但贪心；图式解析把树的分数分解到弧上，在一阶分解下全局精确，其中Eisner算法限于投射树，Chu--Liu/Edmonds允许非投射；双仿射网络用神经弧评分代替手工特征，接图解码；增强依存图取消单头限制，表达共享论元与省略。每种方法都要说明它在什么结构空间内保证合法。

\subsubsection{arc-standard与arc-eager转移解析}
\label{sec:nlp-analysis-arc-transition}

两种转移系统都维护栈$\sigma$、缓冲区$\beta$和弧集合$A$，从$([0],[1,\ldots,n],\varnothing)$开始，到缓冲区为空、栈中只剩根0时结束\scite{nlp-analysis-nivre2003,nlp-analysis-nivre2008}。它们的差别在于连弧的时机，表~\ref{tab:nlp-analysis-arc-systems}列出两套动作。记栈顶为$s_0$、次顶为$s_1$，缓冲区首词为$b_0$。

\begin{table*}[htbp]
  \centering
  \begin{tabularx}{\linewidth}{@{}l l X X@{}}
    \toprule
    系统&动作&前提&效果\\
    \midrule
    arc-standard&SHIFT&缓冲区非空&$b_0$压栈\\
    &LEFT-ARC-$r$&栈中至少两个元素，$s_1\neq0$&加弧$s_0\xrightarrow{r}s_1$，移除$s_1$\\
    &RIGHT-ARC-$r$&栈中至少两个元素；$s_1=0$时缓冲区须为空&加弧$s_1\xrightarrow{r}s_0$，移除$s_0$\\
    \midrule
    arc-eager&SHIFT&缓冲区非空&$b_0$压栈\\
    &LEFT-ARC-$r$&$s_0\neq0$，且$s_0$尚无中心词&加弧$b_0\xrightarrow{r}s_0$，弹出$s_0$\\
    &RIGHT-ARC-$r$&缓冲区非空；$s_0=0$时根尚无孩子&加弧$s_0\xrightarrow{r}b_0$，$b_0$压栈\\
    &REDUCE&$s_0$已有中心词&弹出$s_0$\\
    \bottomrule
  \end{tabularx}
  \caption{两种依存转移系统的动作。arc-standard只在栈顶两个元素之间连弧，被连上的依存词立即移除，因此必须等它的孩子全部连完；arc-eager在栈顶与缓冲区首词之间连弧，右依存词一出现就连上，之后用REDUCE清理。两套系统都只能产生投射树。}
  \label{tab:nlp-analysis-arc-systems}
\end{table*}

对“小王$_1$／阅读$_2$／通知$_3$”，arc-standard的完整轨迹为：SHIFT，SHIFT，LEFT-ARC-nsubj建立$2\to1$，SHIFT，RIGHT-ARC-obj建立$2\to3$，RIGHT-ARC-root建立$0\to2$。栈依次为
\[
 [0]\to[0,1]\to[0,1,2]\to[0,2]
 \to[0,2,3]\to[0,2]\to[0].
\]
根弧在最后才建立，此时“阅读”的子树已经完整。arc-eager的轨迹为：SHIFT，LEFT-ARC-nsubj建立$2\to1$，RIGHT-ARC-root建立$0\to2$，RIGHT-ARC-obj建立$2\to3$，REDUCE，REDUCE。它在读入“阅读”时就连上了根，比arc-standard早；这种“尽早连弧”让arc-eager的部分结构更早可用，也让它成为最早得到动态oracle的系统之一，Goldberg与Nivre证明了它的弧可以逐条判断能否仍被到达\scite{nlp-analysis-goldberg2012}。

两套系统各有$n$次移进和$n$次连弧，arc-eager另有若干次REDUCE，动作数都与句长成线性。没有额外动作时它们只能产生投射树：被移除或弹出的词不再回到栈中，交叉弧所需的连接顺序无法实现。Nivre在arc-standard中加入SWAP动作，把次栈顶的词移回缓冲区，从而改变词进入栈的顺序；按重排后的顺序连出的投射弧，在原词序中可以交叉，解析时间在期望意义上仍是线性的\scite{nlp-analysis-nivre2009}。

\paragraph{Chen与Manning的前馈动作分类器}
早期转移解析器用人工组合的稀疏特征给动作打分，特征模板多、组合稀疏，特征提取占去大部分解析时间。Chen与Manning把离散特征换成稠密嵌入，用一个小型前馈网络预测arc-standard的动作\scite{v9a-chen2014}。输入取三组元素：栈顶三个词与缓冲区前三个词；栈顶两个词的最左、最右两个孩子；以及这两个词的最左孩子的最左孩子、最右孩子的最右孩子。词特征共18个，对应的词性特征18个，除栈与缓冲区6个词之外的12个元素再取弧标签，共48个特征。每个特征查表得到50维嵌入，按组拼接成$\bm x^{w}\in\mathbb R^{18\cdot50}$、$\bm x^{t}\in\mathbb R^{18\cdot50}$、$\bm x^{l}\in\mathbb R^{12\cdot50}$，隐藏层与输出层为
\begin{equation}
 \bm h=\bigl(\bm W_1^{w}\bm x^{w}+\bm W_1^{t}\bm x^{t}+\bm W_1^{l}\bm x^{l}+\bm b_1\bigr)^{3},\qquad
 p(a\mid c)=\operatorname{softmax}(\bm W_2\bm h)_a,
 \label{eq:nlp-analysis-chen-manning}
\end{equation}
立方逐元素计算，$\bm h\in\mathbb R^{200}$，$\bm W_2$的行数是动作数$2N_r+1$（$N_r$为关系标签数）。立方激活展开后含有三个输入分量的乘积项，相当于自动学习特征两两、三三组合，替代了人工组合模板。训练以oracle动作的交叉熵加$L_2$正则，嵌入用预训练词向量初始化。实现上，作者对高频词预先计算嵌入与$\bm W_1$对应列的乘积并缓存，解析时只需查表相加。在宾州树库（CoNLL依存转换）的开发集上，它以92.2的UAS每秒解析约一千句，准确率与速度同时高于作者实现的稀疏特征转移解析器。

\paragraph{栈LSTM}
Chen与Manning的特征仍是固定窗口：栈中更深处的词和已归约子树的内部结构看不到。Dyer等用\term{栈LSTM}（stack LSTM）编码整个栈\scite{nlp-analysis-dyer2015}。栈LSTM是带栈顶指针的LSTM：压栈时以当前栈顶的隐状态为前一状态计算新状态，弹栈时只把指针退回上一个状态，因此栈顶的输出始终概括了栈中全部元素。解析器用三个这样的LSTM分别编码栈$\sigma$、缓冲区$\beta$和动作历史，状态表示为
\[
 \bm p_t=\max\{\bm 0,\ \bm W[\bm s_t;\bm b_t;\bm a_t]+\bm d\},
\]
再经softmax给出动作分布。归约时，中心词$\bm h$、依存词$\bm d$与关系嵌入$\bm e_r$组合为$\bm c=\tanh(\bm U[\bm h;\bm d;\bm e_r]+\bm e)$，作为新子树压回栈中。与固定窗口相比，栈LSTM能读到完整的解析历史，代价是每步都要更新递归状态。

转移解析的共同限制是贪心：每步只看局部动作分数，一旦错误弹出某个词，涉及它的后续弧就不可达。动态oracle、束搜索与全局归一化的改进见第\ref{sec:nlp-analysis-paradigm-transition}小节。

\subsubsection{Eisner投射树解码}
\label{sec:nlp-analysis-eisner}

图式解析先给每条候选弧$h\to d$打分$s(h,d)$，再在合法树中找总分最大的一棵，$\max_{T}\sum_{(h,d)\in T}s(h,d)$。这种一阶分解是第\ref{sec:nlp-analysis-paradigm-graph}小节图范式的典型例子。逐词选最高入边可能成环；即使打断环，结果也可能不投射。Eisner算法在连续区间上递推，直接在投射树空间内求精确最优\scite{nlp-analysis-eisner1996}。

Eisner算法的四类状态与递推已在第\ref{sec:nlp-analysis-paradigm-graph}小节式~\eqref{eq:nlp-analysis-eisner}给出，本小节沿用其记号：闭区间$[i,j]$表示从词$i$到词$j$，与成分解析的边界区间相区别；完整项$C$的根在区间一端，它在区间内一侧的子孙都已收齐；不完整项$I$刚加入端点弧，被支配的端点还要收齐朝向区间内的孩子。实现时把非初值条目设为$-\infty$，按区间长度从短到长、同一长度先算不完整项后算完整项，并记录每次取最大值的分割点和状态类型，供回溯使用。本小节补上图范式中没有展开的两件事：单根约束和一个完整的手算例子。

单根约束在图表完成之后单独处理。一棵投射树被根词$r$分成左右两半：左半是以$r$为右端根的完整区间$C^{\leftarrow}[1,r]$，右半是以$r$为左端根的完整区间$C^{\to}[r,n]$。因此最优值为
\[
 \max_{1\leq r\leq n}
 \bigl[s(0,r)+C^{\leftarrow}[1,r]+C^{\to}[r,n]\bigr],
\]
并保存最优的根词$r$。若把人工根当作位置0直接放进图表递推，它可能连接多个真实词，得到的是一片森林挂在根下，不满足单根要求。

给“小王／阅读／通知”构造全部有限的弧分数（教学数值）：
\[
 s(0,2)=5,\quad s(2,1)=4,\quad s(2,3)=3,\quad
 s(3,1)=5,\quad s(1,3)=6,
\]
其余弧禁止，包括自环和指向人工根的弧。逐词取最高入边得到$3\to1$、$0\to2$、$1\to3$，总分16，但1与3之间成环。长度为2的图表项为$I^{\leftarrow}[1,2]=C^{\leftarrow}[1,2]=4$与$I^{\to}[2,3]=C^{\to}[2,3]=3$，相反方向均为$-\infty$。弧$1\to3$虽有6分，但$I^{\to}[1,3]$的两个分割分别需要$C^{\leftarrow}[2,3]$与$C^{\to}[1,2]$，即被禁止的$3\to2$或$1\to2$，这一不完整项不可达。唯一可行的根是2，总分$5+4+3=12$；回溯得到$2\to1$、$2\to3$与$0\to2$，正是图~\ref{fig:nlp-analysis-syntax}的无标签树。

Eisner算法的时间为$O(n^3)$，图表与回溯指针为$O(n^2)$。它本身不学习参数：弧分数可以来自结构感知器或大间隔训练\cite{mcdonald-etal-2005-non}，训练时用当前分数解码出竞争树，再比较参考树与竞争树的弧特征；也可以来自下文双仿射网络的局部交叉熵。精确性只针对一阶可加的弧分数和投射树空间；加入兄弟、祖父等高阶部件后，状态要扩大，复杂度随之上升。

\subsubsection{Chu--Liu/Edmonds非投射树解码}
\label{sec:nlp-analysis-mst}

自由语序和远距离依赖会产生交叉弧。取消投射性之后，解码问题变成在带权有向图上求以0为根的\term{最大生成树形图}（maximum spanning arborescence）：每个真实词恰有一条入边，无环且可从根到达。McDonald等把非投射依存解析写成这一问题\cite{v9a-mcdonald2005}，求解用\term{Chu--Liu/Edmonds算法}（Chu--Liu/Edmonds algorithm）\scite{nlp-analysis-chu1965,nlp-analysis-edmonds1967}。

算法的思路是：先让每个真实词贪心地选最大入边；若没有环，这已经是上界，直接返回。若有环，最优树一定恰好从环外进入环一次，于是把环收缩成一个点，进入这个点的弧要反映“替换掉环中哪条边”的代价。设环$Q$中节点$v$的已选前驱为$a(v)$，环内总分$w(Q)=\sum_{v\in Q}s(a(v),v)$。外部弧$u\to v$进入环时替换掉原入边$a(v)\to v$，收缩后的分数为
\begin{equation}
 s'(u,Q)=\max_{v\in Q}
   [s(u,v)-s(a(v),v)+w(Q)].
 \label{eq:nlp-analysis-mst-contract}
\end{equation}
出环弧取$s'(Q,u)=\max_{v\in Q}s(v,u)$。这就是第\ref{sec:nlp-analysis-paradigm-graph}小节给出的收缩规则，只是这里把常数$w(Q)$也加回，使收缩图上的树分直接等于原图上展开后的树分。两种最大值都要保存原始端点，否则展开时无法恢复词之间的弧。算法~\ref{alg:nlp-analysis-cle}给出递归形式。

\begin{algorithm}[Chu--Liu/Edmonds最大生成树形图]
  \label{alg:nlp-analysis-cle}
  \begin{algorithmic}[1]
    \Require 节点集合$V=\{0,1,\ldots,n\}$，弧分数$s(u,v)$（禁止的弧为$-\infty$），根0
    \Ensure 以0为根、总分最大的生成树形图$T$
    \Procedure{CLE}{$V,s$}
      \For{每个$v\in V\setminus\{0\}$}
        \State $a(v)\gets\argmax_{u\neq v}s(u,v)$\Comment{贪心选最大入边}
      \EndFor
      \State $M\gets\{(a(v),v):v\neq0\}$
      \If{$M$中无环}
        \State \Return $M$
      \EndIf
      \State 任取$M$中的一个环$Q$，计算$w(Q)=\sum_{v\in Q}s(a(v),v)$
      \State 把$Q$收缩为新节点$q$，得到$V'=(V\setminus Q)\cup\{q\}$
      \For{每个$u\notin Q$}
        \State $s'(u,q)\gets\max_{v\in Q}[s(u,v)-s(a(v),v)]+w(Q)$，记录取到最大值的$v^\ast(u)$
        \State $s'(q,u)\gets\max_{v\in Q}s(v,u)$，记录取到最大值的源点$v^{\circ}(u)$
      \EndFor
      \State 环外节点之间的弧分数不变；$T'\gets$\Call{CLE}{$V',s'$}
      \State 设$T'$中进入$q$的弧为$(u_0,q)$：加入$u_0\to v^\ast(u_0)$，并加入$Q$中除$a(v^\ast(u_0))\to v^\ast(u_0)$以外的全部弧
      \State 把$T'$中每条$(q,u)$换回$v^{\circ}(u)\to u$，其余弧原样保留
      \State \Return 展开后的弧集合$T$
    \EndProcedure
  \end{algorithmic}
\end{algorithm}

沿用上一组分数，环$Q=\{1,3\}$的总分为$5+6=11$。从2进入1的收缩分数为$4-5+11=10$，进入3的为$3-6+11=8$，从0进入的弧都被禁止。收缩图中$q$的最大入边是$2\to q$（10分），2的最大入边是$0\to2$（5分），无环，递归结束。展开时用$2\to1$替换环边$3\to1$，保留$1\to3$，得到总分$5+4+6=15$的树：$0\to2$、$2\to1$、$1\to3$。这棵树无环，但弧$0\to2$与$1\to3$交叉，因而不投射；Eisner排除了它，最高只能得到12分。两者的分差完全来自允许的树空间不同，不说明15分的那棵树在语言上更合理。

如第\ref{sec:nlp-analysis-paradigm-graph}小节所述，朴素实现每次收缩花$O(n^2)$、最多收缩$n$次，总时间$O(n^3)$，Tarjan的实现在稠密图上为$O(n^2)$。上述算法允许人工根连接多个真实词。若规范只允许一个根词，最直接的做法是枚举根词$r$，只保留$0\to r$这一条根出边，分别求解后取最高分，代价是多出$n$倍的计算；Zmigrod等给出了在$O(n^2)$内直接施加单根约束的做法\scite{nlp-analysis-zmigrod2020}。

非投射解码并非总是更好。McDonald等在捷克语PDT上比较两种解码，生成树解码的中心词准确率为84.4，投射解码为83.3；而在依存树全部投射的英语宾州树库转换数据上，投射解码反而略高\scite{v9a-mcdonald2005}。捷克语中非投射弧不足2\%，收益却已可见；在没有交叉弧的数据上，更大的树空间就没有可恢复的对象。选哪一种应看目标树库中非投射结构的比例。

\subsubsection{双仿射依存解析}
\label{sec:nlp-analysis-biaffine}

上面两个解码器都假定弧分数已经给定。Chen与Manning的分类器给的是动作分数，不是弧分数；早期图式解析器用人工特征给弧打分。Kiperwasser与Goldberg先用BiLSTM的输出向量代替人工特征，再用前馈网络给每对词打分\scite{nlp-analysis-kiperwasser2016}。Dozat与Manning的\term{双仿射依存解析}（biaffine dependency parsing）进一步改变评分函数：先把每个词分别投影成“作为依存词”和“作为中心词”的两种表示，再用双仿射形式一次算出全部$n\times n$的弧分数\scite{nlp-analysis-dozat2017}。图~\ref{fig:nlp-analysis-biaffine}画出了整体结构。

\begin{figure*}[htbp]
  \centering
  \begin{tikzpicture}[
  x=1mm,y=1mm,figure picture,
  input/.style={figure node,minimum height=7mm,minimum width=22mm,
    draw=FigureInput,fill=FigureInputFill},
  layer/.style={figure node,minimum height=8mm},
  mlp/.style={figure node,minimum height=8mm,minimum width=25mm},
  hot/.style={figure node,minimum height=9mm,draw=FigureOutput,
    fill=FigureOutputFill,line width=1pt},
  flow/.style={figure flow},
  shape note/.style={font=\figurefont\scriptsize,text=FigureMuted}
]
  % 输入与编码
  \node[input] (x) at (60,0) {词与词性嵌入 $\bm x_0,\ldots,\bm x_n$};
  \node[layer,minimum width=132mm] (enc) at (60,14) {多层 BiLSTM：$\bm R\in\mathbb R^{(n+1)\times 2h}$};
  \draw[flow] (x) -- (enc);

  % 四个投影
  \node[mlp] (ad) at (8,31) {$\mathrm{MLP}^{\mathrm{arc\mbox{-}dep}}$};
  \node[mlp] (ah) at (42,31) {$\mathrm{MLP}^{\mathrm{arc\mbox{-}head}}$};
  \node[mlp] (rd) at (78,31) {$\mathrm{MLP}^{\mathrm{rel\mbox{-}dep}}$};
  \node[mlp] (rh) at (112,31) {$\mathrm{MLP}^{\mathrm{rel\mbox{-}head}}$};
  \foreach \m in {ad,ah,rd,rh}{\draw[flow] (enc.north -| \m.south) -- (\m.south);}

  % 两个双仿射头
  \node[hot,minimum width=58mm] (sa) at (25,52)
    {$\bm S^{\mathrm{arc}}=\bm H^{(h)}\bm U\bm H^{(d)\top}+(\bm H^{(h)}\bm u)\bm 1^{\top}$};
  \node[hot,minimum width=58mm] (sr) at (95,52)
    {关系双仿射：$\bm S^{\mathrm{rel}}\in\mathbb R^{(n+1)\times(n+1)\times R}$};
  \draw[flow] (ad.north) -- (ad.north |- sa.south);
  \draw[flow] (ah.north) -- (ah.north |- sa.south);
  \draw[flow] (rd.north) -- (rd.north |- sr.south);
  \draw[flow] (rh.north) -- (rh.north |- sr.south);

  % 解码
  \node[layer,minimum width=46mm] (tree) at (25,72) {Eisner 或 CLE 选树};
  \node[layer,minimum width=46mm] (lab) at (95,72) {逐弧取最高关系};
  \draw[flow] (sa) -- (tree);
  \draw[flow] (sr) -- (lab);
  \draw[flow] (tree) -- node[above,shape note] {弧集合} (lab);
\end{tikzpicture}

  \caption{双仿射依存解析器的结构。BiLSTM输出经四个独立的MLP投影为弧依存端、弧中心端（维度$d_a$）与关系依存端、关系中心端（维度$d_r$）表示；弧头用双仿射得到$(n+1)\times(n+1)$的弧分数矩阵并交给树解码器，关系头只在选中的弧上读取$R$个关系分数。行下标为中心词，列下标为依存词，第0行对应人工根。}
  \label{fig:nlp-analysis-biaffine}
\end{figure*}

设句子连同人工根共$n+1$个位置，多层BiLSTM输出$\bm R\in\mathbb R^{(n+1)\times2h}$，第$i$行$\bm r_i$是位置$i$的上下文表示。两个MLP把它们投影到较低的维度$d_a$：
\[
 \bm h^{(d)}_i=\mathrm{MLP}^{\mathrm{arc\mbox{-}dep}}(\bm r_i),\qquad
 \bm h^{(h)}_i=\mathrm{MLP}^{\mathrm{arc\mbox{-}head}}(\bm r_i),
\]
按行堆叠成$\bm H^{(d)},\bm H^{(h)}\in\mathbb R^{(n+1)\times d_a}$。弧分数矩阵为
\begin{equation}
 \bm S^{\mathrm{arc}}=\bm H^{(h)}\bm U\bm H^{(d)\top}+(\bm H^{(h)}\bm u)\bm 1^{\top}\in\mathbb R^{(n+1)\times(n+1)},
 \label{eq:nlp-analysis-biaffine}
\end{equation}
其中$\bm U\in\mathbb R^{d_a\times d_a}$，$\bm u\in\mathbb R^{d_a}$，$\bm 1$是全1列向量。元素$S^{\mathrm{arc}}_{i,j}=\bm h^{(h)\top}_i\bm U\bm h^{(d)}_j+\bm h^{(h)\top}_i\bm u$是“$i$作$j$的中心词”的分数。双线性项衡量两端是否相容；第二项只依赖中心词，表示某个词本身有多大可能带依存词，例如动词比限定词更常作中心词。依存端自身的偏置对同一列的所有候选中心词相同，不影响选择，所以省去。

关系标签用另一对MLP投影到维度$d_r$，得到$\bm g^{(d)}_j$与$\bm g^{(h)}_i$，对每个关系$r$计算
\[
 S^{\mathrm{rel}}_{i,j,r}=\bm g^{(h)\top}_i\bm U^{(r)}\bm g^{(d)}_j+\bm w_r^{\top}[\bm g^{(h)}_i;\bm g^{(d)}_j]+b_r,
\]
$\bm U^{(r)}\in\mathbb R^{d_r\times d_r}$。把$R$个$\bm U^{(r)}$叠成三阶张量，全部分数构成$(n+1)\times(n+1)\times R$的张量。原文的BiLSTM为三层、每方向400维，弧投影$d_a=500$，关系投影$d_r=100$。

为什么要先投影再配对，而不直接用$\bm r_i$？BiLSTM的输出要同时服务于弧、关系和其他信息，维度高且混杂；专用MLP把对当前判断有用的部分抽出，并把双线性项的参数从$(2h)^2$降到$d_a^2$。弧与关系使用不同的投影，是因为“谁支配谁”与“是什么关系”依赖的信息不同。

训练时，对每个依存词$j$在第$j$列上做softmax，以参考中心词$i^\star$训练$-\log p(i^\star\mid j)$；关系头只在参考弧上训练正确关系的交叉熵。推断时先把$\bm S^{\mathrm{arc}}$交给Eisner或Chu--Liu/Edmonds选出合法树，再在选中的每条弧上取最高的关系分数。逐列取最大虽然也能给出每个词的中心词，但可能成环或多根；树解码器保证输出合法。

取$d_a=2$、$\bm U=\bm I$、$\bm u=\bm 0$，设“阅读”的中心端表示为$(1,0)$，“小王”与“通知”的依存端表示为$(4,0)$与$(3,0)$，则$S^{\mathrm{arc}}_{2,1}=4$、$S^{\mathrm{arc}}_{2,3}=3$，正好是Eisner例子中的两条弧分数。选定$2\to1$后，若关系头对nsubj、obj、其他的分数为$2,0,-1$，选nsubj。若两个中心词都正确却把$2\to1$标成obj，计入根弧的三条弧中UAS为1，LAS为$2/3$。

与Chen与Manning相比，双仿射解析把特征从固定窗口的嵌入换成整句的BiLSTM表示，把贪心的动作序列换成全局精确的树解码，代价是$O(n^2)$的弧分数和$O(n^3)$（或$O(n^2)$）的解码，不再是线性时间。Dozat与Manning在宾州树库的斯坦福依存版本上得到约95.7的UAS；其消融中去掉低维投影、直接用BiLSTM输出配对时LAS下降，而UAS几乎不变，支持了“关系判断需要专用表示”的解释\scite{nlp-analysis-dozat2017}。

\subsubsection{增强依存图预测}
\label{sec:nlp-analysis-enhanced}

基本树要求每个词只有一个中心词，但并列共享的论元和省略的谓词需要更多连接。\term{增强依存}（enhanced dependencies）在基本树之上按规范增加关系：允许一个词有多个中心词，允许在省略处插入空节点，某些结构中还可能出现环。它不是把树约束放宽后任意加边，哪些边成立仍由标注规范规定\scite{v9a-ud-enhanced}。

例如教学句“小王购买并阅读通知”，基本树有“购买”到“小王”的nsubj、“购买”到“阅读”的conj、“阅读”到“通知”的obj。增强表示把共享主语传播到第二个谓词，增加“阅读”到“小王”的nsubj；规范允许时还传播共享宾语，增加“购买”到“通知”的obj。“小王”因此有两条入边，这在增强图中是正确的。

增强图不能再用单头树解码，常见做法是把它当作多标签的弧预测。保留双仿射的评分结构，对允许的词对$(i,j)$和关系$r$计算$p_{i,j,r}=\operatorname{sigmoid}(S_{i,j,r})$，以增强弧的多热标签训练二元交叉熵，推断时保留超过阈值的弧。设新增两条弧的概率为$.85$与$.80$，反方向的错误候选为$.15$，验证集确定阈值$.5$，则保留前两条。阈值化不保证连通，可以再强制每个词至少保留一条最高分入边，或先用树解码得到基本树、再在其上加边。另一种做法完全不学习：读取基本树中的并列、控制与省略结构，按规范规则传播关系。

省略结构还需要空节点。对“甲买书，乙报纸”，第二分句缺少谓词“买”，增强表示在该处插入一个空节点，把“乙”“报纸”接到它下面。预测时要先确定空节点的位置与身份，再预测与它相连的弧；空节点不在原文中，不能把一个字符跨度强加给它。

评价改用增强弧指标。IWPT 2020共享任务采用ELAS，即增强图中“中心词、关系”对的F1，因为每个词的入边数不再固定为1，UAS与LAS的分母不再适用\scite{nlp-analysis-bouma2020}。表~\ref{tab:nlp-analysis-dependency-methods}按同一输出比较依存方法。转移系统学习动作，Eisner与Chu--Liu/Edmonds在给定弧分数上选树，双仿射网络负责产生弧分数；后三者常组合使用，并不互斥。

\begin{table*}[htbp]
  \centering
  \begin{tabularx}{\linewidth}{@{}l XXXX@{}}
    \toprule
    方法&解决的问题&核心机制&结构空间与复杂度&新代价\\
    \midrule
    arc-standard / arc-eager&线性时间解析&栈、缓冲区与连弧动作&投射树；$O(n)$，贪心或束&错误不可撤销；需交换动作处理非投射\\
    Eisner&投射树上的精确解码&四类区间状态的动态规划&投射树；$O(n^3)$&只支持一阶（或扩展状态后的高阶）分解\\
    Chu--Liu/Edmonds&非投射树的精确解码&贪心入边、收缩环、递归展开&任意生成树形图；$O(n^2)$至$O(n^3)$&失去投射约束的过滤；单根需额外处理\\
    双仿射&弧分数的学习&专用MLP投影与双仿射评分&接上述解码器；分数$O(n^2)$&需整句编码；非线性时间\\
    增强依存&共享论元与省略&多标签弧预测、空节点&一般有向图&无单头约束可用；需新评价口径\\
    \bottomrule
  \end{tabularx}
  \caption{依存解析方法对照。前三行是结构空间与搜索方式，第四行是评分函数，第五行改变了输出结构本身。}
  \label{tab:nlp-analysis-dependency-methods}
\end{table*}

\section{语义分析}
\label{sec:nlp-analysis-semantics}

句法可以判定“小王”是“交给”的主语，却没有说明谁接收了什么，也没有给出一个可以拿去计算或查询的表达。\term{语义分析}（semantic analysis）要先约定意义以什么形式交付：一个词义编号、一组谓词论元、一张概念图，或一个能在指定环境中求值的表达式。这几种形式回答的问题不同，预测出其中一种并不自动得到其他几种。表~\ref{tab:nlp-analysis-semantic-outputs}汇总了本节的五类输出。

\begin{table*}[htbp]
  \centering
  \small
  \setlength{\tabcolsep}{3pt}
  \begin{tabular}{@{}p{15mm}p{24mm}p{32mm}p{24mm}p{27mm}p{28mm}@{}}
    \toprule
    预测对象&教学输入&具体输出&结构检查&意义检查&未覆盖信息\\
    \midrule
    词义&river bank&bank的河岸义编号&编号属于候选清单&所选释义符合当前用法&未收录的新义\\
    语义角色&小王把书交给小李&交给：施事小王、受事书、接受者小李&跨度与角色符合规范&论元与当前谓词对应&事件是否真实发生\\
    AMR&少年想阅读&want-01、read-01及共享施事节点&引用存在、关系与图合法&意愿与阅读的嵌套正确&实际是否阅读\\
    逻辑式&并非所有来宾已报名&$\neg\forall x(G(x)\Rightarrow R(x))$&变量作用域与类型正确&区别于所有人均未报名&事实库外的信息\\
    SQL&已报名校外来宾有几人&连接来宾与报名表后计数&列、连接、聚合合法&条件与计数对象准确&数据库未记录的状态\\
    \bottomrule
  \end{tabular}
  \caption{语义分析的五类输出与核验。向量可以参与评分，但不是这些任务的最终交付；逻辑式和SQL的求值都依赖明确给定的环境。}
  \label{tab:nlp-analysis-semantic-outputs}
\end{table*}

五个小节按输出形式由局部到整体排列。每个小节都要说明三件事：输出的合法性怎样检查；监督是完整的结构标注，还是只有答案；以及怎样从模型分数恢复出结构。可执行的表达式还须明确执行环境。

\subsection{词义消歧与词汇意义}
\label{sec:nlp-analysis-wsd}

“bank”在“deposit money at the bank”和“sit by the river bank”中用法不同。\term{词义消歧}（word sense disambiguation, WSD）先按目标词的基本形式和词性从\term{词义清单}（sense inventory）中检索候选，再选出当前用法对应的那一项。本小节设清单只有金融机构义$F$与河岸义$R$两项，这是教学编号，不是WordNet的正式义项键；真实评测须固定清单及其版本，WordNet 3.0与其他版本的候选不同，准确率不能直接比较。

词义消歧的挑战有三项。一是清单粒度细，WordNet中常用词往往有十几个义项，彼此只差细微的用法；二是标注稀疏且长尾，最常见义项占了大部分样本，少见义项只有个位数甚至零个训练样本；三是语境可能出现清单之外的新义。方法谱系沿着“利用什么证据”展开：Lesk只比较释义与语境的词重叠，不需要标注，但受词面覆盖限制；上下文表示近邻法用预训练表示比较语境，但没有训练样本的义项无从比较；GlossBERT把释义当作输入，让无样本的义项也能被判定；释义双编码检索把释义向量预先算好，降低每次联合编码的成本；个性化PageRank在词汇语义图上传播语境证据。这些方法都只在固定清单中选择身份，新义需要另设检测。

\subsubsection{Lesk与简化Lesk释义重叠}
\label{sec:nlp-analysis-lesk}

词典释义本身含有可以与语境比较的词。Lesk的原始做法在句中各词的候选释义之间寻找重叠\cite{v9a-lesk1986}；简化Lesk只比较目标词的候选释义与原文语境，不必同时为每个邻词决定词义。设规范化后的语境词集合为$C$，候选义项$s$的释义词集合为$G_s$，打分与选择为
\[
  o(s)=|C\cap G_s|,\qquad \hat s=\argmax_{s\in\{F,R\}}o(s).
\]
规范化要事先固定小写化、词形还原、停用词表以及是否删除目标词。本例删去目标词与停用词，两条释义为
\[
 G_F=\{\code{money},\code{deposit},\code{institution}\},\qquad
 G_R=\{\code{land},\code{river},\code{edge}\}.
\]
两句的语境分别为$C_1=\{\code{deposit},\code{money}\}$与$C_2=\{\code{sit},\code{river}\}$，分数依次为$(2,0)$与$(0,1)$，恢复$F$与$R$。

原始Lesk要为语境中的每个词也选义项，最大化所有释义两两重叠之和$\sum_{i<j}|G_{s_i}\cap G_{s_j}|$。若窗口内有$m$个词、每个词平均有$k$个义项，候选组合有$k^m$种，需要限制窗口或做近似搜索；简化Lesk把这个联合问题拆成只对目标词的独立选择。较长的释义有更多偶然命中的机会，可以改用Jaccard分数$|C\cap G_s|/|C\cup G_s|$，第一句的$F$此时得$2/3$而非2。

Lesk的局限是词面覆盖。语境“withdraw cash”与两条释义都没有共同词，分数同为0；这并不说明两个义项都不成立，只说明释义与语境用了不同的词。零分平局应保留歧义或按事先约定回退到最常见义项，不能把平局的打破方式当作语义证据。Lesk不需要任何标注，判断依据可以逐词追溯，适合作为标注稀少时的起点和下述模型的字面基线。

\subsubsection{上下文表示近邻消歧}
\label{sec:nlp-analysis-wsd-knn}

Lesk的困难在于同义而不同形的词无法匹配。预训练语言模型给出的上下文表示能缓解这一点：同一个词形在不同语境中得到不同的向量，用法相近的语境向量也相近。第~\ref{sec:nlp-classical-word-vectors}小节的静态词向量做不到这一点，bank在两句话中查到的是同一个向量，混合了两种用法的统计。

近邻法直接利用这种性质。设编码器在语境$x$中目标词位置的输出为$\bm q=f(x)\in\mathbb R^d$（目标词被切成多个子词时取平均）。对义项$s$，取训练集中所有标注为$s$的样本，计算它们目标位置向量的平均，作为\term{义项向量}（sense embedding）：
\begin{equation}
 \bm e_s=\frac{1}{|D_s|}\sum_{x\in D_s}f(x),
 \label{eq:nlp-analysis-sense-centroid}
\end{equation}
$D_s$是义项$s$的训练样本集合。测试时按余弦相似度取最近的义项，$\hat s=\argmax_{s}\cos(\bm q,\bm e_s)$。这里没有可训练的分类器参数，编码器也保持冻结；学习只发生在求平均这一步。

用二维向量手算（教学数值）。金融义的两个训练语境向量为$(1,0)$与$(0.8,0.2)$，平均为$\bm e_F=(0.9,0.1)$；河岸义的两个为$(0,1)$与$(0.2,0.8)$，平均为$\bm e_R=(0.1,0.9)$。测试句“sit by the river bank”中bank的向量为$\bm q=(0.2,0.7)$，于是
\[
 \cos(\bm q,\bm e_F)=\frac{0.25}{0.906\times0.728}\approx0.38,\qquad
 \cos(\bm q,\bm e_R)=\frac{0.65}{0.906\times0.728}\approx0.99,
\]
选$R$。Peters等用ELMo的上下文表示做了这种一近邻消歧，在统一的英语全词消歧评测上F1为69.0\cite{nlp-classical-peters2018}；第~\ref{sec:nlp-classical-elmo}小节介绍了ELMo的结构。

近邻法的缺口是没有训练样本的义项：式~\eqref{eq:nlp-analysis-sense-centroid}的分母为0，义项向量不存在，这些义项永远不会被选中。Peters等对这类词回退到WordNet的第一义项。Loureiro与Jorge的LMMS改为补全义项向量：沿WordNet的同义词集、上位词和词典类别关系，用邻近义项向量的平均填补缺失者，再把释义文本经同一编码器得到的向量拼接进来，使清单中每个义项都有向量\scite{nlp-analysis-loureiro2019}。没有样本也没有关系邻居的义项，最终只能依赖释义向量，这就自然引出下一种把释义作为模型输入的方法。

\subsubsection{GlossBERT语境--释义判定}
\label{sec:nlp-analysis-glossbert}

近邻法只在训练时间接用到释义。GlossBERT把释义直接放进输入，让模型判断“这条释义是否适用于这个语境”\scite{huang-etal-2019-glossbert}。对目标词的每个候选义项$s$，构造一个句对样本：
\[
 \code{[CLS] context [SEP] gloss}_s\ \code{[SEP]},
\]
并在语境中目标词两侧加引号等约定标记，指明是哪一次出现。正确义项的标签为yes，其余为no。BERT的句对编码与分类头见第~\ref{sec:nlp-classical-bert}小节；这里取\code{[CLS]}位置的表示$\bm h_s$，输出$p_s=\operatorname{sigmoid}(\bm w^{\top}\bm h_s+b)$，以所有正负候选的二元交叉熵训练。推断时取$p_s$最大的候选。

设河岸句对两个候选的logit为$-1$与$2$，yes概率约为$.269$与$.881$，选$R$。这两个概率来自两次独立的二分类，不要求相加为1。若句子是“the bank financed repairs to the river bank”，两个bank的位置标记不同，同一语境会得到两组不同的判定；位置标记正是为这种情况准备的。

与近邻法相比，GlossBERT的收益在长尾：没有训练样本的义项也有释义，模型能对它做判定，因为它学到的是“语境与释义是否匹配”这一跨义项共享的能力，而不是每个义项各自的分类器。其带目标标记的句对版本在统一评测ALL上F1为77.0，同一表中最常见义项基线为65.5。代价是计算：一个目标词有$k$个候选就要做$k$次完整的联合编码。

\subsubsection{释义双编码检索}
\label{sec:nlp-analysis-gloss-retrieval}

GlossBERT每次都要重读语境与释义。候选很多或需要频繁调用时，可以把两边分开编码。Blevins与Zettlemoyer的双编码方法用一个编码器读语境、取目标词位置的表示$\bm q=f(x,\text{目标跨度})$，另一个编码器读释义、取\code{[CLS]}位置的表示$\bm g_s=g(\text{释义}_s)$，按内积在候选义项$S(w)$上做softmax\cite{blevins-zettlemoyer-2020-moving}：
\[
 p(s\mid x)=\frac{\exp(\bm q^{\top}\bm g_s)}
                  {\sum_{t\in S(w)}\exp(\bm q^{\top}\bm g_t)}.
\]
两个编码器以交叉熵联合训练。用$\bm g_F=(1,0)$、$\bm g_R=(0,1)$与河岸语境$\bm q=(0,2)$，两候选得分为0与2，概率约为$.119$与$.881$，返回第二个向量所对应的义项$R$；金融语境$\bm q=(2,0)$给出相反的排序。

释义向量与语境无关，可以对整个清单离线算好。推断时语境只编码一次，再付出$|S(w)|$次$d$维内积。这与第~\ref{sec:nlp-analysis-wsd-knn}小节近邻法的计算方式相同，区别是“义项向量”不再由训练样本平均得到，而由释义编码得到，因此天然覆盖没有样本的义项。与GlossBERT相比，它放弃了语境与释义之间的逐词交叉注意力。Blevins与Zettlemoyer报告双编码在ALL上F1为79.0，提升主要来自少见义项：在训练中出现较少的义项上为52.6，而BERT分类基线为37.0。

\subsubsection{词汇语义图上的个性化PageRank}
\label{sec:nlp-analysis-wsd-ppr}

释义之外，义项之间的关系也携带信息：贷款义与金融机构义相连，河流义与河岸义相连。把义项作为节点、词汇语义关系作为边，再把语境中的词映射到它们的候选节点，就可以从语境出发在图上随机游走，看目标词的哪个候选最容易被走到。\term{个性化PageRank}（personalized PageRank）用非均匀的重启分布表达这种偏置\scite{agirre-soroa-2009-personalizing}。

令$\bm P$为按列归一化的转移矩阵，$\bm v$为非负且和为1的语境重启分布，迭代
\begin{equation}
 \bm r^{(t+1)}=\rho\bm P\bm r^{(t)}+(1-\rho)\bm v,
 \qquad 0<\rho<1.
 \label{eq:nlp-analysis-personalized-pagerank}
\end{equation}
每一步以概率$\rho$沿边走一步，以概率$1-\rho$跳回语境节点。没有出边的节点按约定转向$\bm v$，使每列仍和为1。从$\bm r^{(0)}=\bm v$开始，直到相邻两次的$L_1$差小于阈值。由于$\bm P$是列随机矩阵，映射$\bm r\mapsto\rho\bm P\bm r+(1-\rho)\bm v$在$L_1$范数下是系数为$\rho$的压缩映射，迭代收敛到唯一的不动点，误差每步至少缩小为原来的$\rho$倍。

构造四节点图，顺序为金融义$F$、贷款义$L$、河岸义$R$、河流义$W$，只有无向边$F$--$L$与$R$--$W$。贷款语境映射为$\bm v=(0,.9,0,.1)$，取$\rho=.5$。一次迭代得$(.45,.45,.05,.05)$，第二次为$(.225,.675,.025,.075)$，极限为$(.3,.6,1/30,1/15)$。只在目标词的候选$F,R$之间比较，选$F$。河流语境把$.9$与$.1$的位置交换，极限也相应交换，选$R$；若用均匀重启，四个节点的值都是$.25$，两种语境无法区分。

个性化PageRank不需要逐义项的标注，所需知识是WordNet这类语义关系图和词到节点的映射。它的判断质量取决于图的覆盖：词形映射错误、图中缺边或语境窗口过宽都会改变排序，而且传播只能在已有节点之间进行。表~\ref{tab:nlp-analysis-wsd-methods}把五种方法放在同一“固定清单内选择义项”的任务下比较。

\begin{table*}[htbp]
  \centering
  \begin{tabularx}{\linewidth}{@{}l XXXX@{}}
    \toprule
    方法&解决的问题&核心机制&监督与计算&新代价\\
    \midrule
    简化Lesk&无标注时的判定依据&语境与释义的词重叠&无需训练；集合运算&同义而不同形时失效\\
    上下文近邻&同形词在不同语境中的区分&义项向量取样本平均；余弦近邻&只需标注样本；编码一次&无样本义项缺向量\\
    GlossBERT&长尾与无样本义项&语境--释义句对二分类&逐候选联合编码&$k$个候选需$k$次编码\\
    释义双编码&联合编码的计算成本&语境与释义分别编码；内积检索&释义向量可预计算&失去逐词交互\\
    个性化PageRank&利用义项关系&语境重启的随机游走&无需义项标注；迭代&依赖图的覆盖\\
    \bottomrule
  \end{tabularx}
  \caption{词义消歧方法对照。比较判定结果前，须统一词义清单版本与目标位置。}
  \label{tab:nlp-analysis-wsd-methods}
\end{table*}

\subsection{谓词、论元与语义角色}
\label{sec:nlp-analysis-srl}

在“小王昨天把一本书交给小李”中，“交给”描述一个事件，其他片段说明参与者与时间。\term{谓词}（predicate）是事件或关系的表达，\term{论元}（argument）是与之对应的参与对象，\term{语义角色}（semantic role）说明对象怎样参与。\term{语义角色标注}（semantic role labeling, SRL）对句中每个谓词找出论元并标注角色。

角色体系有两种主要规范。PropBank按谓词义项定义编号角色：每个谓词义项（如“交给.01”）有自己的角色框架，ARG0通常是施事、ARG1通常是受事，但具体含义以该谓词的框架文件为准；时间、地点等附加角色用ARGM-TMP、ARGM-LOC等标签\scite{nlp-analysis-palmer2005}。FrameNet则以\term{语义框架}（semantic frame）组织：一个框架描述一类情境及其参与者，称为\term{框架元素}（frame element），能唤起该框架的词称为词汇单元\scite{nlp-analysis-baker1998}。“交给”唤起Giving框架，框架元素有给予者Donor、接受者Recipient和所给之物Theme，它们对“给”“赠送”“递交”等词共享。本节教学规范采用PropBank式标签：施事A0、受事A1、接受者A2、时间TMP。

给定分词“小王／昨天／把／一本书／交给／小李”，各词的字符跨度依次为$[0,2)$、$[2,4)$、$[4,5)$、$[5,8)$、$[8,10)$、$[10,12)$。期望输出是以“交给”$[8,10)$为谓词的四条记录：A0小王、TMP昨天、A1一本书、A2小李。评价按“谓词、论元跨度、角色”三元组计算F1；是否给定谓词位置是实验条件的一部分，必须说明。

SRL的挑战有五项：谓词识别与谓词义项（或框架）消歧；论元边界；同一谓词下核心角色唯一、论元互不重叠的全局约束；一句多个谓词时各谓词的论元集合不同；以及没有在句中出现的隐含论元。方法谱系为：成分候选分类加整数线性规划依赖句法树，在全局施加约束；BIO--CRF去掉对句法的依赖，端到端标注，但每个谓词要重新编码一次；联合跨度法一次枚举谓词与论元，同时预测所有谓词的角色；依存式只标论元的中心词，框架语义解析把谓词所唤起的框架与框架元素联合预测。这里恢复的是语言层面的参与关系；按业务事件本体组织事实由第~\ref{chap:nlp-information-extraction}章承担。

\subsubsection{成分候选分类与整数线性规划约束}
\label{sec:nlp-analysis-srl-constituents}

论元大多对应句法树中的一个成分。Gildea与Jurafsky据此把SRL写成成分分类问题：对谓词所在句子的每个成分，预测它承担的角色或NONE\scite{nlp-analysis-gildea2002}。成分数与句长同阶，其中绝大多数与谓词无关。Xue与Palmer的剪枝规则从谓词节点出发，收集它的所有兄弟节点作为候选，然后移到父节点，再收集父节点的兄弟，如此直到根；兄弟是介词短语时，还收集它的孩子\scite{nlp-analysis-xue2004}。这条规则利用了论元通常是谓词路径上某个节点的兄弟这一句法观察。

分类器的经典特征包括：短语类型（NP、PP等）；从候选到谓词的\term{句法路径}（syntactic path），例如$\mathrm{NP}\uparrow\mathrm S\downarrow\mathrm{VP}\downarrow\mathrm V$，记录在树中先向上、再向下经过的类别；候选在谓词之前还是之后；谓词的语态；候选的中心词。神经版本把候选表示与谓词表示拼接，再接softmax。训练标签由参考论元与候选跨度精确匹配得到，未匹配的候选标NONE。

逐候选独立分类会产生冲突的输出。Punyakanok等把这些候选的分类概率交给整数线性规划，在全局约束下联合选择\scite{v9a-punyakanok2008}。设候选集合为$\mathcal A$，角色集合为$\mathcal R$（含NONE），二元变量$z_{a,r}=1$表示候选$a$取角色$r$，$p(a,r)$为分类器给出的概率。问题写为
\begin{equation}
 \begin{aligned}
  \max_{z}\ &\sum_{a\in\mathcal A}\sum_{r\in\mathcal R}p(a,r)\,z_{a,r}\\
  \text{s.t.}\ &\textstyle\sum_{r\in\mathcal R}z_{a,r}=1,&&\forall a\in\mathcal A,\\
  &\textstyle\sum_{a\in\mathcal A}z_{a,r}\leq1,&&\forall r\in\mathcal R_{\mathrm{core}},\\
  &\textstyle\sum_{r\neq\mathrm{NONE}}(z_{a,r}+z_{a',r})\leq1,&&\forall\,a,a'\text{重叠},\\
  &z_{a,r}\in\{0,1\}.
 \end{aligned}
 \label{eq:nlp-analysis-srl-ilp}
\end{equation}
第一组约束让每个候选恰取一个标签；第二组要求核心角色（A0至A5）在同一谓词下至多出现一次；第三组禁止两个重叠的候选同时作论元。原文还加入了引用型角色R-AX必须有对应的AX、谓词义项只允许其框架中的角色等约束。目标函数是概率之和，约束都是线性的，可以交给通用整数规划求解器；SRL的候选数不大，求解很快。

用一组教学概率说明两类约束的作用。候选“一本书”取A1的概率为$.7$、NONE为$.3$；包含它的父成分“把一本书”取A1为$.6$、NONE为$.4$；“小王”取A0为$.8$、NONE为$.2$；“小李”取A0为$.55$、A2为$.40$、NONE为$.05$。独立取最大会得到两个A1（还彼此重叠）和两个A0。整数规划比较可行的组合：保留“一本书”的A1、舍去“把一本书”，得$.7+.4=1.1$，优于反过来的$.6+.3=.9$；“小王”取A0、“小李”改取A2，得$.8+.40=1.2$，优于“小王”取NONE、“小李”取A0的$.2+.55=.75$。最终恢复四条期望记录。

若只有不重叠约束而没有角色唯一性，问题退化为加权区间调度，可以精确地动态规划。把候选按右端点排序为$a_1,\ldots,a_m$，令$w_k=\max_{r\neq\mathrm{NONE}}p(a_k,r)-p(a_k,\mathrm{NONE})$为把$a_k$选作论元相对于NONE的收益，$\pi(k)$为右端点不超过$a_k$左端点的最后一个候选，则
\[
 \mathrm{OPT}(k)=\max\{\mathrm{OPT}(k-1),\ w_k+\mathrm{OPT}(\pi(k))\},\qquad\mathrm{OPT}(0)=0.
\]
总分再加上所有候选的NONE概率之和。核心角色唯一性把不同位置的选择耦合起来，这个递推就不再成立，需要在状态中记录已用的角色，或回到整数规划。

成分路线的上限由句法决定：若解析器漏掉了“一本书”这个成分，即使分类和推断都正确，A1也不可能恢复。Punyakanok等比较了自动句法与参考句法下的结果，句法错误是最终角色错误的重要来源，这正是下一种方法去掉句法依赖的动机。

\subsubsection{BIO--CRF语义角色标注}
\label{sec:nlp-analysis-srl-crf}

不依赖成分候选时，可以对每个谓词给词序列打BIO角色标签，把SRL变成第~\ref{sec:nlp-analysis-paradigm-tagging}小节的序列标注。标签集为O以及每个角色的B与I；合法转移的约束与第~\ref{sec:nlp-analysis-paradigm-tagging}小节式~\eqref{eq:nlp-analysis-legal-mask}相同，这里的掩码规定I-A1只能跟在B-A1或I-A1之后，其他角色同理。

Zhou与Xu的模型是这一路线的代表\scite{nlp-analysis-zhou2015}。每个词的输入拼接四部分：词嵌入；谓词的词嵌入；谓词前后窗口的词嵌入；以及一个指示当前词是否落在谓词窗口内的标记。这样同一句话对不同谓词得到不同的输入。输入经多层交替方向的LSTM得到每个位置的表示，线性层给出各标签的发射分数，顶层接线性链CRF，训练目标是参考标签序列的条件对数似然。He等的深层BiLSTM去掉了CRF，改用逐位置交叉熵训练、解码时加BIO约束，单模型在CoNLL-2005的WSJ测试集上F1为83.1\scite{nlp-intro-he2018}。

对上面的六个词，目标序列为B-A0、B-TMP、O、B-A1、O、B-A2，谓词自身记O，谓词位置由输入标记提供。设这些目标标签的发射分数均为2，其他标签为0，另外令第四词的I-A1得分为3（教学数值）。逐位置贪心会在第四词取I-A1，但它前面是O，违反约束；带掩码的Viterbi取目标序列，总分12，而把第三词改为B-A1、第四词改为I-A1的合法竞争路径总分为11。回溯后经词到字符的映射恢复四个角色跨度。

一句含多个谓词时，必须为每个谓词生成一个条件序列。“小王答应小李阅读通知”中，对“答应”，小李是被承诺的一方，“阅读通知”是承诺的内容；对“阅读”，小王是阅读者，通知是受事。一个BIO序列无法同时表达这两组关系，所以$p$个谓词要编码$p$次。此外，BIO掩码只约束相邻标签，不禁止远处再出现一个B-A0；核心角色唯一性需要扩展状态或解码后的结构选择。重叠的论元也无法用单层BIO表示。

\subsubsection{联合跨度式语义角色标注}
\label{sec:nlp-analysis-srl-span}

BIO方法每个谓词编码一次，谓词还要事先给定。He等的联合模型改为一次编码整句，同时枚举候选谓词与候选论元跨度，直接为“谓词、论元、角色”三元组打分\scite{v9a-he2018}。它属于第~\ref{sec:nlp-analysis-paradigm-span}小节的跨度范式：论元跨度$a$的表示用式~\eqref{eq:nlp-analysis-span-representation}，即两端表示、内部注意力汇聚与宽度嵌入的拼接；谓词$p$的表示取其位置的编码器输出。对非空角色$r$，
\[
 s(p,a,r)=s_p(p)+s_a(a)+s_r(p,a),\qquad
 s(p,a,\mathrm{NONE})=0,
\]
$s_p$与$s_a$是一元的前馈打分，衡量一个位置像不像谓词、一个跨度像不像论元；$s_r$是对谓词与论元表示拼接后的前馈打分，给出每个角色的分数。以NONE的0分为参照，对每个候选对做角色类别上的交叉熵。

全部配对有$O(n)$个谓词乘$O(n^2)$个跨度，再乘角色数，无法枚举。模型按一元分数剪枝：只保留得分最高的$\lambda_p n$个谓词和$\lambda_a n$个论元跨度，并限制论元的最大宽度。剪枝前要检查训练参考的覆盖，被剪掉的正例不会参与关系学习。原文在CoNLL-2012开发集上用每词0.4个谓词的预算，保留了超过99.7\%带论元的参考谓词，说明一元分数足以在这类语料中控制候选量。

对“交给”与“一本书”，设三项分数为$1,1,2$，A1总分4；与“书”配对为$1,.5,1.5$，总分3；与“把一本书”配对总分2。独立地取标签会保留三个重叠的A1。联合模型本身对每个候选对独立做softmax，并不保证不重叠；本例按A1唯一、论元不重叠的约束选择，保留4分的完整跨度$[5,8)$。这一步可以用第~\ref{sec:nlp-analysis-srl-constituents}小节的整数规划或区间调度实现。

与BIO方法相比，联合跨度法一次编码就能处理全部谓词，并能输出重叠论元；代价是$O(n^2)$的跨度候选和额外的剪枝超参数。连续跨度的假设不覆盖不连续论元，最大宽度限制可能漏掉很长的论元。

\subsubsection{依存式语义角色标注与框架语义解析}
\label{sec:nlp-analysis-srl-dependency}

前三种方法都输出论元的完整跨度。CoNLL-2009共享任务采用另一种规范：论元只用其句法中心词表示，输出是从谓词词到论元中心词的带角色边，同时还要预测谓词的义项编号\scite{nlp-analysis-hajic2009}。这种规范把边界问题移交给句法，适合以依存树库为基础的多语言资源。

实现上，编码每个词后，用双仿射（第~\ref{sec:nlp-analysis-biaffine}小节）或拼接后的前馈网络计算$s(p,j,r)$，在含NONE的角色上训练交叉熵；谓词义项用谓词位置的表示在其候选义项上分类。把例句细分为“小王／昨天／把／一／本／书／交给／小李”时，完整受事仍是“一本书”，中心词却是“书”。设“交给”到“书”的A1、A2、NONE分数为$3,0,0$，概率约为$.909,.045,.045$，恢复中心词角色边。下游若需要完整短语，要再从中心词的句法子树或另设的边界预测头扩展，扩展本身也需要评价。跨度式与中心词式的F1分母不同，不能直接比较。

\term{框架语义解析}（frame-semantic parsing）把谓词消歧与角色标注都换成FrameNet规范。Das等的SEMAFOR把它分成三步\scite{nlp-analysis-das2014}。第一步识别\term{目标词}（target），即句中唤起框架的词，如“交给”。第二步是框架识别：目标词的词汇单元限定了候选框架集合$\mathcal F(t)$，模型在其中选择，$\hat f=\argmax_{f\in\mathcal F(t)}\bm\theta^{\top}\bm\phi(f,t,x)$；“交给”的候选可能有Giving与Delivery等。第三步是论元识别：对选中框架的每个框架元素，在候选跨度中选择一个跨度或空，并施加同一框架元素不重复、跨度互不重叠的约束，用束搜索或对偶分解联合解码。对本例，Giving框架给出Donor小王、Theme一本书、Recipient小李，“昨天”标为非核心框架元素Time。

框架的收益在于跨谓词共享：“交给”“赠送”“递交”都唤起Giving，其框架元素可以直接对齐，而PropBank的A2在不同谓词下含义未必相同。代价是FrameNet的标注规模远小于PropBank，词汇单元覆盖不全，许多目标词找不到候选框架。表~\ref{tab:nlp-analysis-srl-methods}比较四条路线。

\begin{table*}[htbp]
  \centering
  \begin{tabularx}{\linewidth}{@{}l XXXX@{}}
    \toprule
    方法&解决的问题&核心机制&监督与推断&新代价\\
    \midrule
    成分分类+ILP&候选空间与全局角色约束&句法剪枝；路径特征；整数规划&角色／NONE分类；约束下联合选择&受句法错误限制\\
    BIO--CRF&去掉句法依赖&谓词条件的序列标注；合法转移掩码&参考标签序列；Viterbi&每个谓词编码一次；不能表达重叠\\
    联合跨度&谓词与论元的联合识别&跨度表示；三项分数相加；一元剪枝&谓词、跨度、角色；配对选择&$O(n^2)$候选；剪枝预算\\
    依存式／框架语义&中心词规范；跨谓词共享的角色&谓词到论元头的边；框架识别+框架元素&义项或框架分类；约束解码&需恢复完整短语；框架覆盖有限\\
    \bottomrule
  \end{tabularx}
  \caption{语义角色标注方法对照。完整跨度F1、中心词F1与框架元素F1回答不同的问题，比较时须固定谓词是否给定。}
  \label{tab:nlp-analysis-srl-methods}
\end{table*}

\subsection{句子语义图}
\label{sec:nlp-analysis-amr}

多个谓词可能谈到同一个对象。“少年想阅读”中，少年既是“想”的主体，也是所期望的“阅读”的主体。语义角色标注会为两个谓词分别输出一个“少年”，却不说明它们是同一个对象。\term{抽象意义表示}（abstract meaning representation, AMR）把整句的意义写成一张有根有向图：节点是概念，边是带标签的关系，变量标识节点的身份，同一变量可以作为多条关系的端点。共享节点的多条入边称为\term{重入}（reentrancy）\scite{nlp-analysis-banarescu2013}。

采用英语AMR的概念标签，这句话可以用PENMAN记法写成：
\begin{codeblock}{共享论元的AMR}
(w / want-01
   :ARG0 (b / boy)
   :ARG1 (r / read-01
            :ARG0 b))
\end{codeblock}
斜线后是概念，冒号后是关系，末尾的b引用已经建立的节点。图中有三个概念节点和三条关系；原文只有一次“少年”，图中也只有一个b。AMR抽象掉了时态、冠词等表层形式，也不断言阅读实际发生，\code{read-01}处在意愿的内容中。

AMR的评价指标Smatch先在两张图之间搜索一对一的变量对齐，再比较概念、属性与关系三元组的F1\scite{cai-knight-2013-smatch}。只比较变量名会把相同的图误判为不同，所以对齐是必需的。若预测只省掉了$r\xrightarrow{\mathrm{ARG0}}b$，不计TOP三元组时参考有6个三元组、预测有5个，最佳对齐下匹配5个，精确率1、召回率$5/6$，F1为$10/11$。对齐搜索是NP难的，Smatch用多次重启的爬山法近似，可能错过最佳对齐。

AMR解析的挑战有四项：概念与词不一一对应，一个词可能对应零个或多个概念；重入使输出是图而非树，第~\ref{sec:nlp-analysis-dependency}小节的树解码器不能直接使用；图上没有现成的精确解码算法；训练需要文本与图之间的对齐，而AMR语料并不标注对齐。方法谱系为：CAMR从依存树出发，用转移动作把树改写成图，依赖句法质量；JAMR先识别概念、再求最大生成连通子图，需要外部对齐器；SPRING把图线性化后用预训练的序列到序列模型生成，不需要显式对齐，但须修复生成出的非法序列。

\subsubsection{CAMR转移式AMR解析}
\label{sec:nlp-analysis-camr}

依存树与AMR图在结构上很接近：多数依存弧对应一条语义关系，主要的差别是功能词要删掉、部分方向要反转、共享论元需要加边。CAMR据此把AMR解析写成对依存树的改写，属于第~\ref{sec:nlp-analysis-paradigm-transition}小节的转移范式\scite{wang-etal-2015-transition}。状态由当前的部分图、按依存树后序排列的待处理节点队列$\sigma$、当前节点$\sigma_0$的待处理孩子边队列$\beta$组成，$\beta_0$是当前孩子。初始时部分图就是依存树，节点与边都还没有语义标签。表~\ref{tab:nlp-analysis-camr-actions}列出全部动作。

\begin{table*}[htbp]
  \centering
  \begin{tabularx}{\linewidth}{@{}l X X@{}}
    \toprule
    动作&前提&效果\\
    \midrule
    NEXT-EDGE-$l_r$&$\sigma_0$有待处理的孩子$\beta_0$&把边$\sigma_0\to\beta_0$标为关系$l_r$，转向下一个孩子\\
    SWAP-$l_r$&同上&把边反转为$\beta_0\to\sigma_0$并标为$l_r$，$\beta_0$成为$\sigma_0$的父节点\\
    REATTACH$_k$-$l_r$&同上；$k$是图中另一节点，且不产生环&删去边$\sigma_0\to\beta_0$，改接为$k\to\beta_0$并标为$l_r$\\
    REPLACE-HEAD&同上&删去$\sigma_0$，由$\beta_0$继承$\sigma_0$的全部入边与出边\\
    REENTRANCE$_k$-$l_r$&同上；$k\to\beta_0$尚不存在&增加边$k\to\beta_0$并标为$l_r$，原边保留，当前孩子不变\\
    MERGE&同上&把$\sigma_0$与$\beta_0$合并为一个节点，用于多词概念\\
    NEXT-NODE-$l_c$&$\sigma_0$没有待处理的孩子&把$\sigma_0$标为概念$l_c$，转向队列中的下一个节点\\
    DELETE-NODE&$\sigma_0$没有待处理的孩子&删去$\sigma_0$及其入边，用于不对应概念的功能词\\
    \bottomrule
  \end{tabularx}
  \caption{CAMR的转移动作。前六个动作处理当前边，后两个处理当前节点；队列为空时得到AMR图。REENTRANCE不推进当前边，需要限制重复使用以免轨迹不终止。}
  \label{tab:nlp-analysis-camr-actions}
\end{table*}

对“少年$_1$／想$_2$／阅读$_3$”，依存边为$2\to1$与$2\to3$，后序队列为$[1,3,2]$。节点1没有孩子，执行NEXT-NODE-boy；节点3同样执行NEXT-NODE-read-01。处理节点2的第一条边$2\to1$时，执行REENTRANCE$_3$-ARG0，增加$3\to1$而保留原边，当前边不变；随后NEXT-EDGE-ARG0给$2\to1$标注并推进，NEXT-EDGE-ARG1处理$2\to3$，最后NEXT-NODE-want-01标注节点2，队列清空。把节点1、2、3重命名为b、w、r，即得上面的图。

训练轨迹由oracle产生：先用对齐把参考AMR的概念对应到词，得到与依存树共享节点的“跨度图”，再比较当前部分图与参考图，确定性地选出下一个动作。原文用对数线性（最大熵）分类器在每个状态上选择动作，特征来自当前节点与孩子的词、词性、依存标签和路径。CAMR在LDC2013E117上的Smatch比JAMR高约五个百分点。它的优点是每个动作都对应一次可解释的结构改写；代价是依赖依存解析的质量，依存树中错误的结构会限制可恢复的目标，也依赖对齐来产生oracle。

\subsubsection{JAMR概念识别与最大生成连通子图}
\label{sec:nlp-analysis-amr-node-edge}

CAMR借用句法树作为图的初稿。JAMR不依赖句法，把解析分成两步：先从词跨度识别概念，再在识别出的概念之间选关系边\scite{v9a-flanigan2014}。

概念识别写成半马尔可夫标注。把句子切成若干段$(b_k,e_k)$，每段标为一个概念片段$c_k$或空，片段可以是单个概念，也可以是一小块子图，例如人名“小王”对应\code{(p / person :name (n / name :op1 "小王"))}。候选片段取自训练集上由对齐抽出的概念词典。整句的分数是各段分数之和$\sum_k\bm\theta^{\top}\bm f(b_k,e_k,c_k)$，用动态规划在所有切分上求最大，复杂度与句长、最大段长和候选片段数成正比。

关系识别在概念节点之间选边。设识别出的节点集合为$V$，每对节点$(u,v)$与每个关系标签$l$都有分数$\psi(u,l,v)=\bm\theta^{\top}\bm g(u,l,v)$，片段内部的边必须保留。目标是选出边集$E$使$\sum_{e\in E}\psi(e)$最大，且图连通、任意两点之间至多一条边。JAMR称这样的图为\term{最大生成连通子图}（maximum spanning connected subgraph, MSCG），求解与Kruskal最小生成树算法相似：
\begin{enumerate}
  \item 加入片段内部的边，以及所有分数为正的候选边（同一对节点只取最高分的标签）；
  \item 计算当前的连通分量；把剩余候选边按分数从高到低排序，依次检查，若一条边连接两个不同的分量，就加入并合并分量；
  \item 全图连通时停止。
\end{enumerate}
第一步加入正分边只会增加总分；第二步要用尽量少、尽量高分的负分边把分量连起来，这与Kruskal的贪心论证相同。原文证明了这一过程给出满足连通约束的最优解。

AMR还要求满足\term{确定性约束}（determinism constraint）：一个节点的ARG0、ARG1等核心关系出边各至多一条。这类约束把不同边的选择耦合起来，MSCG算法无法直接处理。JAMR用拉格朗日松弛：为每条确定性约束引入乘子$\lambda\geq0$，把违反约束的代价并入边分数，用调整后的分数运行MSCG；若结果仍违反约束，按次梯度增大相应乘子，再重复。若某次调整后的解满足全部确定性约束，且取正值的乘子只出现在恰好取等号的约束上，这个解就是原问题的最优解；迭代在预定次数内达不到这一条件时，只能得到近似解。

用三个节点手算（教学分数）。设$w\xrightarrow{\mathrm{ARG0}}b$为$2.2$，$w\xrightarrow{\mathrm{ARG1}}r$为$-0.2$，$r\xrightarrow{\mathrm{ARG0}}b$为$-0.4$，其余边更低。第一步只加入$w\to b$，分量为$\{w,b\}$与$\{r\}$。第二步按分数检查剩余边，$w\to r$连接两个分量，加入；此时已连通，停止。结果有$w\to b$与$w\to r$，却没有重入边$r\to b$：MSCG只加“连通所需”的负分边，模型若没有给重入边正分，它就会丢失。反过来，若$r\to b$得分为$0.3$，第一步就会把它加入，重入得以保留。

JAMR的训练需要概念与词的对齐，而AMR语料不提供对齐。作者设计了基于规则的对齐器，按精确匹配、命名实体、日期、词形变化等规则依次贪心地把概念对齐到词跨度。对齐错误会同时损害概念词典与两个阶段的训练数据；概念阶段漏掉的节点，关系阶段也无从连接。

\subsubsection{SPRING图线性化与序列到序列解析}
\label{sec:nlp-analysis-spring}

CAMR与JAMR都需要对齐来产生训练信号，又都有多个阶段，前一阶段的错误传给后一阶段。SPRING去掉对齐与分阶段的接口：把AMR图线性化为符号序列，用预训练的编码器--解码器BART直接从句子生成这个序列\scite{nlp-analysis-spring2021}。它属于第~\ref{sec:nlp-analysis-paradigm-generation}小节的约束生成范式。

线性化必须保留节点身份与重入。SPRING按深度优先遍历图，用专用的指针符号代替变量名：
\begin{codeblock}{带节点指针的图序列}
( <R0> want-01
  :ARG0 ( <R1> boy )
  :ARG1 ( <R2> read-01 :ARG0 <R1> ) )
\end{codeblock}
指针第一次与概念一起出现时建立节点，之后单独出现的\code{<R1>}引用同一节点。这个序列与图可以互相恢复；若把末项换成\code{(<R3> boy)}，就变成了另一个少年。指针符号是新加入词表的特殊词元，避免普通字母既可能是常量又可能是变量的混淆。训练目标是参考序列每一步的交叉熵，括号、关系、指针与结束符都计入。

解码完成后，用栈恢复图：读到左括号压栈，建立指针到节点的映射；读到关系时记录源节点，再解析目标子图或指针引用；读到右括号退栈。全部读完后解析所有引用，检查每个指针恰有一个概念定义、图是否连通。生成可能出错：括号不配对、引用了从未定义的指针、出现非法延续。SPRING用轻量的规则修复这些情况，例如补齐括号、删去悬空引用。束搜索比较完整序列的概率；在“正确引用\code{<R1>}”与“新建节点”之间若条件概率为$.8$与$.2$，这一步就偏向共享节点，但最终排名仍取决于整条序列。

SPRING在英语AMR 2.0上的Smatch比当时最强的比较系统高3.6点。同一个模型把输入与输出对调，就成为AMR到文本的生成器，这是名称中“对称”的含义。表~\ref{tab:nlp-analysis-amr-methods}比较三条路线。

\begin{table*}[htbp]
  \centering
  \begin{tabularx}{\linewidth}{@{}l XXXX@{}}
    \toprule
    方法&解决的问题&核心机制&监督与推断&新代价\\
    \midrule
    CAMR&从已有句法结构构图&依存树后序队列上的改写动作&对齐产生oracle；逐状态分类&依赖依存解析与对齐\\
    JAMR&图上的全局选边&半马尔可夫概念识别；MSCG加拉格朗日松弛&两阶段分别训练；精确或近精确解码&需规则对齐器；阶段间错误传递\\
    SPRING&去掉对齐与多阶段接口&深度优先线性化；指针表示重入；BART生成&序列交叉熵；束搜索后恢复与修复&线性化顺序偏差；需修复非法序列\\
    \bottomrule
  \end{tabularx}
  \caption{AMR解析方法对照。三者的Smatch来自不同版本的AMR语料，不宜按绝对值直接排列。}
  \label{tab:nlp-analysis-amr-methods}
\end{table*}

\subsection{语义解析与形式表达}
\label{sec:nlp-analysis-formal}

如果目标是计算一个数量或执行一次查询，概念图还不够，需要类型、变量作用域和求值规则都明确的形式语言。\term{语义解析}（semantic parsing）在这里指把文本映射为指定形式语言中的表达式。对输出要做三层不同的判断：结构是否合乎形式语言的语法；能否在给定环境中求值；求值的结果是否回答了原问题。前一层成立不蕴含后一层。

本小节固定一个封闭、完整的有限事实环境。实体域$E=\{a,b,c\}$表示三位来宾，$G=E$为来宾集合，$R=\{a,b\}$为已报名者，关系$\mathrm{help}$的外延为$\{(a,b),(c,b)\}$；未列出的事实一律为假。问题是“有几位来宾已报名”。现实资料若不完整，缺记录只能得出未知，不能直接套用这种封闭世界解释。

语义解析的挑战有四项：输出必须是形式语言中合法的表达式；训练中没见过的组合（如新的谓词与量词搭配）要能正确组合，即组合泛化；监督可能是标注好的逻辑式，也可能只有执行答案；只有答案时，存在碰巧得到正确答案的错误表达式。方法谱系为：类型化文法限定可以构造的表达式；CCG让词项同时携带句法范畴与语义函数，组合过程可追踪；Seq2Tree与粗到细解码用神经网络生成树形逻辑式，减少括号与变量错误；答案监督下的边缘似然与策略梯度去掉对逻辑式标注的依赖。

\subsubsection{类型化文法约束的逻辑式构造}
\label{sec:nlp-analysis-typed-logic}

令$e,t,\mathbb N$分别表示实体、真值和自然数类型。$G,R:e\to t$是一元谓词，$\mathrm{help}:e\to e\to t$是二元关系，$\mathrm{count}:(e\to t)\to\mathbb N$对一元谓词计数；否定只作用于真值，量词绑定实体变量。构造器从类型为$\mathbb N$的槽开始，展开为$\mathrm{count}(P)$，再把类型为$e\to t$的槽$P$展开为$\lambda x.G(x)\land R(x)$。每一步只允许填入类型匹配的节点，变量必须已被绑定。

“有几位来宾已报名”的构造依次建立COUNT、LAMBDA、AND、G、R以及两次变量引用，得到
\[
 z_1=\mathrm{count}(\lambda x.G(x)\land R(x)).
\]
在三个实体上求值，谓词的真值向量为$(1,1,0)$，计数为2。若有“文本--逻辑式”的标注，可以对参考构造动作做交叉熵训练，束搜索累加动作的对数概率，并屏蔽类型不匹配的展开：COUNT之后不能填实体常量a，因为a不是一元谓词。

类型检查只能排除一部分错误。$\mathrm{count}(\lambda x.G(x))$类型正确，却遗漏了报名条件，得到3。否定的作用域同样改变意义：“并非所有来宾都已报名”是$\neg\forall x(G(x)\Rightarrow R(x))$，在当前环境中为真；“所有来宾均未报名”是$\forall x(G(x)\Rightarrow\neg R(x))$，为假；两者都通过类型检查。量词顺序也如此：若帮助关系为$\{(a,b),(b,c),(c,a)\}$，“每位来宾都帮助某位来宾”$\forall x\exists y\,\mathrm{help}(x,y)$为真，“有一位来宾受到所有人帮助”$\exists y\forall x\,\mathrm{help}(x,y)$为假。文法要能表达这些区别，模型才有可能在监督下学会选择。

\subsubsection{CCG与类型化\texorpdfstring{$\lambda$}{lambda}表达式}
\label{sec:nlp-analysis-ccg}

类型化文法规定了输出空间，却没有说明逻辑式的各部分来自句中哪些词。\term{组合范畴语法}（combinatory categorial grammar, CCG）给每个词项同时分配句法范畴与语义函数：$A/B$表示向右接收一个B后成为A，$A\backslash B$表示向左接收一个B\scite{nlp-analysis-ccg2005}。函数应用在组合范畴的同时组合语义，最终关系的每一部分都可以追溯到某个词项。

对“小王帮助小李”，两个名字为$\mathrm{NP}:a$与$\mathrm{NP}:b$，“帮助”为
\[
 (\mathrm S\backslash\mathrm{NP})/\mathrm{NP}
 :\lambda y.\lambda x.\mathrm{help}(x,y).
\]
先向右应用b，得到$\mathrm S\backslash\mathrm{NP}:\lambda x.\mathrm{help}(x,b)$；再向左应用a，$\beta$归约得到$\mathrm{help}(a,b)$，在给定外延中为真。若词项把两个参数的顺序写反，范畴照样能组合，却得到为假的$\mathrm{help}(b,a)$；句法组合成功不能代替词项语义的核验。

函数组合规则可以推迟参数的接入：
\[
 A/B:f,\quad B/C:g
 \quad\Longrightarrow\quad A/C:\lambda z.f(g(z)).
\]
考虑问句“哪些来宾帮助小李”。记$\mathrm{VP}=\mathrm S\backslash\mathrm{NP}$，$\mathrm{S_q}$表示答案为实体集合的问句。“哪些来宾”为$\mathrm{S_q}/\mathrm{VP}:f=\lambda P.\lambda x.G(x)\land P(x)$，“帮助”为$\mathrm{VP}/\mathrm{NP}:g=\lambda y.\lambda x.\mathrm{help}(x,y)$。按组合规则先得到$\mathrm{S_q}/\mathrm{NP}:\lambda y.\lambda x.G(x)\land\mathrm{help}(x,y)$，再对小李的常量b做函数应用，得到$\lambda x.G(x)\land\mathrm{help}(x,b)$，答案集合为$\{a,c\}$。另一条推导先让“帮助”应用b再交给$f$，语义相同。同一逻辑式有多条推导，训练时要把它们合并。

Zettlemoyer与Collins从句子与逻辑式的配对中学习CCG\scite{v9a-zettlemoyer2005}。词项由模板从训练对中生成候选，概率CCG对逻辑式$z$与推导$T$联合打分，$p(z,T\mid x)\propto\exp(\bm w^{\top}\bm f(x,z,T))$；训练只给逻辑式时，对推导求和得到$p(z\mid x)=\sum_Tp(z,T\mid x)$，最大化参考逻辑式的条件对数似然。推断时在跨度图表上保留范畴与语义，束裁剪控制候选数。CCG的优点是组合过程透明，能处理词序变化带来的长距离依赖；代价是词典需要模板或人工设计，缺少词项时整句无分析。

\subsubsection{Seq2Tree与粗到细解码}
\label{sec:nlp-analysis-logic-generation}

CCG依赖人工设计的词项模板。神经语义解析把逻辑式当作目标序列，由编码器--解码器直接生成，例如把$z_1$写成\code{count(and(guest,registered))}。扁平的序列到序列模型容易犯两类错误：括号不配对，以及把变量或常量填到错误的位置。Dong与Lapata的Seq2Tree用层次解码器应对前一类\scite{nlp-analysis-dong2016}。

Seq2Tree按层生成。解码器先生成顶层序列，遇到需要展开的子树时输出一个非终结符标记$\langle n\rangle$，例如顶层为\code{count(}$\langle n\rangle$\code{)}。顶层结束后，对每个$\langle n\rangle$启动一次子序列解码，以该标记处的隐状态作为子解码器的初始状态，并在每一步输入中拼接父节点的隐状态，使子树记住自己处在什么位置。子序列\code{and(guest,registered)}中若还有$\langle n\rangle$，再递归展开。每一层的序列都以结束符终止，括号由层次结构隐含给出，不会不配对。训练目标是所有层序列的交叉熵之和，编码器对输入的注意力在每一层都可用。

粗到细解码应对后一类错误\scite{nlp-analysis-dong2018}。它先生成一个\term{意义骨架}（meaning sketch），即去掉变量名、常量与参数的逻辑式框架，再以输入和骨架为条件生成完整的逻辑式：
\[
 p(z\mid x)=p(a\mid x)\,p(z\mid x,a),
\]
$a$是骨架。对本例，骨架为\code{count(and(@,@))}，其中@是待填的谓词槽；细解码器读取骨架的编码，在每个@处从“guest”“registered”等谓词中选择，在骨架的固定符号处直接复制。两步分别以参考骨架与参考逻辑式的交叉熵训练。骨架只表达结构，模型可以先在较小的空间中确定“计数一个合取”，不被具体谓词的选择干扰；细节生成又能利用已确定的结构。

两种方法都可以与文法约束结合：解码器只在类型允许的位置生成相应的符号，括号与变量作用域由第~\ref{sec:nlp-inference-grammar}小节的约束机制维护。它们的共同前提是有“文本--逻辑式”的标注，而这类标注需要专家编写，成本很高。

\subsubsection{答案监督下的边缘似然与策略梯度}
\label{sec:nlp-analysis-answer-supervision}

标注逻辑式昂贵，标注答案便宜得多：问“有几位来宾已报名”，只需写下“2”。Berant等在Freebase问答上系统地研究了这种只用问题--答案对训练语义解析器的设置\scite{nlp-analysis-berant2013}。逻辑式$z$成了隐变量，监督只告诉我们执行结果$\operatorname{eval}(z,D)$应当等于答案$y$，$D$是固定的事实环境。

第一种目标是\term{边缘似然}（marginal likelihood）：把所有执行出正确答案的逻辑式的概率加起来，
\begin{equation}
 \mathcal J_{\mathrm{ML}}=\log\sum_{z\in\mathcal Z(y)}p_\theta(z\mid x),\qquad
 \mathcal Z(y)=\{z:\operatorname{eval}(z,D)=y\}.
 \label{eq:nlp-analysis-marginal-likelihood}
\end{equation}
这个和无法枚举，实际只在束搜索找到的候选中计算。它的梯度为
\[
 \nabla_\theta\mathcal J_{\mathrm{ML}}=\sum_{z\in\mathcal Z(y)}q(z)\,\nabla_\theta\log p_\theta(z\mid x),\qquad
 q(z)=\frac{p_\theta(z\mid x)}{\sum_{z'\in\mathcal Z(y)}p_\theta(z'\mid x)},
\]
即对每个正确候选的对数概率梯度按它在正确候选中的相对概率加权。

第二种目标把语义解析看作一步决策：采样一个逻辑式，执行正确得奖励1，否则得0，最大化期望奖励$\mathcal J_{\mathrm{RL}}=\mathbb E_{z\sim p_\theta}[R(z)]$。这里用到策略梯度的基本结论：期望奖励对参数的梯度等于奖励加权的对数概率梯度的期望，$\nabla_\theta\mathbb E_{z\sim p_\theta}[R(z)]=\mathbb E_{z\sim p_\theta}[R(z)\nabla_\theta\log p_\theta(z)]$，用样本平均估计时无偏但方差大，常减去一个基线来降低方差。这一结论及其推导可参考\BookChapterRef{chap:rl-learning}{卷七的强化学习方法章}。用这一估计更新参数即REINFORCE\scite{rl1-williams1992-reinforce}。代入0/1奖励，
\[
 \nabla_\theta\mathcal J_{\mathrm{RL}}=\sum_{z\in\mathcal Z(y)}p_\theta(z\mid x)\,\nabla_\theta\log p_\theta(z\mid x)
 =p_\theta(\mathcal Z(y)\mid x)\cdot\nabla_\theta\mathcal J_{\mathrm{ML}}.
\]
两种梯度的方向相同，只差一个标量：策略梯度多乘了正确候选的总概率$p_\theta(\mathcal Z(y)\mid x)$。训练初期这个总概率很小，策略梯度几乎为零，学习很慢；边缘似然先做归一化，不受这一影响。Guu等据此分析了两者的关系，并提出在两者之间插值的权重$q(z)^\beta$与带随机探索的束搜索，缓解两种目标都容易锁定在早期找到的伪解上的问题\scite{nlp-analysis-guu2017}。

沿用报名计数的例子（教学数值）。设束中有三个候选：正确的计数式$z_1$概率$.4$；常量表达式\code{2}概率$.3$；遗漏报名条件的$\mathrm{count}(\lambda x.G(x))$概率$.3$，执行得3。观察答案为2时，$\mathcal Z(y)=\{z_1,\code{2}\}$，边缘似然损失为$-\log .7\approx.357$，梯度权重为$q(z_1)=.4/.7\approx.571$、$q(\code{2})\approx.429$；策略梯度的权重是未归一化的$.4$与$.3$。两种目标都在提高常量\code{2}的概率，因为它碰巧给出了正确答案。这类执行正确而语义错误的表达式称为\term{伪解}（spurious program）。

伪解要靠额外的信息排除。把报名集合改为$\{a\}$，参考答案变为1，正确的计数式仍然成立，常量\code{2}则失败；同一逻辑式在多个环境中都正确，才能说明它不是碰巧对。多环境的答案、对逻辑式形状的先验，以及要求表达式使用问题中出现的词对应的谓词，都能减少伪解。单次答案相等不构成语义正确的证明。

\subsection{文本到SQL}
\label{sec:nlp-analysis-text-to-sql}

\term{文本到SQL}（text-to-SQL）是语义解析在关系数据库上的实例：输入自然语言问题与数据库模式，输出一条SQL查询。与上一小节的逻辑式相比，SQL的执行环境是真实的数据库，形式语言有固定的方言，问题中的词还要先对应到具体的表和列。

本小节使用表~\ref{tab:nlp-analysis-sql-db}中的两张教学表。people的主键是pid；registrations的主键是$(\code{event},\code{pid})$，其中pid是引用people的外键，OTHER表示另一个活动。问题是“活动ML-0418中已确认报名的校外来宾有几人”，期望答案为1。

\begin{table}[htbp]
  \centering
  \small
  \begin{tabular}{@{}lll@{}}
    \toprule
    \multicolumn{3}{@{}l}{\code{people}}\\
    \code{pid}&\code{name}&\code{affiliation}\\
    \midrule
    1&小王&\code{external}\\
    2&小李&\code{internal}\\
    3&小赵&\code{external}\\
    \bottomrule
  \end{tabular}\qquad
  \begin{tabular}{@{}lll@{}}
    \toprule
    \multicolumn{3}{@{}l}{\code{registrations}}\\
    \code{event}&\code{pid}&\code{status}\\
    \midrule
    \code{ML-0418}&1&\code{confirmed}\\
    \code{ML-0418}&3&\code{pending}\\
    \code{OTHER}&3&\code{confirmed}\\
    \bottomrule
  \end{tabular}
  \caption{文本到SQL的教学数据库。两表通过pid连接；同一人在同一活动至多一条报名记录，状态均非空。}
  \label{tab:nlp-analysis-sql-db}
\end{table}

文本到SQL的挑战有五项。一是\term{模式链接}（schema linking）：确定问题中的“校外”“已确认”对应哪张表的哪一列、哪个取值。二是跨数据库泛化：测试时的数据库在训练中从未出现，模型不能记住特定的列名。三是多表连接与嵌套查询。四是值绑定：条件中的常量要从问题中恢复，并与库中的实际取值一致。五是两层正确性：查询合法可执行，与执行结果回答了原问题，是两件不同的事。

评测基准体现了这些挑战的演进。WikiSQL只含单表查询，按表划分训练与测试\cite{nlp-analysis-zhong2017}；Spider按数据库划分，测试库在训练中从未出现，查询含连接、嵌套与集合运算\cite{yu-etal-2018-spider}；BIRD使用大规模、取值较脏的真实数据库，并提供外部知识说明\cite{nlp-analysis-li2023bird}；Spider 2.0进一步纳入企业数据仓库中的复杂模式、方言与多步工作流\cite{nlp-analysis-spider2025}；SParC与CoSQL考查多轮上下文中的查询与澄清\scite{yu-etal-2019-sparc,yu-etal-2019-cosql}。常用指标有两种：精确匹配把预测与参考SQL拆成子句集合比较，不要求值正确；执行准确率比较两者在数据库上的执行结果，但不同的查询可能碰巧返回相同结果。Zhong等的测试套件方法在多个自动构造的数据库变体上比较执行结果，减少这种碰巧一致\scite{nlp-analysis-zhong2020testsuite}。

多轮基准还要求承接上下文。上一轮问“ML-0418已确认报名的校外来宾有几人”，下一轮“那校内呢”应当保留活动与confirmed条件，只把affiliation改为internal，结果为0；若上一轮问的是“活动何时举行”，同一句话就没有足够的条件恢复计数查询，应当请求澄清。省略的恢复方式见第~\ref{sec:nlp-analysis-utterance-rewriting}小节。

方法谱系为：SQLNet用固定骨架和槽填充处理单表查询；RAT-SQL用关系感知编码解决模式链接与跨库问题，并以语法树解码保证合法；PICARD给预训练生成模型加增量解析约束；执行引导用试运行剔除出错的候选；大语言模型以模式提示、任务分解、自一致性与执行反馈完成任务。第~\ref{chap:nlp-question-answering-dialogue}章的表格问答回指本小节，跨步骤的执行与恢复见第~\ref{chap:nlp-solving-and-execution}章。

\subsubsection{SQLNet查询骨架与槽填充}
\label{sec:nlp-analysis-sqlnet}

WikiSQL上早期的Seq2SQL把查询当作序列生成，遇到一个问题：WHERE中AND连接的条件交换顺序不改变结果，但逐词交叉熵会惩罚与参考顺序不同的正确查询。Seq2SQL为此用强化学习以执行结果为奖励训练条件部分\scite{nlp-analysis-zhong2017}。SQLNet换了一种做法：既然条件是集合，就直接预测集合，不再生成序列\scite{nlp-analysis-sqlnet2017}。

SQLNet固定单表查询的骨架：
\begin{codeblock}[language=SQL]{SQLNet的查询骨架}
SELECT $AGG $COL FROM $TABLE
WHERE $COL $OP $VALUE (AND $COL $OP $VALUE)*
\end{codeblock}
其中聚合\$AGG取自$\{\text{无},\mathrm{MAX},\mathrm{MIN},\mathrm{COUNT},\mathrm{SUM},\mathrm{AVG}\}$，运算\$OP取自$\{=,<,>\}$。每个槽由一个专门的小网络预测，共享问题与列名的编码。核心机制是\term{列注意力}（column attention）：对每一列$c$，用列名的编码$\bm c$对问题各位置的编码$\bm h_i$做注意力，得到“与这一列相关的问题表示”：
\[
 \alpha_{c,i}=
 \frac{\exp(\bm c^{\top}\bm W\bm h_i)}
      {\sum_j\exp(\bm c^{\top}\bm W\bm h_j)},\qquad
 \bm q_c=\sum_i\alpha_{c,i}\bm h_i.
\]
event列会关注“ML-0418”，status列会关注“已确认”。WHERE的列用多标签方式预测：每列独立给出入选概率$\operatorname{sigmoid}(\bm u_1^{\top}\bm c+\bm u_2^{\top}\bm q_c)$，另有一个分类器预测条件个数$K$，取概率最高的$K$列。这就是“序列到集合”的预测，与条件顺序无关。每个选中的列再预测运算符，值则用指针网络从问题中复制一段词。SELECT的列在全部列上做softmax，聚合由分类器预测。

先在registrations表上问“ML-0418有几条已确认报名记录”。设条件数预测为2，event、status、pid的入选概率为$.9,.8,.1$，取前两列；运算符均为等号，值从问题复制出ML-0418，并按给定的状态词典把“已确认”规范化为confirmed（这是本例补充的元信息）。SELECT列为pid、聚合为COUNT，得到
\begin{codeblock}[language=SQL]{骨架恢复的单表查询}
SELECT COUNT(pid) FROM registrations
WHERE event = 'ML-0418' AND status = 'confirmed';
\end{codeblock}
三行中只有第一行满足两个条件，结果为1。交换两个条件的顺序结果不变，集合预测也不会惩罚这种交换。

SQLNet在WikiSQL上的测试执行准确率为68.0，Seq2SQL为59.4。它的局限也来自骨架：原问题需要检查“校外来宾”，要访问people表并做连接，而单表骨架中没有这样的槽；嵌套与分组同样无法表达。增加训练数据不能弥补骨架没有覆盖的结构。

\subsubsection{RAT-SQL模式关联与语法树解码}
\label{sec:nlp-analysis-ratsql}

Spider要求跨数据库的多表查询。模型面对从未见过的模式，必须在推断时自己确定问题中的词对应哪张表、哪一列，以及表之间怎样连接。RAT-SQL把这些已知或可计算的关系显式地注入编码器\scite{v9a-wang2020ratsql}。

RAT-SQL把问题词、列和表共同看作一张图的节点，图~\ref{fig:nlp-analysis-ratsql-graph}画出了本例的部分节点与关系。节点之间的关系分三类。模式内部的关系来自数据库定义：两列是否同表、某列是否为某表的主键、外键指向哪一列、列属于哪张表。问题与模式之间的关系由字符串匹配计算：问题中的$n$元组与列名或表名完全匹配、部分匹配或不匹配；问题中的词是否出现在某列的取值中（值匹配）。问题词之间的关系是相对距离，截断到$\pm2$。每种关系类型都有一个可学习的嵌入。

\begin{figure*}[htbp]
  \centering
  \begin{tikzpicture}[
  x=1mm,y=1mm,figure picture,
  qword/.style={figure node,minimum height=6.5mm,minimum width=17mm,
    draw=FigureInput,fill=FigureInputFill},
  col/.style={figure node,minimum height=6.5mm,minimum width=34mm,
    font=\figurefont\scriptsize},
  tab/.style={figure node,minimum height=8mm,minimum width=24mm,
    draw=FigureOutput,fill=FigureOutputFill},
  link/.style={draw=FigureOutput,line width=1pt},
  fk/.style={draw=FigureStroke,line width=.85pt,dash pattern=on 2.5pt off 1.5pt,
    {Stealth[length=1.6mm,width=1.1mm]}-},
  own/.style={figure link},
  lab/.style={font=\figurefont\scriptsize,text=FigureMuted,fill=white,inner sep=1pt}
]
  % 问题词
  \node[figure heading] at (0,44) {问题词};
  \node[qword] (q1) at (0,34) {ML-0418};
  \node[qword] (q2) at (0,23) {已确认};
  \node[qword] (q3) at (0,12) {校外来宾};
  \node[qword] (q4) at (0,1) {几人};

  % 列
  \node[figure heading] at (52,44) {列};
  \node[col] (c1) at (52,37) {registrations.event};
  \node[col] (c2) at (52,29) {registrations.status};
  \node[col] (c3) at (52,21) {registrations.pid};
  \node[col] (c4) at (52,10) {people.pid};
  \node[col] (c5) at (52,2) {people.affiliation};

  % 表
  \node[figure heading] at (100,44) {表};
  \node[tab] (t1) at (100,29) {registrations};
  \node[tab] (t2) at (100,6) {people};

  % 问题—模式链接
  \draw[link] (q1.east) -- node[lab,pos=.45] {值匹配} (c1.west);
  \draw[link] (q2.east) -- node[lab,pos=.45] {值匹配} (c2.west);
  \draw[link] (q3.east) -- node[lab,pos=.55] {值匹配} (c5.west);

  % 模式内部关系
  \draw[own] (c1.east) -- node[lab,pos=.45] {属于} (t1.west);
  \foreach \c in {c2,c3}{\draw[own] (\c.east) -- (t1.west);}
  \draw[own] (c5.east) -- node[lab,pos=.45] {属于} (t2.west);
  \draw[own] (c4.east) -- (t2.west);
  \draw[fk] (c4.north east) .. controls (78,13) and (78,18) .. node[lab,right,pos=.5] {外键} (c3.south east);
\end{tikzpicture}

  \caption{RAT-SQL编码时使用的部分关系。左列是问题词，中列是数据库的列，右列是表；粗线为问题与列之间的值匹配，细线为列属于表的关系，虚线箭头为外键。图中只画出与本例相关的关系；实际上每一对节点都有一种关系类型（包括“无匹配”），并在全部节点之间做自注意力。“校外来宾”与external的对应来自本例补充的业务词典。}
  \label{fig:nlp-analysis-ratsql-graph}
\end{figure*}

关系通过\term{关系感知自注意力}（relation-aware self-attention）进入编码。设输入共有$N$个节点（问题词、列、表），初始表示$\bm x_i\in\mathbb R^{d_x}$由BiLSTM分别编码问题和各列名、表名得到（或由BERT联合编码）。$\bm r^K_{i,j},\bm r^V_{i,j}\in\mathbb R^{d_z/H}$是节点$i$与$j$之间关系类型的嵌入。第$h$个头的计算为
\begin{equation}
 \begin{aligned}
  e^{(h)}_{i,j}&=\frac{\bm x_i\bm W^{(h)}_Q\bigl(\bm x_j\bm W^{(h)}_K+\bm r^K_{i,j}\bigr)^{\top}}{\sqrt{d_z/H}},\qquad
  \alpha^{(h)}_{i,j}=\operatorname{softmax}_j\bigl(e^{(h)}_{i,j}\bigr),\\
  \bm z^{(h)}_i&=\sum_{j=1}^{N}\alpha^{(h)}_{i,j}\bigl(\bm x_j\bm W^{(h)}_V+\bm r^V_{i,j}\bigr),
 \end{aligned}
 \label{eq:nlp-analysis-rat-attention}
\end{equation}
其中$\bm W^{(h)}_Q,\bm W^{(h)}_K,\bm W^{(h)}_V\in\mathbb R^{d_x\times d_z/H}$，$\bm x_i$按行向量书写。各头输出拼接为$\bm z_i\in\mathbb R^{d_z}$，再接残差、层归一化与前馈网络：
\[
 \tilde{\bm y}_i=\operatorname{LayerNorm}(\bm x_i+\bm z_i),\qquad
 \bm y_i=\operatorname{LayerNorm}\bigl(\tilde{\bm y}_i+\mathrm{FC}(\operatorname{ReLU}(\mathrm{FC}(\tilde{\bm y}_i)))\bigr).
\]
这样的层堆叠多层（原文为8层）。与普通自注意力相比，唯一的改动是在键和值上加了关系嵌入：两个节点即使内容表示相近，只要关系类型不同，注意力分数和传递的信息就不同。所有节点之间仍然两两交互，关系是偏置，不是只允许沿外键边传递的硬约束。

手算一个分数（教学数值）。设问题词“校外来宾”对affiliation与status两列的内容点积都是0，affiliation与它之间有值匹配关系，使$\bm x_i\bm W_Q\bm r^{K\top}_{i,j}$一项为2，status之间为无匹配，该项为0。两列的注意力分数为2与0，归一化后为$.881$与$.119$：关系让“校外来宾”的表示主要从affiliation列获取信息。原始RAT-SQL的值匹配只看问题与库中取值的字面重合，中文“校外”到external的对应要由业务词典或翻译提供。

解码器按Yin与Neubig的语法树解码方式，以深度优先顺序生成SQL的抽象语法树\scite{nlp-analysis-yin2017}。每一步执行三类动作之一：ApplyRule$[\,r\,]$按SQL文法的一条产生式展开当前非终结符；SelectColumn$[\,c\,]$在当前列槽中选择一列；SelectTable$[\,t\,]$在当前表槽中选择一张表。查询的概率是各步动作概率之积，训练最小化参考动作序列的负对数似然。解码器是LSTM，每步输入拼接上一步动作的嵌入、对编码器输出的注意力上下文、父节点的状态与动作嵌入，以及当前节点类型的嵌入。ApplyRule在可用产生式上做softmax；SelectColumn先让解码器状态对编码器输出的全部节点（问题词、列与表）做注意力，再经一个“节点到列”的对齐矩阵把这一注意力转成列上的分布，SelectTable同理；对齐矩阵本身也由关系感知注意力算出。指针指向的是模式中唯一的列身份，而不是列名字符串，同名的两列不会混淆。

本例的查询按一种简化文法展开为如下动作序列：ApplyRule[查询$\to$选择 来源 条件]；ApplyRule[选择$\to$COUNT(DISTINCT 列)]，SelectColumn[people.pid]；ApplyRule[来源$\to$表 表]，SelectTable[people]，SelectTable[registrations]；ApplyRule[条件$\to$AND(条件,条件,条件)]；三次ApplyRule[条件$\to$列 = 值]，分别接SelectColumn[registrations.event]、SelectColumn[registrations.status]、SelectColumn[people.affiliation]。连接条件\code{r.pid = p.pid}可以在序列化时由外键关系补出，三个常量值由值链接填入。确定性的序列化得到
\begin{codeblock}[language=SQL]{连接与聚合的语义恢复}
SELECT COUNT(DISTINCT p.pid)
FROM people AS p
JOIN registrations AS r ON r.pid = p.pid
WHERE r.event = 'ML-0418'
  AND r.status = 'confirmed'
  AND p.affiliation = 'external';
\end{codeblock}
连接得到3条记录，活动条件留下2条，状态条件留下1条，该人的affiliation为external，计数为1。语法树解码保证每条输出都符合SQL文法，但不保证语义正确：遗漏活动条件时，小赵在OTHER中的确认记录会被计入，结果变成2，查询依然合法。

RAT-SQL的消融直接支持模式链接的作用：在Spider开发集上，去掉问题与模式之间的匹配关系后，精确匹配从55.1降到40.4。它的代价是需要大量“问题--SQL”的标注来训练编码器与解码器，SQL文法与值链接规则也要针对目标方言设计。

\subsubsection{PICARD增量解析约束}
\label{sec:nlp-analysis-picard}

RAT-SQL的合法性来自专门设计的语法树解码器。另一条路线直接微调T5这样的预训练模型，把模式与问题拼成文本输入，逐词元生成SQL；这样能利用预训练的语言知识，但输出可能不合法，非法候选还会占据束中的位置。PICARD在解码时加入增量解析，拒绝不可能完成为合法查询的扩展\scite{scholak-etal-2021-picard}。其机制与第~\ref{sec:nlp-inference-grammar}小节的文法约束解码相同：在每一步用解析器检查候选前缀，把不能继续的候选分数设为$-\infty$；本小节只说明它针对SQL做了哪些检查。

PICARD有三档检查。词法档只检查关键字、列名与表名能否被识别；语法档检查前缀能否按SQL文法继续；带守卫的档还检查作用域，例如引用的列必须属于FROM中出现的某张表，别名必须先定义。为了控制开销，它不屏蔽整个词表，而是在每步只检查模型给出的前$k$个候选词元。

设前缀为\code{SELECT COUNT(}，模型提出\code{FROM}、\code{*}、\code{pid}，概率为$.50,.30,.20$。当前解析状态需要聚合函数的参数，\code{FROM}不能使前缀继续合法，被拒绝；保留另两项的原有分数，束宽为1时选择\code{*}。词元可能只是SQL单词的一部分，解析器要区分“尚未读完”与“已经不可能完成”。带守卫的检查会拒绝\code{SELECT p.status FROM people AS p}：status虽然存在，却不属于p所指的people表。但\code{SELECT COUNT(*) FROM people}完全合法，它回答的是“共有几人”，不是原问题。

PICARD的同模型对照隔离了约束的作用：在Spider开发集上，T5-3B产生执行错误的查询比例由12\%降到2\%。代价是每步的解析检查，以及正确词元不在前$k$个候选中时无法挽回。

\subsubsection{执行引导查询解码}
\label{sec:nlp-analysis-execution-guided}

语法与作用域检查发现不了所有运行期问题，例如类型不匹配的比较、对不存在函数的调用，或在当前数据上必然为空的条件。Wang等的执行引导解码把数据库的反馈接入候选筛选：对束中已经完整的查询（或可以独立执行的子句）试运行，排除执行出错的候选，在任务保证答案非空时也排除结果为空的候选\scite{v9a-wang2018execution}。试运行使用只读连接、固定的数据快照，并设置超时和行数上限；未写完的SQL前缀不能直接交给数据库，否则“尚未写完”的语法错误会被误当作候选不可行。它不需要重新训练生成器，可以接在任何能给出多个候选的模型之后，原文在WikiSQL上把它接到多种已有模型，执行准确率都有提升。

固定本节数据库，设三个候选为：把COUNT误写成不存在的函数的查询；连接正确但漏掉活动条件的查询；完整的参考查询。第一项在数据库中报错，被排除；后两项可运行，结果分别为2与1。执行引导无法仅凭“成功执行”判断后两项哪个回答了原问题。非空约束也只在任务保证答案非空时才合理：“有多少已确认的校内来宾”的正确答案在当前表中就是0。

碰巧相同的结果更隐蔽。若删去OTHER那条记录，漏掉活动条件的错误查询也会返回1；只有在含OTHER记录的数据库上，两者才显出差别。这正是本小节开头所述测试套件评测的出发点：在多个数据库变体上执行，才能区分语义不同而结果碰巧相同的查询。

\subsubsection{大语言模型文本到SQL}
\label{sec:nlp-analysis-llm-text-to-sql}

前四种方法都需要在目标任务上训练。大语言模型在预训练中见过大量SQL与代码，可以通过提示直接完成文本到SQL，把工作重心从训练转到“给模型看什么、让它分几步做、怎样检查结果”。

第一个问题是怎样把模式写进提示。Rajkumar等比较了多种模式表示，发现直接给出建表语句（\code{CREATE TABLE}，含列类型、主键与外键），并附上每张表的前三行样例数据，效果最好\scite{nlp-analysis-rajkumar2022}。建表语句与模型在预训练中见过的代码格式一致；样例行让模型看到取值的实际形式，例如status列存的是confirmed而不是“已确认”。在此之上可以加入少样本示例，Gao等的DAIL-SQL按问题相似度与SQL骨架相似度两方面选择示例，并比较了多种问题表示方式\scite{nlp-analysis-gao2024dail}。示例的选择与顺序对结果的影响见第~\ref{sec:nlp-inference-icl-shots}小节。数据库很大时，全部模式放不进上下文，需要先检索与问题相关的表和列，这一步本身就是模式链接。

第二个问题是复杂查询。一次生成连接与嵌套查询时错误较多，Pourreza与Rafiei的DIN-SQL把任务分解为四个依次执行的提示模块\scite{nlp-analysis-pourreza2023}。模式链接模块用逐步推理的提示找出问题涉及的表、列和值，例如把“校外来宾”对应到条件\code{people.affiliation = 'external'}。分类与分解模块把问题分为三类：单表无嵌套的简单查询、需要连接但无嵌套的查询、含子查询或集合运算的嵌套查询，对嵌套类再拆出子问题。生成模块按类别使用不同的提示与示例，连接类先生成一种省去连接条件的中间表示再转为SQL，嵌套类先解决子问题。自修正模块把生成的SQL交回模型，要求它检查并修正错误。四步都只靠提示，不训练参数；它用GPT-4在Spider测试集上达到85.3的执行准确率。

第三个问题是单次采样的随机错误。第~\ref{sec:nlp-execution-self-consistency}小节的自一致性在这里有一个自然的实现：以较高温度采样$K$条SQL，全部在数据库上执行，按执行结果分组，选择最大的一组中的一条\scite{nlp-llm-selfconsistency2023}。执行出错的候选直接丢弃。投票的对象是执行结果而不是SQL字符串，因为写法不同的查询可以等价。若$K=5$条候选中有3条返回1、1条返回2、1条报错，就输出返回1的那组。这种投票同样受“碰巧相同”的限制：多数候选犯同一种错误时，投票会把错误选出来。

第四个问题是怎样利用执行反馈修正。Chen等的自调试让模型先逐行解释自己生成的SQL，再结合执行结果或错误信息判断是否需要修改，迭代若干轮\scite{nlp-analysis-chen2024selfdebug}。例如数据库返回“no such function: COUNTT”，模型据此改写；即使没有报错，解释过程也可能让模型发现漏掉了活动条件。反馈的上限是可观测的信号：报错与空结果能直接看到，语义遗漏则只能靠模型自己的检查。

最后，也可以微调专门的代码模型。Li等的RESDSQL先用交叉编码器给表和列打相关性分数，只把排名靠前的模式项放进输入，再让解码器先生成SQL骨架、后填细节，这与第~\ref{sec:nlp-analysis-logic-generation}小节的粗到细解码思路相同\cite{nlp-analysis-li2023resdsql}；Li等的CodeS在开源代码模型上继续以SQL相关语料预训练，再在文本到SQL数据上微调，得到可以本地部署的较小模型\scite{nlp-analysis-li2024codes}。微调模型的推断成本低、行为稳定，但需要标注数据，迁移到新方言或新领域时要重新适配。

与RAT-SQL相比，大语言模型方法几乎不需要目标任务的训练数据，能处理更多样的SQL写法，并能利用模式说明与外部知识；代价是可控性：输出合法性要靠解析检查或约束解码补上，提示与示例的细微变化会改变结果，调用成本也高得多。在Spider 2.0这类企业级数据仓库任务上，模式规模、方言差异与多步工作流仍使成功率明显下降。表~\ref{tab:nlp-analysis-text-to-sql-methods}比较本小节的方法。

\begin{table*}[htbp]
  \centering
  \begin{tabularx}{\linewidth}{@{}l XXXX@{}}
    \toprule
    方法&解决的问题&核心机制&监督与推断&新代价\\
    \midrule
    SQLNet&条件顺序无关性&固定单表骨架；列注意力；集合预测&槽位监督；各槽独立预测&无法表达连接与嵌套\\
    RAT-SQL&模式链接与跨库泛化&关系感知自注意力；语法树动作解码&参考动作序列；语法合法的束搜索&需大量标注；文法与值链接需定制\\
    PICARD&预训练生成模型的非法输出&每步对前$k$个候选增量解析&不改训练；推断时屏蔽&解析开销；正确词元须在前$k$内\\
    执行引导&运行期错误&只读快照上试运行完整候选&不改训练；候选过滤&执行开销；无法识别碰巧正确\\
    大语言模型&减少任务训练；复杂查询&模式提示；分解；执行投票；执行反馈&提示或微调；多次调用&可控性弱；成本高；对提示敏感\\
    \bottomrule
  \end{tabularx}
  \caption{文本到SQL方法对照。PICARD、执行引导与自一致性都是推断阶段的附加机制，可以与不同的生成器组合，并不互斥。}
  \label{tab:nlp-analysis-text-to-sql-methods}
\end{table*}

\section{篇章分析}
\label{sec:nlp-analysis-discourse}

把每句话分析正确，仍可能不知道后一句的“他”是谁，也不知道它是在解释、反驳还是补充前一句。\term{篇章分析}（discourse analysis）处理跨句子的联系。图~\ref{fig:nlp-analysis-discourse}用补充段落“小王借到一本书。他读完后，归还了这本书。虽然封面已经磨损，但内容仍然清晰”展示三种不同的联系。“小王”与“他”指同一个人，“一本书”与“这本书”指同一本书，这是身份上的同一；“封面”与书是部分与整体，这是\term{桥接}（bridging），两者相关却不是同一对象；“虽然……但……”连接的是两个陈述，表达让步，联系的端点是内容单元而不是实体。

\begin{figure*}[htbp]
  \centering
  \begin{tikzpicture}[
  x=1mm,y=1mm,
  every node/.append style={font=\figurefont\footnotesize,align=center},
  mention/.style={draw=FigureInput,fill=FigureInput!10,
    rounded corners=1mm,inner sep=1.5mm,minimum height=8mm},
  identity/.style={draw=FigureInput,line width=1.1pt},
  bridge/.style={->,draw=FigureOutput,dashed,line width=1.1pt},
  discourse/.style={->,draw=FigureModel,dash dot,line width=1.1pt}
]
  \node[mention] (wang) at (20,0) {小王};
  \node at (42,0) {借到};
  \node[mention] (book) at (71,0) {一本书};
  \node at (86,0) {。};
  \node[mention] (he) at (20,-20) {他};
  \node at (42,-20) {读完后，};
  \node[circle,draw=FigureInput,line width=1.1pt,
    minimum size=3mm,inner sep=0] (zero) at (67,-20) {};
  \node at (90,-20) {归还了};
  \node[mention] (book2) at (128,-20) {这本书};
  \node at (145,-20) {。};

  \draw[identity] (wang.south) -- (he.north);
  \draw[identity] (wang.south) -- (20,-10) -- (67,-10) -- (zero.north);
  \draw[identity] (book.north) -- (71,9) -- (128,9) -- (book2.north);
  \node[anchor=south] at (113,1) {同一本书};
  \node[anchor=north] at (67,-23) {省略位置};

  \draw[draw=FigureModel,fill=FigureModel!5,rounded corners=1mm]
    (2,-59) rectangle (84,-43);
  \draw[draw=FigureModel,fill=FigureModel!5,rounded corners=1mm]
    (94,-59) rectangle (162,-43);
  \node at (12,-51) {虽然};
  \node[mention,draw=FigureOutput,fill=FigureOutput!10] (cover) at (32,-51) {封面};
  \node at (62,-51) {已经磨损，};
  \node at (128,-51) {但内容仍然清晰。};
  \draw[bridge] (cover.north) -- (32,-35) -- (128,-35) -- (book2.south);
  \node[fill=white,inner sep=1mm] at (102,-35) {PART-OF};
  \draw[discourse] (43,-59) -- (43,-69) -- (128,-69) -- (128,-59);
  \node[fill=white,inner sep=1mm] at (85,-69) {让步：两个陈述};

  \draw[identity] (3,-82) -- (15,-82);
  \node[anchor=west] at (18,-82) {同一身份};
  \draw[bridge] (55,-82) -- (67,-82);
  \node[anchor=west] at (70,-82) {桥接};
  \draw[discourse] (109,-82) -- (121,-82);
  \node[anchor=west] at (124,-82) {篇章关系};
\end{tikzpicture}

  \caption{同一段落中的身份共指、桥接与篇章关系。实线连接同一对象的提及，虚线连接“封面”与书的部分--整体关系，点划线连接篇章单元。空心位置标记不是原文字符，而是“归还”前检测出的省略主语位置，先行项为小王。共指可以形成等价簇，桥接和篇章关系不能并入该簇。}
  \label{fig:nlp-analysis-discourse}
\end{figure*}
\FloatBarrier

三种边的端点和性质不同，输出结构也就不同：身份共指满足传递性，可以合并为等价簇；篇章关系有方向和类型，还可能组织成覆盖全文的树；论证关系同样有方向，但关心的是一个单元是否为另一个单元提供理由。本节据此分为四个小节：共指处理提及的身份，篇章关系处理内容单元之间的语义联系，论证结构处理主张与理由，连贯性则不再输出具体的边，而是给整段文字打分或排定句序。长文档带来的截断和候选筛选问题会影响这四类任务的计算方式，但不构成新的关系类型。

\subsection{提及、先行项与共指}
\label{sec:nlp-analysis-coreference}

\term{提及}（mention）是文本中指向某个对象的一次表达，可以是“机器学习读书会”这样的名词短语，也可以是代词。\term{共指}（coreference）关系连接指向同一对象的提及，所有互相共指的提及构成一个共指簇。\term{先行项}（antecedent）是用来解释当前表达的较早提及。簇与先行项是两种数据结构：簇描述身份，先行项指针描述一条恢复路径；只要把指针看作无向边并取连通分量，指针就能还原为簇，而同一个簇可以对应许多组不同的指针。

有三类现象需要先划出边界。“封面”与书的桥接不合并身份；“活动改期了。这让小王调整了行程”中的“这”指向改期这一事件，属于话语指示，可能没有对应的名词跨度；中文还常省略主语，图中“归还”前的空位需要先检测出位置，再为它选先行项。本节主要讨论名词性提及的身份共指；零代词可视为在字符边界上增加的候选提及，检测与选择先行项分开计分。

共指消解面对的挑战可以归为五项。其一是提及检测：哪些跨度是提及，单例提及（只出现一次的实体）是否要输出。其二是长距离先行项：代词的先行项可能在几句之前，候选数量随文档长度增长。其三是一致性与语义兼容：数、性、生命度要相容，“那家公司”与“它”可以同指，与“他”通常不能。其四是簇级一致性：两两判断各自合理，合起来却可能违背传递性。其五是评价：误合并两个实体与漏连一条边，对不同指标的影响差别很大。

方法的演进大致对应这些挑战。Hobbs算法在句法树上按固定顺序搜索代词的先行项，依靠句法结构和一致性规则处理第三项；提及对模型用分类器判断两两是否同指，可以融合大量特征，但各次判断互相独立；提及排序模型对一个提及的全部候选先行项归一化，只选最好的一个，缓解了独立二分类的冲突；端到端跨度模型再把提及检测并入同一个目标，消除检测与连接之间的流水线误差，并用剪枝和粗到细推断控制候选规模。第~\ref{sec:nlp-analysis-coreference-metrics}小节给出共指评价指标，第~\ref{chap:nlp-information-extraction}章的事件共指沿用这些指标。

\subsubsection{Hobbs句法先行项搜索}
\label{sec:nlp-analysis-hobbs}

在没有标注语料可学的年代，代词消解的主要依据是句法结构：英语中先行项更常出现在主语位置、更常出现在离代词较近的句子里，而且同一小句里的宾语代词一般不与本句主语同指（“John likes him”中him不指John）。Hobbs的朴素算法把这些倾向写成一条在句法树上的固定搜索顺序，按顺序提出候选，用数、性等一致性条件逐个筛选，第一个通过的候选就是结果\scite{nlp-analysis-hobbs1978}。算法的输入是当前句与前文各句的表层句法树，步骤如下：
\begin{enumerate}
  \item 从直接支配代词的NP节点出发。
  \item 向上走到遇到的第一个NP或S节点，记为X，并记下走过的路径p。
  \item 在X之下、路径p左侧的分支中，按从左到右的广度优先顺序遍历，把遇到的、与X之间还隔着某个NP或S节点的NP提为候选。
  \item 若X是本句最高的S节点，转去按由近及远的顺序遍历前文各句的句法树，每棵树内部从左到右广度优先，遇到的NP依次提为候选；否则继续第5步。
  \item 从X继续向上，走到下一个NP或S节点，把它记为新的X，路径记为p。
  \item 若X是NP，且路径p没有经过X直接支配的$\bar N$节点，把X本身提为候选。
  \item 在X之下、路径p左侧的分支中从左到右广度优先遍历，把遇到的NP提为候选。
  \item 若X是S，再遍历X之下路径p右侧的分支，但遇到NP或S节点时不再深入其内部，把遇到的NP提为候选。
  \item 回到第4步。
\end{enumerate}
第3步中“中间隔着NP或S”这一条件，排除的正是同一小句内的主语，它对应上面所说的宾语代词约束；第4步体现由近及远，广度优先使主语先于宾语被提出，体现主语倾向。每个候选都要通过一致性检查，以及已知的非共指约束，才被接受。

以“Mary met John. She thanked him.”为例，求him的先行项，并设已知Mary用she指称、John用he指称。第1、2步从him的NP上升到当前句的S；第3步在路径左侧遇到She的NP，但它与S之间没有其他NP或S，不提出。当前S已是最高S，第4步转到前一句，广度优先先提出主语Mary，性不匹配被拒绝，随后提出宾语John并接受。输出是him指向John这次提及的指针。若求She的先行项，同样的过程会在前一句先提出Mary并接受。

Hobbs算法不需要训练，结果完全由句法树、搜索顺序和一致性词表决定，这使它便于解释，也让它成为后来学习方法的常用基线和特征来源。代价也来自同一处：句法树错了，候选次序随之改变；两个候选都满足形式条件时，搜索顺序可能选到语义上不通的一个，例如“The trophy didn't fit in the suitcase because it was too big”中，it指trophy还是suitcase要靠常识判断。算法只处理代词，不处理名词短语之间的共指；中文零代词没有可以出发的代词节点，需先检测省略位置。

\subsubsection{提及对与提及排序模型}
\label{sec:nlp-analysis-mention-ranking}

规则方法难以综合字符串匹配、中心词、距离、语法角色、语义类别等大量线索，而且规则之间的优先次序要人工调定。\term{提及对模型}（mention-pair model）把共指改写为监督分类：对每一对提及$(j,i)$（$j$在$i$之前），用特征向量$\bm x_{j,i}$预测两者是否同指，再由这些两两判断构造簇\scite{nlp-analysis-soon2001}。

训练样本的构造决定了模型学到什么。如果把所有提及对都作为样本，负例会远多于正例。Soon等的做法是：对每个有先行项的提及$i$，只取它在同一簇中最近的先行项$j^*$构成正例，再把$j^*$与$i$之间的每个提及与$i$配对作为负例。设$\sigma$为sigmoid函数，$f(\bm x_{j,i})$为分类器打分，以二元交叉熵为目标函数：
\begin{equation}
  L_{\mathrm{pair}}=-\sum_{(j,i)}\Bigl[y_{j,i}\log\sigma\bigl(f(\bm x_{j,i})\bigr)
   +(1-y_{j,i})\log\bigl(1-\sigma(f(\bm x_{j,i}))\bigr)\Bigr],
  \label{eq:nlp-analysis-mention-pair}
\end{equation}
其中$y_{j,i}\in\{0,1\}$表示两者是否同指。Soon等的分类器是C5决策树，特征只有12个，可以分为四组：距离（两者相隔的句子数）；提及形式（$j$、$i$各自是否为代词，$i$是否为有定名词短语或指示性名词短语，两者是否都是专名）；一致性（数、性、语义类别是否相容）；字符串与句法关系（去掉限定词后字符串是否相同，一个是否为另一个的别名或缩写，两者是否构成同位语）。字符串匹配与别名特征承担了专名之间的大部分连接，代词则主要依靠距离与一致性。后来的工作在此基础上增加了中心词、语法角色、句法路径等上百个特征，模型框架保持不变。

推断时，对每个提及$i$从近到远检查前面的提及，第一个概率超过0.5的就作为先行项，称为最近优先聚类；Ng与Cardie改为在所有超过阈值的候选中取概率最高者，称为最佳优先聚类\scite{nlp-analysis-ng-cardie2002}。最后把选中的连接取传递闭包，得到簇。

两两判断互相独立，带来局部决策之间的冲突。设A、B、C三个提及，分类器给出$p(A,B)=0.9$、$p(B,C)=0.8$、$p(A,C)=0.1$。若对所有超过阈值的对取连通分量，A--B与B--C两条边会把三者合成一簇，而模型其实相当确信A与C不同指；传递闭包修复了输出格式，却放大了一次错连。另一个问题是，二分类只问“这对是否同指”，不问“在所有候选中哪个最好”，因此一个代词可能有多个候选都略高于阈值，模型并不比较它们。

\term{提及排序模型}（mention-ranking model）直接回答后一个问题：对提及$i$，把所有较早提及与一个虚拟先行项$\epsilon$放在一起，做一次归一化，选出最合适的一个\scite{nlp-analysis-denis-baldridge2008}。$\epsilon$表示$i$不指向任何较早提及，即开启一个新簇（或$i$不是提及）。记候选集合$Y(i)=\{\epsilon,1,\ldots,i-1\}$，打分函数$s(i,j)$，并约定$s(i,\epsilon)=0$作为基准，则
\begin{equation}
  p(j\mid i)=\frac{\exp s(i,j)}{\sum_{j'\in Y(i)}\exp s(i,j')}.
  \label{eq:nlp-analysis-antecedent-softmax}
\end{equation}

训练时遇到一个新问题：标注给出的是簇，而不是唯一的先行项。“王老师……小王……他”中，“他”的前两个提及都同簇，两个都是正确先行项。若任意指定其中一个为唯一标签，会惩罚同样合理的另一个。Durrett与Klein把正确先行项当作隐变量，最大化正确先行项的概率之和\scite{nlp-analysis-durrett-klein2013}。设$G(i)\subseteq Y(i)$为与$i$同簇的较早提及集合，若$i$是所在簇的第一个提及则$G(i)=\{\epsilon\}$，边缘对数似然目标为
\begin{equation}
  L_{\mathrm{rank}}=-\sum_i\log\sum_{j\in G(i)}p(j\mid i).
  \label{eq:nlp-analysis-coreference-marginal}
\end{equation}
梯度只要求把概率质量从错误候选移到正确候选集合，至于集合内部怎样分配，模型可以自行决定，通常会集中到最容易判断的那个先行项上。

另一种目标是带代价的间隔损失。Wiseman等区分三类错误：本应开新簇却连到某个先行项（误连，false link），本应连接却选了$\epsilon$（误开新簇），以及连错了先行项（错连），分别赋予代价$\alpha_{\mathrm{FL}}$、$\alpha_{\mathrm{FN}}$、$\alpha_{\mathrm{WL}}$，记为$\Delta(i,y)$，正确选择的代价为0\scite{nlp-analysis-wiseman2015}。令$\ell_i=\arg\max_{j\in G(i)}s(i,j)$为当前得分最高的正确先行项，损失为
\begin{equation}
  L_{\mathrm{margin}}=\sum_i\max_{y\in Y(i)}\Delta(i,y)\bigl(1+s(i,y)-s(i,\ell_i)\bigr).
  \label{eq:nlp-analysis-coreference-margin}
\end{equation}
它要求每个错误候选比最好的正确候选至少低1分，代价越大的错误类型受到的惩罚越重；调节三个代价可以在“多连”与“少连”之间权衡，这对后面讨论的评价指标有直接影响，因为MUC与B³对两类错误的敏感程度不同。

提及排序仍是逐提及的局部决策：“他”选了“小王”，并不知道“小王”所在的簇里还有哪些提及，簇层面的性别、数等属性无法汇总。实体级模型让候选单位从单个提及换成已经形成的部分簇，用簇的汇总表示打分，代价是推断顺序影响结果、训练更复杂。后文的高阶推断用另一种方式引入簇信息：迭代地用先行项分布更新提及表示。

\subsubsection{端到端跨度共指与粗到细推断}
\label{sec:nlp-analysis-span-coreference}

前面的模型都假设提及已经由一个独立的检测器给出。检测器漏掉的提及在后续连接中无从补回，检测器多给的伪提及又会成为干扰候选，检测与连接各自训练，彼此不知道对方的需要。Lee等的端到端模型把全部跨度都视为可能的提及，用同一个目标同时学习“哪些跨度是提及”和“提及指向谁”\scite{lee2017coreference}。

\paragraph{跨度表示与打分}
设文档有$n$个词元，编码器给出上下文表示$\bm h_1,\ldots,\bm h_n\in\mathbb R^{d}$。对宽度不超过$L$的每个跨度$i$，记起点与终点为$s_i,e_i$，把跨度简记为$i$，按第~\ref{sec:nlp-analysis-paradigm-span}小节式~\eqref{eq:nlp-analysis-span-representation}构造跨度表示
\[
  \bm g_i=\bigl[\bm h_{s_i};\bm h_{e_i};\hat{\bm h}_i;\phi(i)\bigr],
\]
其中$\hat{\bm h}_i$是跨度内部词元按注意力加权的平均，作用近似于找出中心词；$\phi(i)$是宽度嵌入。在此基础上定义两个分数。提及分数衡量跨度$i$本身像不像一个提及：
\begin{equation}
  s_m(i)=\bm w_m^{\top}\operatorname{FFNN}_m(\bm g_i).
  \label{eq:nlp-analysis-mention-score}
\end{equation}
先行项分数衡量$j$作为$i$的先行项是否合适：
\begin{equation}
  s_a(i,j)=\bm w_a^{\top}\operatorname{FFNN}_a\bigl([\bm g_i;\bm g_j;\bm g_i\circ\bm g_j;\phi(i,j)]\bigr),
  \label{eq:nlp-analysis-antecedent-score}
\end{equation}
其中$\circ$为逐元素乘积，用于直接表达两个跨度在各维上的相似性；$\phi(i,j)$编码两者的距离（按区间分桶后嵌入），以及说话人是否相同、文档体裁等元信息。$\operatorname{FFNN}$是带ReLU的两层前馈网络，$\bm w_m,\bm w_a$把隐藏层映射为标量。三项相加得到配对分数：
\begin{equation}
  s(i,j)=s_m(i)+s_m(j)+s_a(i,j),\qquad s(i,\epsilon)=0.
  \label{eq:nlp-analysis-span-coreference-score}
\end{equation}
按张量形状看，若$\bm h_t\in\mathbb R^{d}$、宽度嵌入$\phi(i)\in\mathbb R^{d_w}$，则$\bm g_i\in\mathbb R^{3d+d_w}$，保留的$M$个跨度堆叠为$\bm G\in\mathbb R^{M\times(3d+d_w)}$，$s_m$对$\bm G$逐行计算得到长度为$M$的向量；先行项打分的输入$[\bm g_i;\bm g_j;\bm g_i\circ\bm g_j;\phi(i,j)]$长度为$3(3d+d_w)+d_\phi$，对每个提及及其$K$个候选计算一次，得到$M\times K$的分数矩阵，再在每行末尾补上$\epsilon$的0分，按行做softmax。

这一分解的意义在于，$i$或$j$只要有一个不像提及，$s_m$就会把配对分数拉低到$\epsilon$的0分以下，模型便选择$\epsilon$。于是“不是提及”和“是提及但开新簇”都由同一个$\epsilon$吸收，提及检测不再需要单独的标签：训练仍用式~\eqref{eq:nlp-analysis-coreference-marginal}的边缘似然，监督只来自共指簇，提及分数作为配对分数的一部分间接得到梯度。

\paragraph{计算量与剪枝}
直接计算的代价是不可接受的。宽度不限时候选跨度有$O(n^2)$个，任意两个跨度都可能构成先行关系，配对数为$O(n^4)$；即使限制宽度$L$，跨度数为$O(nL)$，配对仍有$O(n^2L^2)$个，每对都要跑一次前馈网络。Lee等用两级剪枝把规模降下来，第一级就是第~\ref{sec:nlp-analysis-paradigm-span}小节所说的按一元分数剪枝，$s_m(i)$即该节的一元分数：先按$s_m(i)$只保留得分最高的$\lambda n$个跨度（并去掉与更高分跨度交叉的跨度），再对每个保留的跨度只考虑最近的$K$个先行项。原文取$L=10$、$\lambda=0.4$、$K=250$，配对数降为$O(\lambda nK)$。剪枝发生在训练和推断两个阶段，被剪掉的跨度得不到连接；式~\eqref{eq:nlp-analysis-coreference-marginal}在这种情况下取$G(i)=\{\epsilon\}$，不会强迫模型连到已不在候选中的正确先行项。

图~\ref{fig:nlp-analysis-span-coreference}概括了这一流程。从左到右，候选数量在每一级被压缩，越靠右的打分越昂贵。

\begin{figure*}[htbp]
  \centering
  % 端到端跨度共指的计算流程：跨度枚举、提及打分剪枝、粗粒度先行项筛选与精细打分。
\begin{tikzpicture}[
  x=1mm,y=1mm,figure picture,
  stage/.style={figure node,minimum width=26mm,minimum height=15mm,
    align=center,font=\figurefont\scriptsize},
  size note/.style={font=\figurefont\scriptsize,text=FigureMuted,align=center},
  flow/.style={figure flow}
]
  \node[stage,draw=FigureInput] (enc) at (0,0)
    {编码器\\BiLSTM或SpanBERT\\$\bm H\in\mathbb R^{n\times d}$};
  \node[stage,draw=FigureInput] (span) at (32,0)
    {跨度表示$\bm g_i$\\端点、内部注意力\\与宽度嵌入};
  \node[stage,draw=FigureModel] (ment) at (64,0)
    {提及分数$s_m(i)$\\保留前$\lambda n$个};
  \node[stage,draw=FigureModel] (coarse) at (96,0)
    {粗分$s_c(i,j)$\\每个提及保留前$K$个};
  \node[stage,draw=FigureOutput] (fine) at (128,0)
    {精细分$s_a(i,j)$\\对$Y(i)\cup\{\epsilon\}$\\做softmax};
  \draw[flow] (enc) -- (span);
  \draw[flow] (span) -- (ment);
  \draw[flow] (ment) -- (coarse);
  \draw[flow] (coarse) -- (fine);
  \node[size note] at (32,-13) {约$nL$个候选跨度};
  \node[size note] at (64,-13) {$M=\lambda n$个提及};
  \node[size note] at (96,-13) {$M^2$对只需双线性运算};
  \node[size note] at (128,-13) {$MK$对送入前馈网络};
  \node[size note] at (0,-13) {$n$个词元};
\end{tikzpicture}

  \caption{端到端跨度共指的计算流程。编码器输出逐词元表示，宽度不超过$L$的跨度各得到表示$\bm g_i$；提及分数保留$M=\lambda n$个跨度，粗粒度双线性分数为每个提及保留$K$个候选先行项，只有这$MK$对才进入前馈网络计算精细分数，并与虚拟先行项$\epsilon$一起归一化。下方文字标出每一级的候选规模。}
  \label{fig:nlp-analysis-span-coreference}
\end{figure*}

\paragraph{粗到细推断与高阶推断}
按距离保留最近的$K$个先行项是一种粗糙的筛选：远处的正确先行项会被剪掉，近处大量无关跨度却保留下来。Lee等在后续工作中加入一个廉价的粗粒度分数\cite{nlp-analysis-lee2018higher}：
\begin{equation}
  s_c(i,j)=\bm g_i^{\top}\bm W_c\bm g_j,
  \label{eq:nlp-analysis-coarse-score}
\end{equation}
它对全部$M$个保留跨度两两计算，只是一次$M\times M$的矩阵乘法，不需要逐对运行前馈网络。以$s_m(i)+s_m(j)+s_c(i,j)$为依据为每个$i$选出前$K$个候选（此时$K$可以小得多，如50），再对这些候选计算$s_a$，最终分数为$s_m(i)+s_m(j)+s_c(i,j)+s_a(i,j)$。粗分负责召回，精细分负责区分，二者联合训练。

同一篇工作还引入高阶推断，弥补提及排序缺少簇信息的问题。第$t$轮先按当前表示计算先行项分布$p^{(t)}(j\mid i)$，再用它对候选先行项的表示加权，得到“$i$所在簇的期望表示”，并用门控与原表示融合：
\[
  \begin{aligned}
    \bm a_i^{(t)}&=\sum_{j\in Y(i)}p^{(t)}(j\mid i)\,\bm g_j^{(t)},\qquad
    \bm f_i^{(t)}=\sigma\bigl(\bm W_f[\bm g_i^{(t)};\bm a_i^{(t)}]\bigr),\\
    \bm g_i^{(t+1)}&=\bm f_i^{(t)}\circ\bm g_i^{(t)}+(1-\bm f_i^{(t)})\circ\bm a_i^{(t)},
  \end{aligned}
\]
论文只规定$s(i,\epsilon)=0$，公开实现中把$\epsilon$的表示取为$\bm g_i^{(t)}$本身，即选择$\epsilon$的概率越大，$i$的表示越保持不变。迭代一两轮后，“他”的表示混入了“王老师”的信息，再与后续提及打分时便能利用整个簇的属性。

算法~\ref{alg:nlp-analysis-span-coreference}汇总了推断过程。

\begin{algorithm}[粗到细的端到端跨度共指推断]
  \label{alg:nlp-analysis-span-coreference}
  \begin{algorithmic}[1]
    \Require 文档词元$x_1,\ldots,x_n$；最大宽度$L$，保留比例$\lambda$，候选数$K$，高阶迭代轮数$T$
    \Ensure 共指簇集合
    \State 编码得到$\bm h_1,\ldots,\bm h_n$，为全部宽度不超过$L$的跨度计算$\bm g_i$与$s_m(i)$
    \State 按$s_m$降序选取跨度，跳过与已选跨度交叉者，直到选满$M=\lambda n$个，按位置排序
    \State 计算全部$M\times M$的$s_c(i,j)$，对每个$i$只保留$j<i$中$s_m(i)+s_m(j)+s_c(i,j)$最高的$K$个，记为$Y(i)$
    \For{$t=1,\ldots,T$}
      \State 对$j\in Y(i)$计算$s(i,j)=s_m(i)+s_m(j)+s_c(i,j)+s_a(i,j)$，令$s(i,\epsilon)=0$，得到$p^{(t)}(j\mid i)$
      \State 按高阶推断更新全部$\bm g_i$
    \EndFor
    \State 对每个$i$取$\hat y_i=\arg\max_{j\in Y(i)\cup\{\epsilon\}}s(i,j)$
    \State 把$\hat y_i\neq\epsilon$的指针看作无向边，取连通分量，丢弃只含一个跨度的分量，输出簇
  \end{algorithmic}
\end{algorithm}

\paragraph{算例}
设文本已说明“小王”是“王老师”的别称，依次出现$m_1$“王老师”、$m_2$“小王”、$m_3$“他”，以下分数为教学构造。$m_2$到$m_1$的分数为2，到$\epsilon$为0，选$m_1$。$m_3$到$m_1,m_2,\epsilon$的分数为$2,3,0$，按式~\eqref{eq:nlp-analysis-antecedent-softmax}，概率约为$0.259,0.705,0.035$。$G(m_3)=\{m_1,m_2\}$，边缘概率为$0.965$，损失$-\log0.965\approx0.036$；推断选$m_2$，连通分量为$\{m_1,m_2,m_3\}$。若随后出现“另一位王老师”$m_4$，它到前三者的分数为$-2,-3,-1$，均低于$\epsilon$的0分，于是开启新簇。只比较名字会把$m_4$错并到第一个簇，这里是$s_a$中的上下文表示与修饰语“另一位”压低了分数。

\paragraph{编码器的更替}
原始模型的编码器是以GloVe与Turian词向量、字符卷积为输入的BiLSTM。把编码器换成预训练的BERT后，长文档需要切成不超过512个子词（实验中较短的片段反而更好）的不重叠片段分别编码，片段之间不共享注意力，跨片段的先行项只能依靠跨度表示本身连接\scite{nlp-analysis-joshi2019bert}。SpanBERT在预训练中遮蔽连续跨度并用边界表示预测跨度内容（第~\ref{sec:nlp-classical-spanbert}小节），正好强化了式中$\bm h_{s_i},\bm h_{e_i}$对整个跨度的概括能力；以它为编码器，在SpanBERT原论文的对照中，同一共指模型在OntoNotes英文测试集上的CoNLL F1从使用BERT时的77.1提高到79.6\scite{nlp-classical-joshi2020}。

端到端模型的优势是不需要独立的提及检测器和人工特征，代价是显存随文档长度快速增长，剪枝参数需要权衡召回与计算量；OntoNotes不标注单例，模型因而学会不输出单例，迁移到标注单例的数据集时需要调整。

共指也可以按第~\ref{sec:nlp-analysis-paradigm-generation}小节的约束生成范式求解。Bohnet等把共指写成转移系统的动作序列，再线性化为文本：模型逐句读入文档，对当前句输出一串动作：把某个提及追加到已有的编号簇，或把它与某个较早的提及连接成新簇；动作中的提及以原文字符串加其后的三个词表示，生成后再对齐回原文跨度\scite{nlp-analysis-bohnet2023}。这种做法不需要跨度枚举与剪枝，可以直接利用大型序列到序列模型，代价是解码按句串行、速度较慢，输出的提及字符串必须能在原文中唯一定位。

中文零代词需要在上述框架之外先解决位置问题。常见做法是在谓词前的主语与宾语空位上生成候选零代词位置，判断该位置是否确有省略，再像代词一样从前文的名词短语中选先行项；Chen与Ng用前馈网络分别学习零代词位置与候选先行项的表示并打分\scite{nlp-analysis-chen-ng2016}。位置检测与先行项选择两步的错误会叠加，评价时须说明零代词位置是人工给定还是系统检测的。

\subsubsection{共指评价指标}
\label{sec:nlp-analysis-coreference-metrics}

共指的输出是一组簇，参考答案也是一组簇，比较两组划分没有唯一自然的方式。一个错误的连接可能把两个大实体合并，也可能只多连了一个代词，不同指标对这两类错误给出的惩罚差别很大。设参考簇（key）为$K_1,\ldots,K_a$，系统簇（response）为$R_1,\ldots,R_b$，下面用同一个例子计算三种常用指标。例子有五个提及：$a$“王老师”、$b$“小王”、$c$“他”、$d$“李同学”、$e$“她”。参考答案为$K_1=\{a,b,c\}$、$K_2=\{d,e\}$；系统把“他”错连到李同学，输出$R_1=\{a,b\}$、$R_2=\{c,d,e\}$。

\paragraph{MUC}
MUC指标以连接为单位\scite{nlp-analysis-vilain1995}。一个含$|K|$个提及的簇至少需要$|K|-1$条连接才能连通；系统输出把它切成$|p(K)|$块，其中$p(K)$表示用系统簇划分$K$得到的部分（系统中没有出现的提及各自成为一块），于是系统恢复了$|K|-|p(K)|$条必要连接。召回率与精确率为
\begin{equation}
  R_{\mathrm{MUC}}=\frac{\sum_k\bigl(|K_k|-|p(K_k)|\bigr)}{\sum_k\bigl(|K_k|-1\bigr)},\qquad
  P_{\mathrm{MUC}}=\frac{\sum_l\bigl(|R_l|-|p'(R_l)|\bigr)}{\sum_l\bigl(|R_l|-1\bigr)},
  \label{eq:nlp-analysis-muc}
\end{equation}
其中$p'(R)$是用参考簇划分$R$得到的部分。例子中，$K_1$被系统分成$\{a,b\}$与$\{c\}$两块，恢复$3-2=1$条连接；$K_2$完整地落在$R_2$中，恢复1条，召回率$(1+1)/(2+1)=2/3$。反过来，$R_1$落在$K_1$中，贡献1；$R_2$被参考分成$\{c\}$和$\{d,e\}$，贡献$3-2=1$，精确率$2/3$，F1为$0.667$。

MUC只数连接，带来两个已知缺陷：单例簇不需要任何连接，对分子分母都没有贡献，因此MUC完全忽略单例；把所有提及并成一个大簇只增加少量“多余连接”，精确率下降有限，召回率却达到1，误合并受到的惩罚偏轻。

\paragraph{B³}
B³以提及为单位\scite{nlp-analysis-bagga1998}。对每个提及$m$，记$K(m)$与$R(m)$分别为它所在的参考簇和系统簇，比较两者的重合程度，再对全部$N$个提及取平均：
\begin{equation}
  R_{\mathrm{B^3}}=\frac1N\sum_m\frac{|K(m)\cap R(m)|}{|K(m)|},\qquad
  P_{\mathrm{B^3}}=\frac1N\sum_m\frac{|K(m)\cap R(m)|}{|R(m)|}.
  \label{eq:nlp-analysis-bcubed}
\end{equation}
例子中，召回率各项为：$a,b$各$2/3$，$c$为$1/3$，$d,e$各$2/2=1$，平均$11/15\approx0.733$；精确率各项为：$a,b$各1，$c$为$1/3$，$d,e$各$2/3$，平均同样为$11/15$。B³让每个提及都参与，单例也计入；大簇里的提及多，因而在平均中权重也大。

\paragraph{CEAF}
CEAF先在参考簇与系统簇之间求一个一对一的最优对齐，再按对齐结果计算得分，每个簇最多被使用一次，避免了一个系统簇同时“匹配”多个参考簇\scite{nlp-analysis-luo2005}。设$\phi(K,R)$为一对簇的相似度，$g$为从参考簇到系统簇的一对一映射，$g^*$使总相似度$\Phi(g)=\sum_k\phi(K_k,R_{g(k)})$最大，可用匈牙利算法求得。召回率与精确率为
\begin{equation}
  R_{\mathrm{CEAF}}=\frac{\Phi(g^*)}{\sum_k\phi(K_k,K_k)},\qquad
  P_{\mathrm{CEAF}}=\frac{\Phi(g^*)}{\sum_l\phi(R_l,R_l)}.
  \label{eq:nlp-analysis-ceaf}
\end{equation}
相似度有两种常用取法。基于提及的$\phi_3(K,R)=|K\cap R|$数共有提及，记为CEAF$_m$；基于实体的$\phi_4(K,R)=2|K\cap R|/(|K|+|R|)$对每对簇归一化，记为CEAF$_e$，此时$\phi_4(K,K)=1$，分母就是参考簇数与系统簇数，每个实体同等重要，不论大小。例子中，$\phi_4(K_1,R_1)=2\cdot2/5=0.8$，$\phi_4(K_1,R_2)=2/6\approx0.333$，$\phi_4(K_2,R_1)=0$，$\phi_4(K_2,R_2)=2\cdot2/5=0.8$。最优对齐为$K_1\to R_1$、$K_2\to R_2$，$\Phi=1.6$，召回率与精确率都是$1.6/2=0.8$。

\paragraph{CoNLL F1与对比}
三种指标各有偏向，CoNLL-2011/2012共享任务取MUC、B³、CEAF$_e$三者F1的算术平均作为主指标，称为CoNLL F1\scite{nlp-analysis-pradhan2012}。例子的CoNLL F1为$(0.667+0.733+0.8)/3\approx0.733$。不同实现对未匹配到参考的系统提及处理不一，Pradhan等给出了参考实现，报告结果时应使用同一评分脚本\scite{nlp-analysis-pradhan2014}。

表~\ref{tab:nlp-analysis-coreference-metrics}把同一参考答案下的三种系统输出放在一起：上面的错连输出、把五个提及全部并成一簇的误合并输出，以及每个提及各自成簇的全单例输出。

\begin{table}[htbp]
  \centering
  \small
  \begin{tabular}{@{}lcccc@{}}
    \toprule
    系统输出 & MUC F1 & B³ F1 & CEAF$_e$ F1 & CoNLL F1\\
    \midrule
    错连：$\{a,b\},\{c,d,e\}$ & 0.667 & 0.733 & 0.800 & 0.733\\
    全部合并：$\{a,b,c,d,e\}$ & 0.857 & 0.684 & 0.500 & 0.680\\
    全单例：$\{a\},\{b\},\{c\},\{d\},\{e\}$ & 0 & 0.571 & 0.333 & 0.302\\
    \bottomrule
  \end{tabular}
  \caption{同一参考答案$\{a,b,c\},\{d,e\}$下三种系统输出的共指指标（教学算例）。}
  \label{tab:nlp-analysis-coreference-metrics}
\end{table}

表中最值得注意的是第二行：全部合并在MUC上的得分高于只错连一个代词的系统，因为它召回了全部连接，只多出一条；B³与CEAF$_e$则明显惩罚这种误合并，CEAF$_e$因为只能把一个系统簇对齐到一个参考簇，惩罚最重。第三行相反，全单例在MUC上为0，却在B³上拿到0.571，因为每个提及至少与自己重合。没有哪一个指标单独可靠，这正是取三者平均的理由；分析错误时仍应分别查看三个指标，以及提及检测本身的召回率与精确率。Moosavi与Strube指出，B³与CEAF在系统提及与参考提及不一致时也会出现反直觉的结果，并提出按实体大小加权、以连接为单位计分的LEA指标\cite{nlp-analysis-moosavi2016}；目前英文共指的主要报告口径仍是CoNLL F1。

这些指标最常用的数据是CoNLL-2012共享任务的OntoNotes~5.0语料，它包含英语、汉语和阿拉伯语三部分，英语部分约有2800篇训练文档、约340篇开发文档与约350篇测试文档，体裁覆盖新闻、广播对话、电话交谈、网络文本等。OntoNotes不标注单例，也不标注与通指名词短语的共指，因此在它上面训练的系统默认不输出单例。

表~\ref{tab:nlp-analysis-coreference-methods}对照本小节的四类方法。

\begin{table*}[htbp]
  \centering
  \small
  \begin{tabularx}{\linewidth}{@{}l XXX@{}}
    \toprule
    方法 & 解决的问题 & 核心机制 & 新的代价或局限\\
    \midrule
    Hobbs算法 & 无标注数据时消解代词 & 句法树上的固定搜索顺序加一致性筛选 & 依赖句法树质量；只处理代词；不用语义\\
    提及对模型 & 综合大量特征 & 两两二分类，最近优先或最佳优先聚类 & 判断互相独立，传递闭包放大错连；候选间不比较\\
    提及排序模型 & 候选间的比较 & 对全部候选与$\epsilon$做softmax，正确先行项作隐变量 & 仍是逐提及决策，缺少簇级信息\\
    端到端跨度模型 & 提及检测与连接的流水线误差 & 跨度表示、$s_m+s_a$分解、剪枝与粗到细筛选 & 显存随长度增长；剪枝损失召回\\
    \bottomrule
  \end{tabularx}
  \caption{共指消解方法对照。从上到下，每一行解决上一行留下的主要问题。}
  \label{tab:nlp-analysis-coreference-methods}
\end{table*}

\subsection{篇章单元与关系}
\label{sec:nlp-analysis-discourse-relations}

身份共指连接对象，篇章关系连接承载内容的单元。\term{基本篇章单元}（elementary discourse unit, EDU）是标注规范选定的最小分析片段，通常接近一个分句。用“下雨了，所以路面湿了。车辆减速。”作为本小节的统一材料，切成三个单元：$u_1$“下雨了，”、$u_2$“所以路面湿了。”、$u_3$“车辆减速。”。$u_1$与$u_2$之间有连接词“所以”，表达结果关系；$u_2$与$u_3$之间没有连接词，读者仍会理解为结果关系。这里分析的是文字所表达的关系，不是对现实因果的验证。

篇章关系有两种主要的标注框架。宾州篇章树库（Penn Discourse Treebank, PDTB）以连接词为中心，标注一个连接词及其两个论元（Arg1、Arg2）和关系类别，不要求全文形成一棵树\scite{nlp-analysis-prasad2008}。PDTB的关系类别分三级：一级为时序（Temporal）、偶然（Contingency）、比较（Comparison）、扩展（Expansion）四类，二级细分为十余种类型（隐式关系分类常用其中11类），三级再区分方向等细节。没有显式连接词的相邻句对，标注者会补写一个最合适的连接词，再按它标注关系，这类关系称为\term{隐式篇章关系}（implicit discourse relation）。\term{修辞结构理论}（rhetorical structure theory, RST）则要求用关系把全部EDU递归地组合成覆盖全文的树，每个内部节点记录关系，以及子节点的\term{核性}（nuclearity）：核心（nucleus）承载作者的主要意图，附属（satellite）为核心提供原因、背景、阐述等支持\scite{nlp-analysis-mann-thompson1988}。RST篇章树库（RST-DT）在《华尔街日报》的385篇文章上按这一理论标注，78种细粒度关系被归为16类\scite{nlp-analysis-carlson2001}；解析研究通常再加上文本组织（Textual-Organization）与同一单元（Same-Unit）两种结构关系，按18类评价。

篇章关系分析的挑战可以列为五项：单元切分，即EDU边界在哪里；连接词识别与消歧，同一个词可能不在篇章层面起连接作用，起连接作用时也可能表达多种关系（英语since既可表原因也可表时间）；隐式关系，缺少连接词这一最强信号；树结构，RST要求层次决策互相一致；标注稀缺，RST-DT只有数百篇文档。方法谱系与之对应：边界标注加显式连接词分类处理有信号的情形；隐式关系分类用论元对编码器从两段的语义推断关系；RST移进--归约按第~\ref{sec:nlp-analysis-paradigm-transition}小节的转移范式自底向上构建整棵树；递归分割则自顶向下决定树的划分，避免早期局部合并的错误一直保留到最后。

\subsubsection{篇章单元切分与显式关系分类}
\label{sec:nlp-analysis-discourse-boundary}

EDU切分可以写成序列标注：每个字或词标B（单元开始）或I（单元内部），按第~\ref{sec:nlp-analysis-paradigm-tagging}小节的方法训练，句首位置强制为B。上面16个字符的材料在位置0、4、11取B，其余取I，解码后恢复$[0,4)$、$[4,11)$、$[11,16)$三个单元。标点在本例归入左侧单元，不同规范可以不同，但恢复出的单元必须覆盖全部原文，且彼此不重叠。英语RST-DT上的EDU切分主要依赖句法与标点，用预训练编码器加序列标注已经相当准确，篇章解析中的主要误差来自后面的结构与关系决策。

显式关系分类分两步。第一步判断候选连接词是否真在篇章层面起作用：“所以”多为连接用法，“而且”“然后”却也常出现在小句内部；英语and有大量非篇章用法。Pitler与Nenkova用连接词本身及其句法上下文（所在成分的类别、父节点和左右兄弟节点）作为特征，区分篇章用法与非篇章用法，并在此基础上判断关系类别\scite{nlp-analysis-pitler-nenkova2009}。第二步给已确认的连接词分配关系类别：大多数连接词的关系分布高度集中，在PDTB上仅用连接词本身作特征，一级四分类的准确率就超过九成\scite{nlp-analysis-pitler2008}。这一事实说明，显式关系的难点不在关系分类本身，而在论元跨度的定位，尤其是Arg1可能在前几句之中。

对$u_1,u_2$，编码两段及连接词，在原因、对比、阐述等关系上做softmax。设三个logit为$2,0,-1$，原因关系的概率约为$0.844$，输出“原因论元$u_1$，结果论元$u_2$，连接词‘所以’”的记录。关系的方向必须与论元对应：把$u_2$当原因论元，类别名正确也是错误的输出。因此评价时应分别报告连接词识别、论元跨度和关系三项。

把这些步骤串起来，就得到端到端的PDTB式篇章解析器。Lin等的流水线依次执行：连接词分类器判断候选词是否为篇章用法；论元位置分类器判断Arg1与连接词在同一句还是在前面的句子中；论元抽取器在句法树上选出两个论元的跨度；显式关系分类器给出类别；最后对相邻句对中没有显式连接词的位置，再判断是否存在隐式关系及其类别\scite{nlp-analysis-lin2014pdtb}。流水线每一步都依赖上一步的输出，端到端结果明显低于各子任务在给定前一步正确输入时的成绩，这也是DISRPT等评测把分割、连接词识别与关系分类分别作为独立任务报告的原因\scite{braud-etal-2023-disrpt}。

\subsubsection{隐式篇章关系分类}
\label{sec:nlp-analysis-implicit-relation}

隐式关系难在什么地方，可以对照显式关系看。显式关系中，连接词几乎决定了类别；隐式关系中，这个信号不存在，模型只能从两段的内容推断联系。“路面湿了”与“车辆减速”之间的结果关系，依赖“湿滑路面导致减速”这一常识；把$u_3$换成“行人打着伞”，关系就变为并列或阐述。推断需要比较两段的语义，而不是只看其中一段。标注本身也更难，标注者在补写连接词时意见分歧较大。因此在PDTB上，隐式关系分类的准确率长期远低于显式关系\scite{nlp-analysis-pitler2009}。

早期方法用词对特征：把Arg1中的每个词与Arg2中的每个词组合成特征，期望“湿—减速”这类词对反映关系。词对数量巨大且稀疏，泛化很差。神经方法改用\term{论元对编码器}（argument pair encoder）：分别编码两个论元，再让它们互相注意，最后分类。下面给出一个代表性结构的逐层运算。

设两个论元的词嵌入为$\bm X^{(1)}\in\mathbb R^{n_1\times d}$与$\bm X^{(2)}\in\mathbb R^{n_2\times d}$。第一层用共享参数的BiLSTM分别编码，得到$\bm H^{(1)}\in\mathbb R^{n_1\times 2h}$与$\bm H^{(2)}\in\mathbb R^{n_2\times 2h}$。第二层计算交互矩阵并做双向软对齐：
\begin{equation}
  \begin{aligned}
    \bm E&=\bm H^{(1)}\bm W\bm H^{(2)\top}\in\mathbb R^{n_1\times n_2},\\
    \tilde{\bm H}^{(1)}&=\operatorname{softmax}_{\mathrm{row}}(\bm E)\,\bm H^{(2)},\qquad
    \tilde{\bm H}^{(2)}=\operatorname{softmax}_{\mathrm{col}}(\bm E)^{\top}\bm H^{(1)},
  \end{aligned}
  \label{eq:nlp-analysis-implicit-interaction}
\end{equation}
其中$\bm E_{p,q}$衡量Arg1第$p$个词与Arg2第$q$个词的关联，$\tilde{\bm H}^{(1)}\in\mathbb R^{n_1\times 2h}$的第$p$行是Arg2中与该词对齐的内容。第三层把原表示与对齐表示组合为$[\bm H^{(1)};\tilde{\bm H}^{(1)};\bm H^{(1)}-\tilde{\bm H}^{(1)};\bm H^{(1)}\circ\tilde{\bm H}^{(1)}]\in\mathbb R^{n_1\times 8h}$，差与积直接表达两侧的差异与一致；对它沿词的维度做最大池化与平均池化，得到$\bm r^{(1)}\in\mathbb R^{16h}$，Arg2同理得到$\bm r^{(2)}$。最后一层把$[\bm r^{(1)};\bm r^{(2)}]$送入多层感知机，在$C$个关系上做softmax，以交叉熵训练。

用一个二维小例子看交互层做了什么（数值为教学构造）。设Arg1为“路面／湿”，Arg2为“车辆／减速”，$\bm H^{(1)}$的两行为$(1,0)$与$(0,1)$，$\bm H^{(2)}$的两行为$(1,0)$与$(0,2)$，取$\bm W=\bm I$。则$\bm E=\bigl(\begin{smallmatrix}1&0\\0&2\end{smallmatrix}\bigr)$，按行softmax得到“路面”对Arg2两词的权重约为$(0.731,0.269)$，“湿”约为$(0.119,0.881)$；对应的对齐表示为$\tilde{\bm h}^{(1)}_1\approx(0.731,0.538)$、$\tilde{\bm h}^{(1)}_2\approx(0.119,1.762)$。“湿”主要对齐到“减速”，后续的差与积特征就围绕这一对词计算，分类器得以利用“湿—减速”这类跨论元的联系，而不必在稀疏的词对特征中查找。

这一结构与第~\ref{sec:nlp-dialogue-decomposable}小节判断证据关系的可分解注意力同属“对齐—比较—汇聚”的句对模型；Liu与Li的多层注意力模型则让两个论元反复互相读取多轮\scite{nlp-analysis-liu-li2016}。

预训练编码器把这一过程交给自注意力。输入“[CLS] Arg1 [SEP] Arg2 [SEP]”，两个论元的词元在每一层都互相注意，取$\bm h_{\mathrm{[CLS]}}$接线性层分类即可。Shi与Demberg在《华尔街日报》文本上继续预训练BERT时比较了保留与去掉下一句预测（NSP）目标两种做法，去掉NSP后隐式关系分类的准确率下降，他们据此认为NSP有助于判断“两段能否前后衔接”\scite{nlp-analysis-shi-demberg2019}。这一比较针对的是继续预训练阶段，原始BERT本身仍经过NSP预训练。

隐式关系分类的评价设置需要说明清楚，否则数字不可比。PDTB 2.0上有两种常用划分：一种以第2--20节为训练集、第0--1节为开发集、第21--22节为测试集，另一种以第2--21节训练、第22节开发、第23节测试；任务可以是一级四类的单一多分类，也可以是对每个一级类别做“一类对其余”的二分类，后者报告各类F1；二级任务通常在11个类型上做多分类并报告准确率。由于类别极不均衡（扩展类占多数），四分类常同时报告宏平均F1。不同论文还在是否包含实体关系（EntRel）、是否把多标签样本拆开计算上有差别，比较前应核对这些设置。

标注数据不足是另一个瓶颈。一个自然的想法是利用大量带显式连接词的文本：去掉连接词，把它当作隐式样本，标签由连接词决定\scite{nlp-analysis-marcu-echihabi2002}。Sporleder与Lascarides的实验说明这样得到的样本与真实隐式样本分布不同：作者省略连接词的句对，往往是关系本来就容易从内容看出的句对，而人为删去连接词的句对，内容未必支持该关系，于是在人造样本上训练的分类器迁移到真实隐式关系时效果很差\scite{nlp-analysis-sporleder2008}。较稳妥的用法是把连接词作为中间预测目标：先让模型预测“这里最适合补写哪个连接词”，再由连接词推断关系，或者把连接词预测作为辅助任务与关系分类联合训练\scite{nlp-analysis-zhou2010}。

\subsubsection{RST移进--归约篇章解析}
\label{sec:nlp-analysis-rst-transition}

RST要求把全部EDU组合成一棵树。自底向上的移进--归约是RST解析的主要做法之一，Ji与Eisenstein把它与可学习的EDU表示结合，在每个状态上用学到的表示选择动作\scite{nlp-analysis-ji-eisenstein2014}。它与第~\ref{sec:nlp-analysis-shift-reduce}小节的移进--归约成分解析在结构上同构：叶节点由词换成EDU，短语类别换成“关系＋核性”，状态仍由子树栈和待读队列组成。因此这里只写两者的差异。

动作集合为SHIFT与REDUCE$(r,\nu)$。SHIFT把队首EDU入栈；REDUCE$(r,\nu)$弹出栈顶两棵相邻子树，以关系$r$和核性$\nu\in\{\mathrm{NS},\mathrm{SN},\mathrm{NN}\}$组成父树后入栈，其中NN对应列举、对比等多核心关系。与成分解析相比，它没有一元动作，篇章树中不存在单子节点的内部节点；归约动作的标签空间是关系数乘以核性组合，约为$18\times3$量级，远大于成分解析的短语类别数，而训练样本少得多。为此，Wang等提出两阶段方法：第一阶段的转移系统只预测不带关系的树结构与核性，动作只有SHIFT与三种REDUCE；第二阶段再为每个内部节点分类关系，并对句内、跨句、跨段落的节点使用不同的分类器\scite{nlp-analysis-wang2017twostage}。

表~\ref{tab:nlp-analysis-rst-actions}列出两种系统动作的对应关系。

\begin{table}[htbp]
  \centering
  \small
  \begin{tabularx}{\linewidth}{@{}l XX@{}}
    \toprule
    & 成分移进--归约 & RST移进--归约\\
    \midrule
    叶节点 & 词 & EDU\\
    SHIFT & 队首词入栈 & 队首EDU入栈\\
    归约 & REDUCE-L/R-$X$：合并栈顶两项，标短语类别$X$，记录中心词方向 & REDUCE$(r,\nu)$：合并栈顶两项，标关系$r$与核性$\nu$\\
    一元动作 & UNARY-$X$ & 无\\
    子树代表 & 中心词 & 核心EDU\\
    终止 & 队列空且栈中一棵树 & 同左\\
    \bottomrule
  \end{tabularx}
  \caption{成分移进--归约与RST移进--归约的动作对照。核性中的N方向对应成分解析中的中心词方向。}
  \label{tab:nlp-analysis-rst-actions}
\end{table}

训练与解码沿用第~\ref{sec:nlp-analysis-paradigm-transition}小节：参考树经后序遍历给出唯一的静态oracle动作序列，在每个状态上对合法动作做softmax，以式~\eqref{eq:nlp-analysis-transition-loss}的交叉熵训练；栈中少于两棵子树时REDUCE非法，队列为空时SHIFT非法。状态特征取栈顶两棵子树与队首EDU的表示，子树表示可由其核心EDU的表示代表，这与成分解析中用中心词代表短语是同一种思路；此外还加入两棵子树是否在同一句、同一段落等结构特征，因为句内关系与跨段落关系的分布差别很大。

对三个单元，参考树为$[[u_1^{\mathrm S}u_2^{\mathrm N}]_{\text{原因}}^{\mathrm S}u_3^{\mathrm N}]_{\text{原因}}$，动作序列为SHIFT、SHIFT、REDUCE(原因,SN)、SHIFT、REDUCE(原因,SN)。另一条合法轨迹SHIFT、SHIFT、SHIFT、REDUCE、REDUCE先合并$u_2,u_3$，得到右分支的树。设前一条轨迹的动作分数之和为8，后一条为6，束搜索在完整轨迹之间比较时会选前者；贪心解码若在第三步就因局部分数选择了SHIFT，此后再无法组成$[u_1u_2]$这一子树。

\paragraph{RST-Parseval评价}
RST树的评价沿用成分句法的Parseval思想，比较参考树与预测树中的跨度。Marcu提出的RST-Parseval把每个非根节点（包括叶节点EDU）看作一个评价单位，并给它附上核性与关系：核性是该节点在其父节点中是核心还是附属，关系是它与兄弟节点之间的关系\scite{nlp-analysis-marcu2000}。在此基础上逐级加严，得到四个指标：只比较跨度（S），跨度与核性都对（N），跨度与关系都对（R），三者都对（F）。EDU切分给定时，两棵树的节点数相同，精确率等于召回率。

用三个单元的例子计算。参考树的四个非根节点为$u_1$（S，原因）、$u_2$（N，原因）、$[u_1u_2]$（S，原因）、$u_3$（N，原因）。设预测树为$[u_1^{\mathrm S}[u_2^{\mathrm N}u_3^{\mathrm S}]_{\text{阐述}}^{\mathrm N}]_{\text{原因}}$，其非根节点为$u_1$（S，原因）、$[u_2u_3]$（N，原因）、$u_2$（N，阐述）、$u_3$（S，阐述）。跨度上，三个叶节点都匹配，$[u_1u_2]$与$[u_2u_3]$不匹配，S为$3/4$；核性上$u_1,u_2$匹配，N为$2/4$；关系上只有$u_1$匹配，R为$1/4$；F同样为$1/4$。预测树唯一的内部结构决策是错的，跨度指标却有0.75，原因是叶节点在EDU给定时总是匹配。Morey等据此指出，RST-Parseval因计入叶节点而抬高了分数，且不同论文混用不同的计算方式，使结果不可直接比较；他们建议改用只评价内部节点的标准Parseval，并统一评分脚本\scite{nlp-analysis-morey2017}。

\subsubsection{递归分割式篇章树恢复}
\label{sec:nlp-analysis-rst-split}

移进--归约在看到全文之前就要作出局部合并，较早的错误合并会一直保留在树中。自顶向下的方法反过来，先在整段上决定最高层的划分，再递归处理两侧。最高层的划分往往对应段落或语义大块之间的边界，可以利用全文信息作出判断\scite{nlp-analysis-lin2019unified,nlp-analysis-zhang2020topdown,nlp-analysis-kobayashi2020}。

设EDU边界区间$[i,j)$包含$j-i$个单元，若$j-i=1$直接返回叶节点；否则在$i<k<j$中选择分割点，再预测核性$\nu$与关系$r$，递归处理$[i,k)$与$[k,j)$。实现上，编码器先给出每个EDU的表示$\bm e_1,\ldots,\bm e_m$，区间表示$\bm z_{i,j}$可由区间内EDU表示汇聚得到（或用解码器的隐状态逐步维护）；分割头对候选边界做指针注意力：
\[
  p(k\mid i,j)=\operatorname{softmax}_{i<k<j}\bigl(\bm v^{\top}\tanh(\bm W_1\bm z_{i,j}+\bm W_2\bm e_k)\bigr),
\]
核性头与关系头在选定的分割下，以两个子区间的表示为输入分类。参考树的每个内部节点提供分割、核性、关系三项监督，各项交叉熵相加作为目标函数。区间按栈顺序依次展开，$m$个EDU共需$m-1$次分割决策。

对$[0,3)$，设分割点$k=1,2$的logit为1与3，选$k=2$，概率约$0.881$，核性SN与关系“原因”得分最高；子区间$[0,2)$只有$k=1$一种分割，仍取原因与SN；$[2,3)$等单元区间直接返回。恢复的树与第~\ref{sec:nlp-analysis-rst-transition}小节的参考树相同。若最高层误选了$k=1$，那么$[0,2)$这一跨度不可能在后续出现，这说明自顶向下方法同样有错误传播，只是方向相反：高层错误限制了所有下层结构。

如果分割分数只依赖当前区间与分割点，不依赖此前的分割历史，那么整棵树的分数是各内部节点分数之和，可以像第~\ref{sec:nlp-analysis-neural-cky}小节的神经CKY（式~\eqref{eq:nlp-analysis-neural-cky}）那样，以EDU为单位用区间动态规划求精确最优：
\[
  D(i,j)=\max_{i<k<j,\,r,\nu}\bigl\{s(i,j,k,r,\nu)+D(i,k)+D(k,j)\bigr\},\qquad D(i,i+1)=0,
\]
保存每个区间的最优分割与标签即可回溯，复杂度为$O(m^3)$。贪心分割只需$O(m)$次决策，但不保证全局最优；解码器若用隐状态记录了此前的分割历史，分数不再可分解，上式就不能直接使用。

表~\ref{tab:nlp-analysis-discourse-methods}对照本小节的方法。

\begin{table*}[htbp]
  \centering
  \small
  \begin{tabularx}{\linewidth}{@{}l XXX@{}}
    \toprule
    方法 & 解决的问题 & 核心机制 & 新的代价或局限\\
    \midrule
    边界标注与显式关系分类 & EDU切分；连接词消歧与关系分类 & B/I标注；连接词及其句法上下文特征 & 论元跨度定位仍难，尤其Arg1远离连接词时\\
    论元对编码器 & 无连接词的隐式关系 & 两论元交互注意力或句对编码，再分类 & 依赖常识；以显式样本扩充时分布不同\\
    RST移进--归约 & 覆盖全文的层次树 & 转移系统，归约动作带关系与核性 & 局部早期合并不可撤回；标签空间大\\
    递归分割 & 自底向上的局部合并错误 & 自顶向下指针选择分割点；可分解时用区间DP & 高层分割错误限制下层；精确解需分数无历史依赖\\
    \bottomrule
  \end{tabularx}
  \caption{篇章单元与关系分析方法对照。}
  \label{tab:nlp-analysis-discourse-methods}
\end{table*}

\subsection{论点、理由与论证结构}
\label{sec:nlp-analysis-argumentation}

篇章关系中的“解释”说明两段内容如何衔接，论证分析关心的则是一段文字是否被用来支持或反驳另一个主张。\term{论证分析}（argument mining）先定位论证单元并判断类型，再恢复单元之间有方向的支持与攻击关系。以下材料含三个单元：$a_1$“应增加自习座位”，$a_2$“晚间常有人找不到空位”，$a_3$“本月晚间到场人数持续上升”。在标注中，$a_1$是主张，$a_2$支持$a_1$，$a_3$支持$a_2$。论证结构描述作者的推理组织，不判断这些陈述是否属实，也不判断理由是否充分。

这一小节的数据主要来自Stab与Gurevych标注的议论文语料（第2版含402篇学生议论文），单元类型为主论点（MajorClaim）、论点（Claim）和前提（Premise），关系为支持与攻击，论点相对主论点另标赞成或反对的立场。

论证结构分析的挑战有三项。单元边界与类型：论证单元不一定与句子或EDU对齐，同一句可能包含主张与理由；类型（主论点、论点、前提）依赖它在整篇中的位置。关系方向：支持关系从理由指向主张，方向错误等于结构错误。全局结构：在Stab与Gurevych的议论文语料中，每段的论证结构是以主张为根的树或森林，局部预测的关系可能成环或给一个前提指向多个目标\scite{stab-gurevych-2017-parsing}。方法谱系有两条：流水线先分类单元与关系，再用整数线性规划（ILP）恢复全局合法的结构；双仿射依存解析把单元当作依存树的节点，用神经弧评分接同一套结构约束，并可与单元识别联合。

\subsubsection{论证单元分类与ILP结构恢复}
\label{sec:nlp-analysis-argument-ilp}

Stab与Gurevych的流水线分四步：用序列标注识别论证单元，分类单元类型，对单元对预测是否存在连接，最后用ILP联合协调前两类局部输出。ILP中边的权重综合了关系分类器的连接概率、由类型分类器给出的“该单元像主张”的程度，以及由其他边推出的入边与出边数。求解后，没有出边的单元为主张，其余为前提；支持或攻击的立场再单独分类\scite{stab-gurevych-2017-parsing}。

用二元变量$x_{i,j}\in\{0,1\}$表示是否选择从理由$i$到目标$j$的边，$w_{i,j}$为边的权重，目标是最大化所选边的总权重：
\[
  \max_{\bm x}\sum_{i\neq j}w_{i,j}x_{i,j}.
\]
基本约束为
\[
  x_{i,i}=0,\qquad \sum_jx_{i,j}\leq1\ \ (\forall i),\qquad \sum_{i,j}x_{i,j}\leq n-1,
\]
分别表示没有自环、每个单元至多指向一个目标、总边数不超过树的边数。这三项还不能排除环，例如$a_2\to a_3$与$a_3\to a_2$同时满足它们。为此引入层次变量$0\leq u_i\leq n-1$，要求
\[
  u_i\geq u_j+1-n(1-x_{i,j}),\qquad i\neq j.
\]
选中边$i\to j$时约束变为$u_i\geq u_j+1$，即理由的层次严格高于目标；沿一个环累加这些不等式会得到$u_i\geq u_i+k$（$k$为环长），矛盾，因此不能成环。未选中时约束右端至多为$u_j+1-n\leq0$，自动满足。这样得到的是有向森林；若标注要求每段恰好一棵树，再加上总边数等于$n-1$即可。

对三个单元，设主张$a_1$不能有出边，其余候选边的权重为$w_{2,3}=6$、$w_{3,2}=5$、$w_{2,1}=4$、$w_{3,1}=2$（教学构造）。各自取最高分的出边，得到$a_2\to a_3$与$a_3\to a_2$，总分11，却成环。合法组合中，$a_2\to a_1,a_3\to a_2$得9分，$a_2\to a_3,a_3\to a_1$得8分，ILP取9分的结构：$a_1$为根，$a_2,a_3$为前提。

ILP的优点是约束可以按标注规范自由增减，局部分类器可以任意替换；代价是求解时间在单元数多时上升，而且流水线中单元识别的错误无法由ILP修复。两个理由共同才能支持一个结论（合取前提）的情形，普通二元边无法表达，需要额外的论证节点。

\subsubsection{双仿射论证依存解析}
\label{sec:nlp-analysis-argument-biaffine}

论证结构是以主张为根、每个前提至多指向一个目标的树或森林，这与依存树“每个词至多一个中心词”的约束相同。因此可以把论证单元当作词，把“前提指向目标”当作“依存词指向中心词”，直接复用第~\ref{sec:nlp-analysis-biaffine}小节的双仿射依存解析\scite{nlp-analysis-eger2017}。每个单元先编码为向量$\bm u_i$（如对单元内词元表示做池化），再经两个多层感知机分别投影为“依据端”表示$\bm h_i^{(d)}$与“目标端”表示$\bm h_j^{(h)}$，弧分数按式~\eqref{eq:nlp-analysis-biaffine}计算为$s_{i,j}=\bm h_j^{(h)\top}\bm U\bm h_i^{(d)}+\bm h_j^{(h)\top}\bm u$。类型头预测主张、前提或非论证，关系头预测支持或攻击。训练对标注的弧、类型与关系分别计算交叉熵；解码时把$s_{i,j}$作为第~\ref{sec:nlp-analysis-argument-ilp}小节整数规划中的边权$w_{i,j}$，施加同一套森林约束，也可以用Chu--Liu/Edmonds算法求最大生成树。

令依据端$\bm h_2^{(d)}=(1,0)^{\top}$、$\bm h_3^{(d)}=(0,1)^{\top}$，目标端$\bm h_1^{(h)}=(4,2)^{\top}$、$\bm h_2^{(h)}=(0,5)^{\top}$、$\bm h_3^{(h)}=(6,0)^{\top}$，取$\bm U$为单位矩阵、$\bm u=\bm 0$，则$s_{2,1}=4$、$s_{3,1}=2$、$s_{3,2}=5$、$s_{2,3}=6$，恰好是该小节算例的边权，约束解码得到同一个9分结构。双仿射评分替换的是边权的来源，结构约束与解码问题不变。

与ILP流水线相比，双仿射解析让单元表示、弧评分和类型分类共享一个编码器联合训练，不必手工设计单元对特征。单元边界未知时，可以把词元级BIO标注与单元间的弧预测放在同一模型中联合学习\cite{nlp-analysis-eger2017}；Eger等比较了把端到端论证分析写成依存解析、序列标注与多任务学习等几种形式，发现在他们的议论文数据上，直接套用依存解析的效果较差，以BiLSTM为基础、把单元类型与关系距离编进标签的序列标注模型在不同设置下表现较稳健，多任务学习也有帮助；Ye与Teufel进一步把论证单元内部用带类型的连接串起来，以单元的末词作为对外端点，在词元层面做依存解析，再把端点间的弧映射回单元间的关系，这一做法允许非树的论证图\scite{ye-teufel-2021-end}。对于本身就是图（一个前提支持多个主张）的语料，Morio等去掉单父约束，对每对单元独立判断是否存在边\scite{nlp-analysis-morio2020}。约束放松后可以覆盖更多结构，但不再能用最大生成树保证输出合法，需要另外检查环。

\subsection{内容衔接与连贯性}
\label{sec:nlp-analysis-coherence}

前三个小节输出具体的边；连贯性分析转而评价整段文字：给段落一个连贯性分数，或者把打乱的句子恢复成最合理的顺序。比较三句材料：A“小王借了一本书”，B“他读完了这本书”，C“这本书被归还”。原顺序ABC读来自然；顺序ACB让归还先于阅读，所有实体词都还在，读者却感到不通。

连贯性的挑战首先是缺少直接标注：很少有语料给段落标连贯分数。常见做法是用原文作为正例，用打乱句序、删去指代、插入无关句等扰动构造负例，训练模型区分两者。这种监督有两个隐患：模型可能只学会识别某一种扰动，而且可能依赖句首编号、段落长度、标点等与连贯无关的捷径。方法谱系为：中心理论描述相邻句之间注意焦点如何转移；实体网格把这种转移变成可学习的特征；句对衔接评分加束搜索直接求句子排列；指针网络以已选句子的全部历史为条件逐句选择。句子排序的结果用Kendall τ与完全匹配率评价。

\subsubsection{中心理论与实体网格}
\label{sec:nlp-analysis-entity-grid}

\term{中心理论}（centering theory）描述局部连贯：相邻两句如果围绕同一个实体展开，并且该实体保持在显著的语法位置，读者处理起来最省力\scite{nlp-analysis-grosz1995}。对每个话语单元$U_n$（这里取句子）定义三个对象。前瞻中心列表$C_f(U_n)$是$U_n$中出现的实体，按显著性排序，通常主语高于宾语，宾语高于其他成分。回指中心$C_b(U_n)$是$C_f(U_{n-1})$中在$U_n$里再次出现的、排位最高的实体，即这句话“接着谈论”的对象。首选中心$C_p(U_n)$是$C_f(U_n)$的第一个元素，预示下一句最可能继续谈论的对象。

相邻句之间的过渡按两个问题分类：回指中心是否与上一句相同，回指中心是否就是本句的首选中心。两者都是为延续（Continue），前者是而后者否为保持（Retain），前者否为转移（Shift）。中心理论还给出一条关于代词的规则：若$C_f(U_{n-1})$中有实体在$U_n$里用代词实现，则$C_b(U_n)$也必须用代词实现。它解释了为什么“小王借了一本书。小王读完了它”显得别扭：书用了代词，排位更高、作为回指中心的小王却用了全名。Brennan等又把转移细分为$C_b=C_p$与$C_b\neq C_p$两种，并提出偏好次序：延续优于保持，保持优于前一种转移，前一种转移优于后一种\scite{nlp-analysis-brennan1987}；后来的文献把这两种转移分别称为平滑转移（smooth shift）与粗糙转移（rough shift）。

对ABC，$C_f(A)=[\text{小王},\text{书}]$；B中两者都再现，$C_b(B)=\text{小王}=C_p(B)$，A没有回指中心，按惯例记为延续；C中只出现书，$C_b(C)=\text{书}\neq C_b(B)$，且$C_p(C)=\text{书}$，为平滑转移。对ACB，A到C为延续，C到B时$C_b(B)=\text{书}=C_b(C)$，但$C_p(B)=\text{小王}$，为保持。按偏好次序，ACB的过渡序列（延续、保持）反而优于ABC（延续、平滑转移）。这个例子说明中心理论只刻画实体焦点的延续，不刻画事件的时间顺序；ABC优于ACB的理由在于“先读后还”，这属于事件知识。

\term{实体网格}（entity grid）把中心理论的思想变成统计特征\scite{barzilay-lapata-2008-modeling}。网格以句子为行、共指实体为列，格值为S（主语）、O（宾语）、X（其他成分）或--（不出现）；同一句多次出现同一实体时，按S优先于O、O优先于X取值。每列从上到下的相邻格值构成长度为2的转移，如SS、S--、OS，共16种；统计全文各类转移所占的比例，得到特征向量$\bm\phi(D)\in\mathbb R^{16}$。与中心理论不同，网格不规定哪种过渡更好，而是从数据中学习每种转移的权重。

对三句材料，把“他”并入小王、“这本书”并入书，网格为
\[
  \begin{array}{c|cc}
     & \text{小王} & \text{书}\\ \hline
   A & \mathrm S & \mathrm O\\
   B & \mathrm S & \mathrm O\\
   C & - & \mathrm S
  \end{array}.
\]
C是被动句，书在语义上是受事，句法上是主语，格值取S。两列各有两次转移，共4次，SS、S--、OO、OS各占$1/4$。

学习用成对排序：原文$D$应比它的扰动版本$D'$得分高，设$f(D)=\bm w^{\top}\bm\phi(D)$，排序SVM的目标函数为
\[
  \min_{\bm w}\ \frac{\mu}{2}\lVert\bm w\rVert_2^2+\sum_{(D,D')}\max\bigl(0,1-f(D)+f(D')\bigr),
\]
其中$\mu$为正则化系数。设教学权重$w_{\mathrm{SS}}=2$、$w_{\mathrm{OO}}=1$、其余为0，原顺序得分$0.75$；ACB产生S--、--S、OS、SO四种转移，得分0，分差$0.75$，铰链损失为$0.25$。网格依赖共指结果：若漏掉“他”与小王的共指，网格变成三列，SS消失，OO只占$1/6$，原文得分降为$1/6$，共指错误由此传入连贯性判断。实体网格只捕捉实体延续，一段始终围绕同一实体、却前后矛盾的文字，网格无法区分。Nguyen与Joty把网格的离散特征换成可学习的表示：每个格值S、O、X、--映射为嵌入向量，在每一列（即每个实体的角色序列）上做一维卷积以捕捉更长的转移模式，池化后打分，仍以原文优于扰动的成对排序损失训练\scite{nlp-analysis-nguyen-joty2017}。卷积窗口覆盖多句，弥补了长度为2的转移只看相邻句的局限，但输入仍是实体网格，共指错误同样会传入。

\subsubsection{句对衔接评分与束搜索排序}
\label{sec:nlp-analysis-pair-ordering}

实体网格给整篇打分，要用它恢复句序就得枚举排列。另一种思路直接学习“哪句适合接在哪句后面”：训练有方向的句对评分$s(a,b)$，正样本来自原文中的相邻句，负样本来自不相邻句或颠倒的句对\scite{nlp-analysis-chen2016ordering}。再学习起始分数$s_0(a)$，表示句子$a$适合作为开头。对排列$\pi=(\pi_1,\ldots,\pi_m)$，总分为
\[
  S(\pi)=s_0(\pi_1)+\sum_{t=2}^{m}s(\pi_{t-1},\pi_t).
\]
求最高分排列等价于在句子构成的完全图上找一条最大权哈密顿路径，一般是NP难的，实践中用束搜索近似（第~\ref{sec:nlp-generation-beam}小节），搜索状态为“已选序列及剩余句子集合”。

取$s_0(A)=2$、$s_0(B)=0$、$s_0(C)=-1$，$s(A,B)=3$、$s(A,C)=4$、$s(B,C)=2$、$s(B,A)=-1$、$s(C,A)=-3$、$s(C,B)=-2$（教学构造）。束宽为2时，第一步保留A（2分）与B（0分）；扩展后AC得6、AB得5、BC得2、BA得$-1$，保留AC与AB；第三步各自只能接剩下的句子，ACB得$6-2=4$，ABC得$5+2=7$，输出ABC。束宽为1时，第二步只保留AC，最终得到ACB。穷举全部6个排列可确认7分是最大值。这个例子说明，句对评分只看相邻两句，局部最高分的选择不保证整条排列最优；束宽越大越接近最优，计算量也随之增加。

句对评分还有另一种用法：不判断“是否相邻”，而判断“$a$是否在$b$之前”。Prabhumoye等用BERT对每一对句子分类先后，得到一张有向图，再对它做拓扑排序得到全序\scite{nlp-analysis-prabhumoye2020}。先后关系比相邻关系更容易从内容判断（例如“读完”应早于“归还”，中间可以隔着别的句子），而且拓扑排序是线性时间的精确算法；代价是成对分类可能互相矛盾而形成环，需要先按置信度删边，并且$O(m^2)$个句对都要过一次编码器。

\subsubsection{指针式句子排序}
\label{sec:nlp-analysis-pointer-ordering}

句对评分只依赖前一句，“B接在A后面合不合适”与A之前已经说过什么无关。指针网络让每一步的选择依赖已选的全部历史\scite{nlp-analysis-gong2016pointer}。编码器先把每个句子编码为$\bm h_j$，再对全部句子做整体编码（如以句子表示为输入的LSTM或自注意力层），得到段落表示作为解码器初始状态；解码器把已选句子依次输入，维护状态$\bm z_t$，在第$t$步对每个剩余句子打分：
\[
  g_{t,j}=\bm v^{\top}\tanh(\bm W\bm h_j+\bm U\bm z_t),\qquad
  p(\pi_t=j\mid\pi_{<t})=\operatorname{softmax}_{j\notin\pi_{<t}}(g_{t,j}),
\]
已选句子的分数设为$-\infty$。训练时把参考顺序作为解码输入（教师强制），对每一步取交叉熵。

设第一步对A、B、C的概率为$0.6,0.3,0.1$，选A；第二步屏蔽A，对B、C的概率为$0.75,0.25$，选B；第三步只剩C。ABC的路径概率为$0.6\times0.75=0.45$，训练损失$-\log0.45\approx0.799$；ACB的路径概率为$0.6\times0.25=0.15$。训练时解码器总是看到正确的历史，推断时则看到自己的预测，首句选错会影响此后每一步的分数，这是第~\ref{sec:nlp-analysis-paradigm-transition}小节所说的暴露偏差在句子排序中的表现。

\paragraph{排序评价}
排序结果用两个指标评价。Kendall τ衡量预测排列与参考排列的一致程度\scite{nlp-analysis-lapata2006}。设$m$个句子有$\binom m2$个句子对，其中在预测中相对次序与参考相反的对数为$I$（逆序对数），则
\begin{equation}
  \tau=1-\frac{2I}{\binom m2},
  \label{eq:nlp-analysis-kendall-tau}
\end{equation}
取值在$-1$到1之间，完全一致为1，完全颠倒为$-1$。完全匹配率（perfect match ratio, PMR）是预测排列与参考完全一致的文档所占比例。对ACB，三个句子对中只有(B,C)颠倒，$I=1$，$\tau=1-2/3=1/3$。若测试集有两篇文档，一篇完全恢复，一篇输出ACB，则平均τ为$(1+1/3)/2\approx0.667$，PMR为$1/2$。τ对局部交换给部分分，PMR只认完全正确，两者结合可以区分“大体对、局部错”和“完全错”。两个指标都假设参考顺序唯一，存在多种合理顺序时会低估模型。

表~\ref{tab:nlp-analysis-coherence-methods}汇总本小节的方法。

\begin{table*}[htbp]
  \centering
  \small
  \begin{tabularx}{\linewidth}{@{}l XXX@{}}
    \toprule
    方法 & 解决的问题 & 核心机制 & 新的代价或局限\\
    \midrule
    中心理论 & 描述局部连贯 & 回指中心与首选中心的过渡类型及偏好次序 & 不可学习；不刻画事件顺序\\
    实体网格 & 把实体转移变为可学习特征 & 角色转移比例特征加成对排序损失 & 依赖共指与句法；只看实体延续\\
    句对评分加束搜索 & 直接恢复句序 & 有方向的相邻句评分，束搜索求排列 & 只依赖前一句；搜索近似\\
    指针网络 & 选择依赖全部历史 & 编码全部句子，逐步指针选择剩余句子 & 暴露偏差；长文档成本高\\
    \bottomrule
  \end{tabularx}
  \caption{连贯性与句子排序方法对照。}
  \label{tab:nlp-analysis-coherence-methods}
\end{table*}

\section{语用分析}
\label{sec:nlp-analysis-pragmatics}

“明天能完成吗”在制定计划时是在询问可行性；对方已经两次延误之后，同一句话更像催促。字面命题几乎相同，在交流中起的作用却不同。\term{语用分析}（pragmatic analysis）研究话语在具体背景下的作用：它在交流中完成了什么行为，省略了哪些可以从背景补出的内容，以及说话者通过字面意义暗示了什么。这里的背景包括说话者与听话者、当前话语回应的对象、此前的对话历史，以及双方共享的常识。

本节的三个小节对应三类问题。话语行为把每一轮话语归入请求、告知、确认等功能类别，是对话系统理解用户与规划回复的基础；省略恢复与话语改写把依赖历史才能理解的话语补全为独立完整的句子，使下游的检索与问答模块可以只看当前句；会话含义与间接言语行为则处理字面意义与说话者意图不一致的情形，这是语用中最依赖推理的部分。三者的共同难点是，所需的信息不在当前句之内，模型必须从背景中取用，又不能补出背景没有支持的内容。

\subsection{话语行为}
\label{sec:nlp-analysis-dialogue-acts}

\term{话语行为}（dialogue act）用请求、告知、确认、承诺等标签表示一次话语在交流中的功能。这一概念来自言语行为理论（见第~\ref{sec:nlp-analysis-indirect-acts}小节）：说话不只是陈述命题，也是在做事，如提问、许诺、请求。话语行为标注把这种思想落实为可操作的标签集。Switchboard对话行为语料（SwDA）在1155段电话交谈上按SWBD-DAMSL规范标注，常用的归并标签集有42类，包括陈述观点、陈述非观点、附和（backchannel，如“嗯”“对”）、是非问、同意等\cite{nlp-analysis-jurafsky1997swbd}；ICSI会议对话行为语料（MRDA）在75场多人会议上标注，原始标签由11个一般标签和39个具体标签组合而成\scite{nlp-analysis-shriberg2004}；后续研究最常用的是随语料发布的五类映射：陈述、提问、话轮控制（floor mechanism，如抢话与占话）、附和与打断。两份语料的单位都是切分后的话语片段，而不是完整话轮：一个话轮可以包含“嗯，对，我觉得周五不行”这样连续的几个片段。多功能问题在标注规范层面也有专门的处理：国际标准ISO 24617-2把话语功能分为任务、反馈、话轮管理、时间管理等多个维度，一个片段可以在每个维度上各有一个功能，例如“嗯，周五不行”同时承担反馈维度的确认收听与任务维度的告知\scite{nlp-analysis-bunt2012iso}。采用这类规范时，识别就成为每个维度各一个分类器的多任务问题。

话语行为识别有四个挑战。同一句式在不同前文下功能不同，“明天能完成吗”既可以是询问也可以是催促。说话者转换：对话行为的依赖往往跨越说话者，问句之后通常是回答，而回答出自另一方。多功能话语：“谢谢，能再发一份吗”同时致谢与请求，标签集是否允许多标签要事先规定。标签集差异：SwDA与MRDA的类别粒度、单位切分都不同，模型和结论不能直接迁移。此外，在线识别只能使用当前及以前的话轮，事后分析可以看到整段对话，两种设置的输入条件不同，结果不可直接比较。

方法谱系按怎样使用上下文展开。单句分类只看当前话语，用词语、句式、说话者等特征送入最大熵、SVM或卷积网络分类器，这是最直接的基线，但它对上面第一项挑战无能为力。跨轮HMM与CRF把行为序列看作马尔可夫链，用行为之间的转移刻画前后依赖。层次对话编码先在话轮内汇聚词语得到话轮向量，再在话轮之间编码对话历史，用学到的表示代替人工特征，并保留CRF对标签依赖的建模。

单句分类器也可以用来扩展标签集。新领域出现训练数据中没有的行为类别时，可以把每个标签的定义与几个例子写成一段描述，用句对编码器计算“当前话语--标签描述”的匹配分数，在所有描述中取最高者。这改变的是类别信息的来源：模型不再从大量样本中学习一个类别，而是从描述中读出它的边界，因此描述必须写清相近类别的区别，例如“询问可行性”与“要求加快进度”。训练与预测时可见的信息也应一致：若标注者看了后续回答才判断当前行为，而系统在预测时只能看到当前及以前的话语，监督信号就含有预测时不可得的信息，这种任务应明确为事后分析。

\subsubsection{跨轮HMM与CRF话语行为标注}
\label{sec:nlp-analysis-dialogue-sequence}

Stolcke等把一段对话的行为序列建模为HMM：每轮行为$a_t$是隐状态，话语$u_t$是观测，联合概率为$\prod_t p(a_t\mid a_{t-1})\,p(u_t\mid a_t)$\scite{stolcke-etal-2000-dialogue}。转移概率由已标注的相邻行为计数估计，相当于一个行为层面的$n$元语言模型；发射概率由每类行为各自的词语$n$元语言模型给出，还可以结合韵律特征。若手头只有判别式分类器给出的$p(a_t\mid u_t)$，可以按贝叶斯公式除以先验$p(a_t)$，换算为与$p(u_t\mid a_t)$成比例的量后再代入，不能直接把后验当作发射概率使用。

固定观测“明天能完成吗”，设教学发射概率为$p(u\mid I)=0.6$、$p(u\mid U)=0.4$，I为询问，U为催促。前轮行为是计划讨论P时，转移到I、U的概率为$0.8,0.2$，乘积为$0.48,0.08$，归一化后I占$6/7$；前轮是延误解释E时，转移概率为$0.2,0.8$，乘积为$0.12,0.32$，U占$8/11$。观测的贡献没有变，改变排序的是行为上下文。

前轮行为本身也不确定时，在线识别用前向算法逐轮更新过滤分布：
\[
  p(a_t\mid u_{\leq t})\propto p(u_t\mid a_t)\sum_{a_{t-1}}p(a_t\mid a_{t-1})\,p(a_{t-1}\mid u_{<t}).
\]
事后分析整段对话时，前向--后向算法给出每轮的边缘分布，Viterbi算法给出整条最可能的行为序列。两者回答不同的问题：前者对前一轮的所有可能行为求和，后者只沿最好的一条路径取最大，因此Viterbi路径上某一轮的行为，不一定是该轮边缘概率最大的行为。看一个两轮的例子，第二轮仍是“明天能完成吗”，沿用上面的转移与发射概率。若第一轮的过滤分布为P占$0.55$、E占$0.45$，第二轮I的未归一化概率为$0.55\times0.48+0.45\times0.12=0.318$，U为$0.55\times0.08+0.45\times0.32=0.188$，边缘偏向I；Viterbi中到达I的最优路径来自P，分数$0.264$，到达U的来自E，分数$0.144$，最优路径P--I，两者一致。若第一轮改为P占$0.35$、E占$0.65$，第二轮I的未归一化概率为$0.35\times0.48+0.65\times0.12=0.246$，U为$0.35\times0.08+0.65\times0.32=0.236$，边缘仍略偏向I（约$0.51$）；Viterbi比较的却是$0.35\times0.48=0.168$与$0.65\times0.32=0.208$，最优路径为E--U。I的概率来自P与E两条前驱路径之和，单看任何一条都不如E--U高。需要给每一轮输出最可能的标签时用边缘分布，需要一条前后一致的行为序列时用Viterbi。HMM的定义与前向后向、Viterbi递推见第~\ref{chap:nlp-classical-models}章。

线性链CRF直接对条件概率建模，可以把话语文字、说话者是否转换、与前一轮的词语重叠等任意特征放进位置得分，并学习行为之间的转移分数，其定义、梯度与解码见第~\ref{sec:nlp-classical-crf}小节。对同一问句设I、U的位置分数为$1,0.5$；P之后的转移分数为$1,0$，两候选总分为$2,0.5$，偏向I；E之后的转移分数为$-0.5,1$，总分为$0.5,1.5$，偏向U。与分词的BMES标注不同，话语行为之间一般没有必须禁止的转移，因此不需要第~\ref{sec:nlp-analysis-paradigm-tagging}小节那样的合法转移掩码，转移矩阵完全由数据学习。

HMM与CRF的局限在于，上下文只通过行为标签传递：两段历史若行为标签序列相同而内容不同，模型无法区分。“明天能完成吗”之前若是“上次你说周三能好”，催促的意味来自具体内容，而这一内容在“告知”这个标签中已经丢失。

\subsubsection{层次对话编码与分层BiLSTM-CRF}
\label{sec:nlp-analysis-dialogue-hierarchical}

层次对话编码让历史话语的内容而不仅是标签进入当前判断。对话天然有两层结构：词组成话语，话语组成对话。Kumar等的分层BiLSTM-CRF按这两层各用一个双向LSTM编码，再接线性链CRF解码整段行为序列\cite{nlp-analysis-kumar2018}，结构如图~\ref{fig:nlp-analysis-hierarchical-dialogue-act}所示。

设一段对话$C$含$R$个话语$u_1,\ldots,u_R$，第$j$个话语有$S_j$个词。第一层是词级编码：每个词经嵌入层得到$\bm e_k\in\mathbb R^{d_e}$，词级BiLSTM输出每个位置的前后向隐状态拼接$\bm h_k\in\mathbb R^{2h_1}$。话语向量取最后一个时间步的表示，$\bm v_j=\bm h_{S_j}\in\mathbb R^{2h_1}$；原文也比较了平均池化，最终采用末状态。第二层是话轮级编码：把$\bm v_1,\ldots,\bm v_R$当作一个序列送入另一个BiLSTM，输出$\bm g_j\in\mathbb R^{2h_2}$，它同时包含$u_j$之前与之后的话语内容。第三层是CRF：设行为标签集大小为$K$，位置分数$\bm o_j=\bm W\bm g_j+\bm b\in\mathbb R^{K}$，转移矩阵$\bm A\in\mathbb R^{K\times K}$与位置无关，行为序列$\bm y$的条件概率为
\begin{equation}
  p(\bm y\mid C)=\frac{\exp\bigl(\sum_{j=1}^{R}(\bm o_{j,y_j}+\bm A_{y_{j-1},y_j})\bigr)}
   {\sum_{\bm y'}\exp\bigl(\sum_{j=1}^{R}(\bm o_{j,y'_j}+\bm A_{y'_{j-1},y'_j})\bigr)}.
  \label{eq:nlp-analysis-hierarchical-crf}
\end{equation}
训练最大化各段对话的对数似然$\sum_C\log p(\bm y\mid C)$，嵌入层、两层BiLSTM与CRF参数端到端联合更新；推断用Viterbi算法求最优行为序列。用两轮对话看CRF层的作用（数值为教学构造）。设标签只有问（Q）与答（A）两类，两轮话语的位置分数为$\bm o_1=(2,0)$、$\bm o_2=(0.4,0.6)$，即第二轮单看自身略偏向答；转移矩阵中$\bm A_{Q,A}=1$、$\bm A_{Q,Q}=-1$、$\bm A_{A,Q}=\bm A_{A,A}=0$，起始转移取0。四条序列的总分为QQ：$2+0.4-1=1.4$，QA：$2+0.6+1=3.6$，AQ：$0+0.4=0.4$，AA：$0+0.6=0.6$。最优序列QA的概率为$e^{3.6}/(e^{1.4}+e^{3.6}+e^{0.4}+e^{0.6})\approx0.83$。若去掉转移矩阵、逐轮独立取最大，第二轮也会选A，但概率只有约$0.55$；转移矩阵把“问之后多为答”的规律加了进来，使判断更确定。若第二轮的位置分数明显偏向Q，例如$(2,0)$，QQ得$3$，QA得$3$，两者持平，此时是否连续两问取决于话语内容与转移规律的权衡。按张量形状看，一个批次中的对话先被展开为话语的批次，形状为（话语数，最大词数，$d_e$），词级编码后得到（话语数，$2h_1$），再按对话重新组织为（对话数，最大话语数，$2h_1$）送入话轮级编码。

\begin{figure*}[htbp]
  \centering
  % 分层BiLSTM-CRF话语行为标注：词级编码汇聚为话轮向量，话轮级编码建模轮间依赖，CRF解码行为序列。
\begin{tikzpicture}[
  x=1mm,y=1mm,figure picture,
  cell/.style={figure node,minimum width=12mm,minimum height=6mm,
    inner sep=1pt,font=\figurefont\scriptsize},
  word/.style={font=\figurefont\scriptsize},
  row label/.style={font=\figurefont\scriptsize,text=FigureMuted,anchor=east},
  flow/.style={figure flow},
  link/.style={draw=FigureStroke,line width=.7pt,{Stealth[length=1.3mm]}-{Stealth[length=1.3mm]}},
  pair/.style={draw=FigureOutput,line width=.9pt}
]
  \foreach \j/\x/\txt in {1/0/{计划定在周五},2/36/{明天能完成吗},3/72/{可以}}{
    \node[word] (u\j) at (\x,0) {$u_{\j}$：\txt};
    \node[cell,draw=FigureInput,minimum width=26mm] (w\j) at (\x,10)
      {词级BiLSTM，取末状态};
    \node[cell,draw=FigureInput] (v\j) at (\x,21) {$\bm v_{\j}$};
    \node[cell,draw=FigureModel] (g\j) at (\x,33) {$\bm g_{\j}$};
    \node[cell,draw=FigureOutput] (y\j) at (\x,46) {$y_{\j}$};
    \draw[flow] (u\j) -- (w\j);
    \draw[flow] (w\j) -- (v\j);
    \draw[flow] (v\j) -- (g\j);
    \draw[flow] (g\j) -- (y\j);
  }
  \draw[link] (g1) -- (g2);
  \draw[link] (g2) -- (g3);
  \draw[pair] (y1) -- node[above,font=\figurefont\scriptsize] {$\bm A_{y_1,y_2}$} (y2);
  \draw[pair] (y2) -- node[above,font=\figurefont\scriptsize] {$\bm A_{y_2,y_3}$} (y3);
  \node[row label] at (-16,21) {话轮向量};
  \node[row label] at (-16,33) {话轮级BiLSTM};
  \node[row label] at (-16,46) {线性链CRF};
\end{tikzpicture}

  \caption{分层BiLSTM-CRF话语行为标注。每个话语先经词级BiLSTM编码，取末状态作为话语向量$\bm v_j$；话轮级BiLSTM在话语向量序列上编码得到$\bm g_j$，使每个位置都含有前后话语的内容；线性链CRF以$\bm g_j$给出位置分数，以转移矩阵$\bm A$给出相邻行为的分数，联合解码整段行为序列。}
  \label{fig:nlp-analysis-hierarchical-dialogue-act}
\end{figure*}

两层编码各自承担一项挑战。词级编码把话语内容压缩为向量，解决了HMM与CRF只能传递标签的问题；话轮级编码让前后话语的内容互相影响，CRF的转移矩阵再刻画“问句之后多为回答”这类标签依赖。原论文在SwDA（42类）与MRDA（5类）上评测，准确率分别为79.2\%与90.9\%，均高于此前报告的结果；训练集分别有1003段与51段对话。两份语料的标签分布都极不均衡，SwDA中陈述非观点与附和两类就占一半以上，MRDA中陈述类也超过一半，因此只看准确率容易高估模型对少数类的识别能力，宜同时报告宏平均F1。

这一模型有两点需要按任务调整。其一，话轮级BiLSTM的后向部分读取了未来话语，只适用于事后分析；在线识别必须改用单向编码，或在每一轮只对已发生的前缀重新编码。其二，原模型没有区分说话者，而说话者是否转换对行为预测很有信息量（同一人连续两句多为陈述的延续，换人之后常是回答或附和）。常见的扩展是在$\bm v_j$上拼接说话者嵌入，或拼接“与上一话语是否同一说话者”的指示特征；在多方对话中，当前话语回应的未必是紧邻的上一句，还需要输入或同时预测回应关系。多功能话语若规范允许多个标签，应把CRF换成逐标签的sigmoid输出，或先切分片段再分别标注。

后续工作主要替换编码器。Raheja与Tetreault在话语层与对话层引入自注意力，让话语向量按注意力汇聚词语表示，而不是只取末状态，并保留CRF解码\cite{nlp-analysis-raheja2019}；用预训练编码器替换词级BiLSTM时，把当前话语与前几轮拼接后编码，可以同时覆盖词级与话轮级两层。这些替换改变的是表示能力，话语—对话的两层结构与CRF的标签依赖仍然沿用。

表~\ref{tab:nlp-analysis-dialogue-act-methods}对照三类方法。

\begin{table*}[htbp]
  \centering
  \small
  \begin{tabularx}{\linewidth}{@{}l XXX@{}}
    \toprule
    方法 & 解决的问题 & 核心机制 & 新的代价或局限\\
    \midrule
    单句分类 & 话语功能的基本识别 & 词语、句式、说话者特征加分类器 & 不用上下文，同一句式只能有一个答案\\
    HMM与CRF序列标注 & 前后话语行为的依赖 & 行为转移概率或转移分数，前向或Viterbi解码 & 上下文只经标签传递，丢失历史内容\\
    分层BiLSTM-CRF & 历史内容进入当前判断 & 词级与话轮级两层编码，CRF建模标签依赖 & 双向编码只适合事后分析；需另加说话者信息\\
    \bottomrule
  \end{tabularx}
  \caption{话语行为识别方法对照。}
  \label{tab:nlp-analysis-dialogue-act-methods}
\end{table*}

\subsection{省略恢复与话语改写}
\label{sec:nlp-analysis-utterance-rewriting}

多轮对话中，用户经常省略前文已经说过的内容。历史H1为“今天能到图书馆借书吗”，历史H2为“今天的空气质量如何”，当前话语都是“那明天呢”。在H1之下，它的意思是“明天能到图书馆借书吗”；在H2之下，是“明天的空气质量如何”。检索系统、问答模型若只拿到“那明天呢”，几乎无法工作。\term{话语改写}（utterance rewriting）把当前话语改写为脱离历史也能理解的完整句子，补出被省略的成分（\term{省略恢复}，ellipsis recovery），并把代词替换为它所指的内容。改写之后，下游模块可以把对话问题当作单轮问题处理。

这一任务有几份常用数据。CANARD由Elgohary等构建：从对话式问答数据QuAC中取出依赖上下文的问题，请标注者把每个问题改写为独立完整的问题，共约4万条\scite{nlp-analysis-elgohary2019}。TREC CAsT（对话式辅助检索评测）面向对话式段落检索，每个话题是一串相关的问题，后面的问题大量省略和指代前面的内容，评测方同时提供人工改写的独立问题，便于比较“用原始问题检索”与“用改写问题检索”的差距\scite{nlp-analysis-dalton2020}。中文方面，Pan等构建了Restoration-200K，在约20万条开放领域多轮闲聊上标注了补全后的话语\cite{nlp-analysis-pan2019restoration}；Su等收集了2万段中文多轮对话并为其中的话语标注改写，他们对对话数据的分析发现，约七成的话语存在指代或省略\scite{nlp-analysis-su2019rewriter}。

省略恢复的挑战有三项。其一，省略的位置不显式：当前话语中没有占位符告诉模型缺了什么、缺在哪里。其二，补全内容必须来自历史：“明天能到图书馆借书吗”中的“图书馆”“借书”都应能在H1中找到来源，模型补出历史中没有的内容，就改变了用户的问题。其三，历史不足时应保留未知：若对话没有历史，“那明天呢”就没有可靠的补全，应当保持原样或请求澄清。评价也因此有两个层面：改写与参考改写的重合度（BLEU、ROUGE、精确匹配，或只看补入词的精确率与召回率），以及改写对下游检索或问答效果的提升。重合度指标中，只看补入词的精确率与召回率最能反映省略恢复本身：设参考改写相对当前话语新增的词集合为$A$，系统改写新增的词集合为$B$，则精确率为$|A\cap B|/|B|$，召回率为$|A\cap B|/|A|$。“那明天呢”在H1下的参考改写新增“能、到、图书馆、借书、吗”，系统若改写为“明天图书馆开门吗”，新增“图书馆、开门、吗”，精确率为$2/3$，召回率为$2/5$。BLEU在这里不够敏感：改写与当前话语大量重叠，照抄当前话语也能拿到不低的分数。下游评价则直接用改写去检索，在TREC CAsT上常报告前3个结果的NDCG@3，并与人工改写的检索效果比较。

方法谱系从结构化场景走向开放领域。改写并不是利用历史的唯一方式。也可以把完整历史与当前话语一起交给下游模型，让它在内部完成补全。显式改写的好处是中间结果可读、可检查，可以接入只接受单轮输入的现成检索系统，改写模型与下游模块可以分别优化；代价是多了一个模块，改写的错误会传给下游，且改写可能丢掉历史中对下游有用、却不属于省略内容的信息。两种方式常结合使用：检索用改写后的问题，生成回答时仍提供完整历史。

候选承接与槽位继承依赖预定义的槽位结构，把上一轮的槽值按规则或分类器继承到当前轮，可靠但覆盖有限；复制式改写用序列到序列模型从历史和当前句两路复制词语，不需要槽位定义，覆盖开放领域；大语言模型改写用提示直接完成，不需要专门的训练数据，但可能补入历史中没有的信息，需要来源检查。

\subsubsection{候选承接与槽位继承}
\label{sec:nlp-analysis-slot-carryover}

在任务型对话中，用户的意图由一组槽位描述，例如查询活动安排时有“活动”“日期”“时段”“时间”“地点”等槽位。每一轮结束后，系统维护一个对话状态，即目前已确定的槽值集合。省略在这种场景下表现为当前话语只给出部分槽值，其余槽值需要从对话状态继承。\term{槽位继承}（slot carryover）把这一判断写成分类问题：对历史中出现过的每个候选槽值$(s,v)$，判断它在当前轮是否仍然有效\scite{nlp-analysis-naik2018carryover}。

设当前话语为$u_t$，候选槽值$(s,v)$最近一次出现在$\delta$轮之前。编码器分别给出当前话语的表示$\bm q_t$、候选槽名与槽值的表示$\bm c_{s,v}$，再加上距离嵌入$\phi(\delta)$，以sigmoid输出继承概率：
\[
  p(\mathrm{carry}=1\mid u_t,s,v)=\sigma\bigl(\bm w^{\top}\operatorname{FFNN}([\bm q_t;\bm c_{s,v};\bm q_t\circ\bm c_{s,v};\phi(\delta)])\bigr),
\]
以标注的继承标签做二元交叉熵训练。例如当前话语“那上午呢”与候选｛时段：下午｝的匹配表示显示两者属于同一槽且取值冲突，网络会给出很低的继承概率；对候选｛活动：读书会｝，当前话语没有提及任何活动，距离$\delta=1$，继承概率高。当前话语若显式给出了某个槽的新值，该槽的旧值直接被替换，不再参与判断。规则方法是这一模型的特例：“当前轮未提及的槽一律继承”相当于对所有候选输出1，它在简单场景下够用，但无法处理槽之间的依赖。

依赖关系最常出现在时间类槽位上。设读书会的通知为“读书会改到周五下午两点”，对话状态为｛活动：读书会，日期：周五，时段：下午，时间：14:00｝。用户接着问“那地点呢”：只是换了询问对象，四个槽全部继承。用户问“改到周六呢”：日期被替换为周六，而“下午两点”是否随之保留，取决于用户是在问“周六同一时间”还是在问“周六有没有安排”，分类器应给出中等的继承概率，系统可以确认。用户问“那上午呢”：日期继承为周五，时段被替换为上午，精确时间14:00与上午矛盾，必须丢弃而不能继承。这说明继承判断不能逐槽独立进行，替换一个粗粒度的槽（时段）会使依附于它的细粒度槽（时间）失效。

槽位继承的长处是输出结构明确、可以逐槽检查，代价是只能处理槽位体系覆盖的内容。“那明天呢”在借书场景下需要补出“到图书馆借书”这一动作，在开放领域对话中省略的可能是任意短语，没有预先定义的槽位可以承接。

\subsubsection{基于复制的话语改写}
\label{sec:nlp-analysis-copy-rewriting}

开放领域的改写需要从历史中取出任意片段。观察可知，改写结果中的词几乎都来自两个地方：要么来自当前话语本身（“那明天呢”中的“明天”），要么来自历史（“图书馆”“借书”）。Su等据此设计了只做复制、不从词表生成的改写模型：解码器每一步从历史或当前话语中各自复制一个词，再由门控决定取哪一路\scite{nlp-analysis-su2019rewriter}。复制机制的一般形式见式~\eqref{eq:nlp-analysis-copy}，这里是它的双源版本。

编码器是Transformer，历史H与当前话语U拼接后共同编码，每个词元除词嵌入与位置嵌入外，还加上一个轮次嵌入表示它属于哪一轮，得到编码$\bm E\in\mathbb R^{(|H|+|U|)\times d}$，再按来源拆成$\bm E_H\in\mathbb R^{|H|\times d}$与$\bm E_U\in\mathbb R^{|U|\times d}$。解码器每层在自注意力之后，对两部分分别做编码器--解码器注意力，得到$\bm c^H_t$与$\bm c^U_t$，拼接后送入前馈层。最后一层在第$t$步对历史词元和当前词元各给出一个注意力分布$\bm\alpha_t\in\mathbb R^{|H|}$与$\bm\beta_t\in\mathbb R^{|U|}$，输出分布为
\begin{equation}
  p(r_t=w\mid H,U,r_{<t})=\lambda_t\sum_{i:\,H_i=w}\alpha_{t,i}+(1-\lambda_t)\sum_{j:\,U_j=w}\beta_{t,j},
  \label{eq:nlp-analysis-two-source-copy}
\end{equation}
\begin{equation}
  \lambda_t=\sigma\bigl(\bm w_d^{\top}\bm d_t+\bm w_H^{\top}\bm c^H_t+\bm w_U^{\top}\bm c^U_t\bigr),
  \label{eq:nlp-analysis-copy-gate}
\end{equation}
其中$\bm d_t$是解码器当前状态，$\lambda_t\in(0,1)$是从历史复制的权重。$\lambda_t$接近0时模型照抄当前话语，接近1时从历史中取词；当前话语既无指代也无省略时，理想的行为是整句照抄、$\lambda_t$始终接近0。训练以参考改写的负对数似然为目标函数；这一模型无法生成历史与当前话语中都没有的词，因此它要求参考改写中的词都能在这两处找到。

用“那明天呢”在H1下的改写看一步计算（数值为教学构造）。解码到“明天能到”之后，下一个词应为“图书馆”。设$\lambda_t=0.8$，历史注意力在“图书馆”上为$0.9$，当前话语中没有这个词，则$p(\text{图书馆})=0.8\times0.9=0.72$；“明天”只在当前话语中出现，若当前注意力在它上面为$0.6$，则$p(\text{明天})=0.2\times0.6=0.12$。模型输出“图书馆”。若门控错误地给出$\lambda_t=0.3$，“图书馆”的概率降为$0.27$，与当前话语中某个词的概率相当，改写就可能停在“明天能到呢”这样的不完整句子上。门控因此承担了“此处是否需要从历史补内容”的判断，相当于隐式地定位了省略位置。

只做复制的设计限制了模型补入新内容，这既是优点也是限制：它天然满足“补全内容来自历史”这一要求，却无法生成连接词、语序调整所需的虚词。其他方法从不同角度处理同一问题。Pan等的方法分两步，先在历史中逐词标注哪些词应被补入，再把选出的词附在当前话语之后，交给一个指针生成网络写出完整句子，选词一步可以单独评价\scite{nlp-analysis-pan2019restoration}。选词写成历史词元上的二分类：用预训练编码器把历史与当前话语一起编码，第$i$个历史词元的表示为$\bm h_i$，补入概率为$\sigma(\bm w^{\top}\bm h_i)$；Liu等把改写看作在“当前话语词×历史词”的矩阵上预测编辑操作（在某处插入历史中的某个片段，或用历史片段替换当前话语中的某个词），用语义分割网络一次预测整个矩阵，避免了逐词解码的错误累积\scite{nlp-analysis-liu2020run}。与槽位继承相比，这些方法不需要槽位定义，覆盖面广，代价是输出不再是可以逐项检查的结构，补入的片段是否放在了正确的语法位置需要另外核对：“图书馆”被复制到人物的位置上，复制来源正确，改写仍然是错的。

\subsubsection{基于大语言模型的查询改写}
\label{sec:nlp-analysis-llm-rewriting}

前两种方法都需要成对的改写训练数据。大语言模型可以直接按提示完成改写：给出若干“历史、当前问题、独立问题”的示例，再给出当前对话，要求模型输出一个不依赖历史也能理解的问题。改写的要求可以写进提示：保留用户原意，补出省略和指代的内容，不添加历史中没有的信息，历史不足以补全时原样输出。

代码~\ref{code:nlp-analysis-rewrite-prompt}给出一个少样本改写提示的骨架，示例中同时包含一个无需改写和一个历史不足的情形，用来教模型在这两种情况下原样输出。

\begin{codeblock}[label={code:nlp-analysis-rewrite-prompt}]{少样本话语改写提示}
把“当前问题”改写为不看对话历史也能理解的独立问题。
只使用历史中出现过的信息；无需改写或历史不足时原样输出。

历史：今天能到图书馆借书吗？——可以，开到晚上九点。
当前问题：那明天呢？
独立问题：明天能到图书馆借书吗？

历史：（无）
当前问题：那明天呢？
独立问题：那明天呢？

历史：读书会改到周五下午两点了。
当前问题：读书会在哪里举行？
独立问题：读书会在哪里举行？

历史：{对话历史}
当前问题：{当前问题}
独立问题：
\end{codeblock}

在对话式检索中，改写的目的是检索效果，这带来两个新的做法。Mao等的LLM4CS让模型对同一问题生成多个改写，并生成一个假想的回答，把这些文本的嵌入加以聚合作为检索查询，以多个样本平均掉单次改写的偏差\scite{nlp-analysis-mao2023llm4cs}。Ye等提出“信息充分的改写”，要求改写不仅补全省略，还要补入历史中有助于检索的背景信息，并以正确、清晰、信息充分、不冗余四项要求编写提示；他们再用大模型的改写结果训练一个较小的改写模型，以降低推断成本\scite{nlp-analysis-ye2023informative}。

大语言模型改写的主要风险是补入历史中没有的内容。模型可能把“那明天呢”改写为“明天上午九点能到市图书馆借书吗”，“上午九点”“市”都是模型自己的猜测。一个简单而有效的检查是来源核对，见算法~\ref{alg:nlp-analysis-rewrite-check}：改写中每个不在当前话语中的实词，都应能在历史中找到对应；找不到来源的新增内容要么删去，要么交给人工或澄清流程。

\begin{algorithm}[改写结果的来源核对]
  \label{alg:nlp-analysis-rewrite-check}
  \begin{algorithmic}[1]
    \Require 对话历史$H$，当前话语$u$，候选改写$r$；可接受的规范化映射$N$（如“明天”到具体日期）
    \Ensure 通过核对的改写，或带标记的无来源片段
    \State 对$H$、$u$、$r$分词，去掉虚词与标点，得到实词集合$W_H$、$W_u$、$W_r$
    \State $D\gets W_r\setminus W_u$\Comment{改写中新增的实词}
    \For{每个$w\in D$}
      \If{$w\in W_H$，或$w$可由$N$从$W_H\cup W_u$中的某个词得到}
        \State 记录$w$的来源位置
      \Else
        \State 把$w$标记为无来源
      \EndIf
    \EndFor
    \If{没有无来源的词}
      \State \Return $r$
    \Else
      \State \Return 删去无来源片段后的$r$，或请求澄清，并附上被标记的词
    \EndIf
  \end{algorithmic}
\end{algorithm}

字符串层面的核对只能发现最明显的捏造：同义改述（历史说“借阅”，改写写“借书”）会被误报，需要配合同义词表或嵌入相似度；把正确来源的词放到错误位置，它也无法发现。本节关注的是语义补全，即改写是否忠实地还原了用户的问题；第~\ref{chap:nlp-question-answering-dialogue}章从检索与问答系统的角度讨论查询改写，那里的评价标准是检索效果，可能有意加入扩展词，两者的目标不同。

表~\ref{tab:nlp-analysis-rewriting-methods}对照三类方法。

\begin{table*}[htbp]
  \centering
  \small
  \begin{tabularx}{\linewidth}{@{}l XXX@{}}
    \toprule
    方法 & 解决的问题 & 核心机制 & 新的代价或局限\\
    \midrule
    槽位继承 & 结构化场景中的省略 & 逐槽判断历史槽值是否继承 & 只覆盖预定义槽位；槽间依赖需单独处理\\
    复制式改写 & 开放领域的任意省略 & 从历史与当前句两路复制，门控$\lambda$选择来源 & 需要改写训练数据；补入位置需另行核对\\
    大语言模型改写 & 无训练数据；面向检索的改写 & 少样本提示，多次采样聚合 & 可能补入无来源内容，需来源核对；推断成本高\\
    \bottomrule
  \end{tabularx}
  \caption{省略恢复与话语改写方法对照。}
  \label{tab:nlp-analysis-rewriting-methods}
\end{table*}

\subsection{会话含义与间接言语行为}
\label{sec:nlp-analysis-implicature}

上一小节补出的是说话者省略、但历史中明确出现过的内容。还有一类意义从未被说出，却是说话者想要传达的。问“同学们都到了吗”，回答“有些同学到了”，听话者通常会理解为“不是所有同学都到了”，尽管“有些”在逻辑上并不排除“全部”。问“你能把窗关上吗”，听话者知道这是请求关窗，而不是询问自己有没有关窗的能力。前者是会话含义，后者是间接言语行为，二者都要求在字面意义之上推断说话者的意图。

这类推断有两个挑战。其一，它依赖合作假设与背景：听话者之所以从“有些”推出“不是全部”，是因为假定对方提供了尽可能充分的信息；对方若不知道其余同学的情况，这一推断就不成立。其二，推断结果本身可能有分歧，不同标注者对同一句话是否带有某种含义的判断并不总是一致。方法的演进从语言学的规则与分类模型，过渡到用自然语言推断的形式检验模型的判断，再到对大语言模型的语用推理能力做专门评测。评测时的核心问题是区分模型是在推理含义，还是在匹配训练数据中的表面模式。

\subsubsection{合作原则与含义推理}
\label{sec:nlp-analysis-grice}

Grice提出，交谈的双方默认遵守\term{合作原则}（cooperative principle）：使自己的话语在当时的交谈目的下恰如其分\scite{nlp-analysis-grice1975}。合作原则展开为四条准则。量的准则：提供当前交谈所需的信息，不多也不少。质的准则：不说自己认为不真的话，不说缺乏证据的话。关系准则：说话要相关。方式准则：避免晦涩和歧义，简洁而有条理。这些准则描述的不是人们实际总是怎样说话，而是交谈双方默认对方会遵守的期望，听话者据此解释话语。当字面意义看起来违反某条准则，而听话者仍相信对方在合作时，就会推出一个使话语恢复合作的额外意义，这就是\term{会话含义}（conversational implicature）。

等级含义是最常见、也最常被计算研究的一类。“有些”与“全部”构成一个按信息强度排序的等级$\langle\text{有些},\text{全部}\rangle$，“全部”蕴含“有些”而反之不然。推理分四步：说话者说了较弱的“有些同学到了”；若全部到了，说较强的“全部到了”信息更充分，按量的准则说话者应当这样说；说话者没有这样说，又假定他在合作且了解情况；因此他认为“全部到了”不成立。“或”与“和”、“可能”与“一定”、“温”与“热”构成类似的等级。关系含义则依赖关系准则：甲问“还有咖啡吗”，乙答“楼下有便利店”，字面上答非所问，听话者假定它相关，推出“这里没有咖啡了，可以去楼下买”。

两类含义的依赖程度不同。Grice把不依赖特定语境、只要用了某种说法就默认出现的含义称为\term{一般会话含义}（generalized conversational implicature），把只在特定语境中才出现的含义称为\term{特殊会话含义}（particularized conversational implicature）。等级含义通常归入前者：只要说了“有些”，在没有相反信息时就推出“不是全部”。关系含义属于后者：“楼下有便利店”在另一场对话中不会有“没有咖啡了”的意思。前者有规律可循，适合用模板构造评测数据，后者必须结合具体对话判断。说话者还可能故意明显地违反某条准则，让听话者意识到违反并推出另一层意义：天下着大雨，有人说“真是出门的好天气”，明显违反质的准则，听话者据此推出反讽；推荐信只写“该生出勤良好、字迹工整”，明显违反量的准则，读者推出“学术能力不值一提”。反讽与言外之意的识别因此也是会话含义推理的一部分。

等级含义的推理可以写成一个可计算的概率模型。理性言语行为（rational speech act, RSA）模型把听话者与说话者的相互推理写成递归\scite{nlp-analysis-frank-goodman2012}。设可能的世界状态为$w$，话语为$u$，$[\![u]\!](w)\in\{0,1\}$表示$u$在$w$中字面上是否为真，$P(w)$为先验。字面听话者只按真值过滤：$L_0(w\mid u)\propto[\![u]\!](w)P(w)$。说话者按听话者能否猜对来选择话语：$S_1(u\mid w)\propto\exp\bigl(\alpha\log L_0(w\mid u)\bigr)$，$\alpha>0$衡量说话者的理性程度。语用听话者再推断说话者为什么选了这句话：$L_1(w\mid u)\propto S_1(u\mid w)P(w)$。

取两个世界“部分到了”与“全部到了”，先验各$0.5$，两句话“有些到了”（两个世界中都为真）与“全部到了”（只在后者中为真），$\alpha=1$。字面听话者听到“有些”时两个世界各$0.5$，听到“全部”时只剩“全部到了”。在“全部到了”的世界里，说话者说“有些”能让听话者猜对的概率为$0.5$，说“全部”为1，于是$S_1(\text{有些}\mid\text{全部})=0.5/1.5=1/3$；在“部分到了”的世界里，“全部”为假，$S_1(\text{有些}\mid\text{部分})=1$。语用听话者听到“有些”后，$L_1(\text{部分}\mid\text{有些})\propto1\times0.5$，$L_1(\text{全部}\mid\text{有些})\propto\frac13\times0.5$，归一化得$0.75$与$0.25$。字面意义对两个世界不作区分，递归一次之后，“不是全部”成为更可能的解读，这正是上面四步推理的概率版本；$\alpha$越大，偏向越明显。RSA的局限在于需要事先列出世界状态与备选话语，适合控制良好的语言游戏，难以直接用于开放文本。

会话含义与逻辑蕴含的区别在于\term{可取消性}（cancellability）：含义可以被后续话语明确否定而不产生矛盾，蕴含不能。“有些同学到了，实际上全部都到了”虽然略显迂回，但并不自相矛盾，说明“不是全部”只是含义；“全部同学都到了，实际上有些没到”则自相矛盾，说明“全部”蕴含“有些”。这一检验给出了标注与评测时区分两者的可操作标准。

\term{预设}（presupposition）是第三类需要区分的意义：它是话语成立所依赖的背景前提。“小王停止抽烟了”预设小王曾经抽烟；“小王知道会议改期了”预设会议确实改期了。引出预设的词语称为预设触发语，常见的有：有定描述（“小王的车”预设小王有车）、叙实动词（“知道”“意识到”）、状态改变动词（“不再”“开始”“停止”）、反复类副词（“又”“再”）以及分裂句（“是小王打碎了杯子”预设有人打碎了杯子）。预设的特征是在否定、疑问、条件等算子下仍然保留：“小王没有停止抽烟”“小王停止抽烟了吗”“如果小王停止抽烟”都仍然预设他曾经抽烟，而蕴含在否定下不保留。预设也不是完全不可取消：在否定句中，说话者可以明确否认它，“小王没有停止抽烟，他根本就没抽过烟”并不矛盾；但在肯定句中加上同样的否认就会自相矛盾。表~\ref{tab:nlp-analysis-meaning-types}用两条检验区分三种意义。

\begin{table}[htbp]
  \centering
  \small
  \begin{tabularx}{\linewidth}{@{}l X cc@{}}
    \toprule
    意义类型 & 例子 & 否定后是否保留 & 能否取消\\
    \midrule
    蕴含 & “全部同学到了”$\Rightarrow$“有些同学到了” & 否 & 否\\
    会话含义 & “有些同学到了”$\leadsto$“不是全部” & 否 & 能\\
    预设 & “小王停止抽烟”$\gg$“小王曾经抽烟” & 是 & 肯定句中不能\\
    \bottomrule
  \end{tabularx}
  \caption{蕴含、会话含义与预设的区分检验。$\leadsto$表示含义，$\gg$表示预设。}
  \label{tab:nlp-analysis-meaning-types}
\end{table}

\subsubsection{间接言语行为识别}
\label{sec:nlp-analysis-indirect-acts}

Austin指出，说话本身就是在做事：说“我保证明天交稿”就完成了一次许诺，这样的话语没有真假，只有是否得当\scite{nlp-analysis-austin1962}。Searle把这一思想系统化为言语行为理论，并用一组\term{适切条件}（felicity conditions）刻画每种行为成功实施的条件。以请求听话者H做动作A为例，四类条件为：命题内容条件，A是H将来的动作；准备条件，H有能力做A，且不请求的话H不一定会去做；真诚条件，说话者想要H做A；本质条件，这句话算作让H做A的一次尝试。Searle后来又按行为的目的把言语行为分为五类：断言类（陈述、报告，使话语与世界相符）、指令类（请求、命令、提问，让听话者去做某事）、承诺类（许诺、答应，说话者约束自己将来的行为）、表达类（感谢、道歉，表达心理状态）、宣告类（任命、宣布，话语本身改变了世界的状态）。第~\ref{sec:nlp-analysis-dialogue-acts}小节的话语行为标签集，大体可以看作这一分类在对话中的细化，例如请求与是非问都属于指令类，感谢与附和更接近表达类。

\term{间接言语行为}（indirect speech act）是通过实施一种行为来完成另一种行为。Searle观察到，间接请求的句式系统地对应于对适切条件的询问或陈述\scite{nlp-analysis-searle1975}。“你能把窗关上吗”询问的是准备条件中H的能力；“我希望你把窗关上”陈述的是真诚条件；“你会把窗关上吗”询问的是命题内容条件中H将来的动作。听话者的推理过程大致是：字面上这是一个关于能力的问题；在当前情境中，对方显然知道我有能力关窗，这个问题作为信息询问没有意义；询问能力是请求的准备条件之一；因此对方很可能是在请求我关窗。这里同样依赖合作假设，并且推理的关键一步是“字面解读在情境中不合理”。若情境改变，例如在讨论受伤后的康复情况，“你能把窗关上吗”就可能是真的在问能力。

把这一推理拆成可检查的步骤，便于看出它在哪里依赖情境：
\begin{enumerate}
  \item 确定字面行为：话语的句式是疑问，字面上询问H是否有能力做A。
  \item 检查字面解读是否合作：在当前情境中，S是否已经知道答案；若已知，作为信息询问就违反了量或关系的准则。
  \item 寻找被询问或陈述的适切条件：能力属于请求的准备条件。
  \item 推断间接行为：S通过询问准备条件来实施请求。
  \item 检查情境是否支持该请求：A是否对S有益、H是否方便做A；不支持时退回字面解读。
\end{enumerate}
第2步与第5步都需要情境知识，这也是计算方法最难自动获得的部分。

计算上的识别方法经历了两个阶段。Perrault与Allen用基于规划的方法建模这一推理：把言语行为表示为带前提与效果的规划算子，听话者从字面行为出发，推断说话者可能的目标，再识别出间接的请求\scite{nlp-analysis-perrault-allen1980}。这种方法可解释，但需要手工编写情境知识。统计方法把识别写成条件分类：输入当前话语和可见的上下文，输出它实际承担的行为，即$p(a\mid u,\text{上下文})$，与第~\ref{sec:nlp-analysis-dialogue-acts}小节的话语行为分类同形，区别只在于标签表示的是实际意图而不是句式。间接回答是一个重要的特例：问“你周末要来聚会吗”，回答“我得赶一篇论文”，要识别的是“不来”。识别时若上下文不足以排除字面解读，模型应当保留字面解读，而不是默认存在间接意图；“这里有点冷”在讨论天气时只是陈述感受，在对方站在窗边时才可能是请求关窗。

\subsubsection{语用推理能力的评测}
\label{sec:nlp-analysis-pragmatic-eval}

语用推理很难直接标注，评测通常把它转写为已有任务的形式。IMPPRES把含义与预设写成自然语言推断问题\scite{nlp-analysis-jeretic2020imppres}。对等级含义，给出前提“有些同学到了”与假设“不是所有同学都到了”，按语用解读应判为蕴含，按逻辑解读应判为中立，模型的选择反映它采用哪一种解读。对预设，数据把触发语放在否定、疑问、情态、条件等算子下，检验模型是否仍然推出预设。数据由模板生成，规模可控，便于逐现象分析。Jeretic等发现，在MultiNLI上训练的BERT稳定地把“some”引出的等级含义当作蕴含，即采用语用解读；对部分预设触发语（如only），即使触发语处在否定之下，它也能推出预设，对另一些触发语则不能。

Circa检验间接回答的理解\scite{nlp-analysis-louis2020circa}。数据含约3.4万个是非问句与间接回答对，由众包人员在给定情境下写出问题和不直接说“是”或“否”的回答，标签包括是、否、有条件的是、介于两者之间等。模型需要从回答推断它对问题的立场。Louis等的基线把问题与回答拼接后送入BERT做句对分类，并比较了先在自然语言推断数据MultiNLI上微调、再在Circa上微调的做法；后者的准确率最高，但与不经迁移的模型相比差异在统计上不显著，因此不能据此断言推断数据带来了稳定收益。他们还比较了只给回答、不给问题的设置：准确率显著下降，说明模型确实利用了问题，但下降后的准确率仍然不低，相当一部分样本单凭回答的表面形式就能判断。在训练中未出现过的情境上测试时，准确率也会下降。

对大语言模型，Ruis等设计了二选一的会话含义评测：把一问一答写成“甲问……，乙答……，乙的意思是”，比较模型续写“是”与“否”的概率\scite{nlp-analysis-ruis2023}。例如“甲问：你周末要来聚会吗？乙答：我得赶一篇论文。乙的意思是”，若模型给“否”的概率为$0.7$、“是”为$0.3$，记为答对。同一样本换用措辞不同的六个模板，若三个模板答对、三个答错，单一模板的结论就完全取决于选了哪一个，因此应报告在所有模板上的平均准确率及其范围。他们比较了只做预训练、对话微调、基准级指令微调与样例级指令微调四类模型，结论是：只做预训练的模型在零样本下接近随机水平，样例级指令微调的模型明显更好，少样本示例进一步提高准确率，但多数模型仍低于人类的一致率。在他们的设置中，六个模板之间的波动不大；这一稳健性依赖于对多个模板取平均的评测设计，换用新的格式时仍需重新检验。Hu等用多项选择形式比较人与语言模型在讽刺、隐喻、间接言语、违反准则等七类现象上的表现，并分析模型的错误类型与人是否相近\scite{nlp-analysis-hu2023pragmatics}。

这些评测的设计给出几条经验。其一，结果可能随提示措辞和示例变化，报告时应给出多个模板的均值与波动，而不是只报最好的模板。其二，需要人类基线：人对同一问题的判断本身存在分歧，模型的准确率应与人类的一致程度对照看。其三，要设计对照样本以区分推理与模式匹配：保持回答不变、改变问题或情境，若模型的判断不随之改变，说明它依赖的是回答的表面形式，而不是推理。语用评测的分数衡量的是模型在特定格式下给出符合多数人解读的答案的能力，能否推广到开放对话中的语用理解，需要结合具体应用检验。


\section*{文献说明}

本章的任务定义和评价对象来自若干标注规范与基准，下列资料分别规定了标注对象、结构形式或任务范围。

词法与句法的标注对象依据宾州树库与Universal Dependencies，后者还给出了允许多入边的增强依存规范。

\begin{fullwidth}
\raggedright
\fullcite{nlp-analysis-marcus1993}\par
\fullcite{nivre2020universal}\par
\fullcite{v9a-ud-enhanced}
\end{fullwidth}

谓词论元与语义角色依据PropBank与FrameNet，句子语义图依据AMR及其Smatch评价，文本到SQL的跨数据库设定依据Spider系列基准。

\begin{fullwidth}
\raggedright
\small
\fullcite{nlp-analysis-palmer2005}\par\medskip
\fullcite{nlp-analysis-baker1998}\par\medskip
\fullcite{nlp-analysis-banarescu2013}\par\medskip
\fullcite{cai-knight-2013-smatch}\par\medskip
\fullcite{yu-etal-2018-spider}\par\medskip
\fullcite{nlp-analysis-spider2025}\par\medskip
\fullcite{yu-etal-2019-sparc}\par\medskip
\fullcite{yu-etal-2019-cosql}
\end{fullwidth}

篇章分析中，共指依据OntoNotes上的CoNLL-2012共享任务，篇章关系依据PDTB与RST篇章树库，篇章单元与关系的跨框架任务依据DISRPT，论证结构依据说服性文章语料。

\begin{fullwidth}
\raggedright
\small
\fullcite{nlp-analysis-pradhan2012}\par\medskip
\fullcite{nlp-analysis-prasad2008}\par\medskip
\fullcite{nlp-analysis-carlson2001}\par\medskip
\fullcite{braud-etal-2023-disrpt}\par\medskip
\fullcite{stab-gurevych-2017-parsing}
\end{fullwidth}

语用分析中，话语行为依据Switchboard与ICSI会议语料的标注体系，上下文中的问题改写依据CANARD与TREC CAsT。

\begin{fullwidth}
\raggedright
\small
\fullcite{stolcke-etal-2000-dialogue}\par\medskip
\fullcite{nlp-analysis-shriberg2004}\par\medskip
\fullcite{nlp-analysis-elgohary2019}\par\medskip
\fullcite{nlp-analysis-dalton2020}
\end{fullwidth}

五种范式的代表方法分别见以下原始论文：线性链CRF、结构化感知机、转移式依存分析的系统化描述、最大生成树依存分析、端到端跨度共指，以及把句法树线性化后用序列到序列模型生成的做法。

\begin{fullwidth}
\raggedright
\small
\fullcite{nlp-classical-lafferty2001}\par\medskip
\fullcite{collins2002}\par\medskip
\fullcite{nlp-analysis-nivre2008}\par\medskip
\fullcite{v9a-mcdonald2005}\par\medskip
\fullcite{lee2017coreference}\par\medskip
\fullcite{nlp-analysis-vinyals2015grammar}
\end{fullwidth}

\section{本章小结}
\label{sec:nlp-analysis-summary}

章首列出的五项挑战，现在可以逐项收拢。分析单位的不确定性由序列标注加合法掩码处理：分词把边界写成BMES标签，预训练编码器与词格模型补充语境和词典证据，所有结果最后恢复为原文字符跨度；模型词元与语言词的对齐由首子词规则或词元合成规则显式处理。结构多样与合法性约束由五种范式分别回答：序列标注加掩码处理平面区间，跨度范式处理嵌套，转移系统以线性时间构建树和图，图与区间动态规划在一阶可加分解下给出精确最优，约束生成以统一接口覆盖逻辑式、语义图和SQL，并用文法掩码保证可解析。远距离歧义主要靠上下文编码器和全局推断消解：全局归一化或精确解码让后面的证据能够修正前面的选择，这是Viterbi、CKY与Eisner共同的机制。标注成本与规范差异部分由预训练和多标准学习缓解，但规范本身决定了什么算正确，模型只能在训练过的规范之间切换。篇章与语用对背景的依赖则只解决到一部分：模型可以利用上文与对话历史，文本之外的意图和共同知识仍须由任务提供或明确标为未知。

各任务的方法演进呈现出相同的模式：从生成式独立假设到判别式全局归一化，从人工特征到上下文编码，从流水线到联合解码，再到以线性化生成统一各种结构。每一步解决前一步的一个具体缺陷，也带来新的代价：全局归一化需要计算配分函数，神经编码增加计算，联合标签扩大输出空间，生成范式失去全局最优保证并增加核验负担。选择范式时，应依据结构是否嵌套或交叉、是否需要精确推断与概率、监督是结构还是答案，以及长度与时延限制来权衡。

不同结构的合法性检查不能互相替代。树的单中心词约束不适用于共享论元的语义图，共指簇的传递性不适用于桥接指代，查询能执行也不表示它符合原问题。实际系统应同时保留原文位置、标注规范、可用背景和恢复轨迹，使候选遗漏、评分错误、结构冲突与依据不足能够分别诊断。下一章在这些结构之上抽取有身份、角色和证据的事实记录，并要求输出同时受任务本体与原文证据约束。

\paragraph{延展阅读}
UD可用于核对跨语言的词法与依存标注对象，自注意成分解析展示了上下文编码与树恢复的分工，DISRPT整理了篇章单元与关系任务，Spider把语义分析落实到跨数据库查询。阅读时宜分别记录单位规范、合法结构、可见模式与评价分母\scite{nivre2020universal,nlp-intro-kitaev2018,braud-etal-2023-disrpt,yu-etal-2018-spider}。
